\documentclass[a4paper,12pt]{report}
\usepackage[left=2.5cm, right=2.5cm, top=2cm, bottom=2.2cm]{geometry}
\usepackage{amssymb}
\usepackage[T1]{fontenc}
\usepackage[utf8]{inputenc}
\usepackage[french]{babel}
\usepackage{hyperref}
\usepackage{fancyhdr}
\usepackage{pdfpages}
\usepackage{amsmath}
\usepackage{supertabular}
\usepackage[table]{xcolor}
\renewcommand{\footrulewidth}{0.4pt}
\fancyfoot[L]{Projet de fin d'études}
\rhead{\leftmark}
\usepackage{graphicx}
\usepackage{float}
\usepackage{longtable}
\usepackage{indentfirst}
\usepackage{subcaption}
\usepackage[export]{adjustbox}
\usepackage{wrapfig}
\usepackage{color}
\usepackage{titletoc}
\usepackage{titlesec}
\usepackage{lastpage}
\usepackage[Lenny]{fncychap}
\usepackage{enumitem}
\usepackage{pifont}
\usepackage{etoc}
\usepackage{mdframed}
\usepackage{tikz}
\usepackage{booktabs}
\usepackage{multirow}
\usepackage{array}
\usepackage{needspace}
\usepackage{caption}
\usepackage{tabularx}

\captionsetup[longtable]{
    labelfont=normalfont,
    textfont=normalfont,
    format=plain
}

% ---------------------------------------------------------------
% \img{largeur}{fichier} : insère Image/fichier s'il existe,
% sinon affiche un cadre « image à insérer » pour que le rapport
% compile même avant d'avoir ajouté toutes les captures.
% ---------------------------------------------------------------
\newcommand{\img}[2]{%
  \IfFileExists{Image/#2}%
    {\includegraphics[width=#1\textwidth]{Image/#2}}%
    {\fbox{\parbox[c][4.5cm][c]{#1\textwidth}{\centering\small\textit{Image à insérer :}\\[4pt]\texttt{\detokenize{#2}}}}}%
}
\newcommand{\logo}[2]{%
  \IfFileExists{Image/#2}%
    {\includegraphics[height=#1]{Image/#2}}%
    {\fbox{\parbox[c][#1][c]{1.4cm}{\centering\tiny logo}}}%
}

\begin{document}

\setcounter{secnumdepth}{3}
\setcounter{tocdepth}{2}

% Page de garde : remplacer par le PDF officiel ESPRIT si disponible
\IfFileExists{firstpage.pdf}
  {\includepdf[pages=-,pagecommand={\thispagestyle{empty}}]{firstpage.pdf}}
  {% Page de garde provisoire : elle est remplacée automatiquement par
% firstpage.pdf si ce fichier (modèle officiel ESPRIT) est ajouté au dossier.
\begin{titlepage}
\thispagestyle{empty}
\begin{center}
\begin{minipage}{0.3\textwidth}
\raggedright\logo{1.6cm}{esprit.png}
\end{minipage}\hfill
\begin{minipage}{0.3\textwidth}
\raggedleft\logo{1.6cm}{Antigonelogo.jpg}
\end{minipage}

\vspace{1.2cm}
{\large \textsc{École Supérieure Privée d'Ingénierie et de Technologies}}\\[0.2cm]
{\large ESPRIT}

\vspace{1.2cm}
{\Large \textbf{Rapport de Projet de Fin d'Études}}\\[0.3cm]
{\large Présenté en vue de l'obtention du}\\[0.2cm]
{\large \textbf{Diplôme National d'Ingénieur en Informatique}}\\[0.2cm]
{\large Option : \textbf{TWIN} (Technologies du Web et de l'Internet)}

\vspace{1.2cm}
\rule{\textwidth}{1.2pt}\\[0.4cm]
{\LARGE \textbf{Conception et Développement d'une Plateforme\\[0.2cm] Transactionnelle Intelligente}}\\[0.4cm]
{\large \textit{Sellio : plateforme SaaS multi-boutiques pour le e-commerce\\ de cosmétiques et de vêtements}}\\[0.3cm]
\rule{\textwidth}{1.2pt}

\vspace{1.2cm}
{\large Réalisé par : \textbf{Rania KHEDHRI}}

\vspace{1cm}
\begin{tabular}{ll}
\textbf{Encadrant académique :} & \textit{[Nom de l'encadrant(e) ESPRIT]} \\[0.2cm]
\textbf{Encadrant professionnel :} & M. Malek NAOUAR \\
\end{tabular}

\vspace{1cm}
{\large Réalisé au sein de \textbf{Antigone}}

\vfill
{\large Année universitaire 2025 -- 2026}
\end{center}
\end{titlepage}
}
\IfFileExists{secondpage.pdf}
  {\includepdf[pages=-,pagecommand={\thispagestyle{empty}}]{secondpage.pdf}}
  {}

\cleardoublepage
\pagenumbering{Roman}
\setcounter{page}{1}
\pagestyle{plain}

\begin{center}
\textbf{\Huge Remerciements}
\end{center}

\vspace{1cm}

\hspace{0.5cm}Avant de présenter ce travail, je tiens tout d'abord à exprimer ma profonde gratitude à mon encadrant(e) académique, \textbf{\textit{[Nom de l'encadrant(e) ESPRIT]}}, pour son accompagnement, sa disponibilité et ses conseils tout au long de ce projet. Ses remarques ont guidé mes choix, en particulier sur la démarche d'évaluation des modèles d'apprentissage automatique.

\vspace{0.5cm}

Je souhaite également adresser mes sincères remerciements à mon encadrant professionnel et directeur de l'entreprise Antigone, \textbf{M. Malek Naouar}, pour son accueil, la confiance qu'il m'a accordée en me confiant un projet de cette ampleur, et pour ses orientations qui ont permis de faire évoluer une boutique en ligne en une véritable plateforme multi-boutiques.

\vspace{0.5cm}

Je tiens aussi à remercier l'ensemble de l'équipe de \textbf{l'entreprise Antigone} pour sa collaboration, sa disponibilité et l'environnement professionnel et bienveillant dans lequel j'ai eu l'occasion d'évoluer.

\vspace{0.5cm}

Enfin, je remercie les membres du jury qui ont accepté d'évaluer ce travail, en espérant qu'ils y trouveront la rigueur et la motivation qui ont guidé sa réalisation.

\vspace{0.5cm}

\begin{center}
Avec toute ma gratitude.
\end{center}

\thispagestyle{plain}
\begin{center}
\textbf{\Huge Dédicace}
\end{center}

% Texte proposé : à personnaliser librement.

\vspace{1.5cm}

\begin{center}
\textit{
À \textbf{mes chers parents},\\
pour leur amour, leurs sacrifices et leur soutien sans faille,\\
qui m'ont toujours donné la force d'aller au bout de mes projets.
}
\end{center}

\vspace{0.8cm}

\begin{center}
\textit{
À \textbf{toute ma famille},\\
pour sa présence et ses encouragements\\
tout au long de mon parcours.
}
\end{center}

\vspace{0.8cm}

\begin{center}
\textit{
À \textbf{mes amis},\\
pour les moments partagés, les nuits de révision\\
et les encouragements dans les moments difficiles.
}
\end{center}

\vspace{0.8cm}

\begin{center}
\textit{
À vous tous,\\
je dédie ce travail avec toute mon affection et ma gratitude.
}
\end{center}
\begin{center}
\raggedleft\textbf{Rania Khedhri}
\end{center}

\cleardoublepage

\tableofcontents
\listoffigures
\cleardoublepage
\listoftables

\cleardoublepage
\pagenumbering{arabic}
\setcounter{page}{1}
\newpage

\vspace*{-1.7cm}
\begin{center}
\noindent\textbf{\rule{\textwidth}{1.2pt}}

\vspace{1.2cm}
{\Huge \textbf{Introduction Générale}}
\vspace{0.5cm}
\end{center}

\addcontentsline{toc}{chapter}{Introduction Générale}

Le commerce électronique n'est plus réservé aux grandes enseignes. En Tunisie comme ailleurs, une part croissante des petites marques de cosmétiques naturels et de prêt-à-porter vendent déjà en ligne, mais le plus souvent à travers une page Instagram, un catalogue partagé sur WhatsApp et une prise de commande manuelle. Ces canaux suffisent tant que le volume reste faible ; dès que les commandes se multiplient, ils montrent leurs limites : stocks tenus à la main, paiements à la livraison difficiles à suivre, aucune donnée exploitable sur le comportement des clients, et une image de marque qui dépend d'une plateforme tierce. Les solutions de création de boutique en ligne comme Shopify ou WooCommerce existent, mais elles sont pensées pour un marché international, facturées en devises, peu adaptées aux spécificités locales (TVA tunisienne, zones de livraison, dinar) et elles réservent les fonctions avancées, recommandation ou analyse prédictive, à des extensions payantes.

\vspace{0.5cm}
Antigone, agence tunisienne de conseil en stratégie et en marketing digital, accompagne justement ce type de marques. Elle a d'abord développé pour l'une d'elles, NaturEssence, une boutique en ligne de cosmétiques naturels. L'expérience a montré que la plupart des besoins de cette marque se retrouvaient chez les autres clients de l'agence : un catalogue avec des variantes (tailles, couleurs, matières), une vitrine personnalisable, un tunnel de commande fiable, des outils marketing et une compréhension fine des clients. D'où l'idée de transformer ce développement unique en une plateforme mutualisée.

\vspace{0.5cm}
C'est dans ce cadre que s'inscrit notre projet de fin d'études, réalisé au sein d'Antigone dans le cadre du cycle ingénieur à l'École ESPRIT. Il consiste à concevoir et développer \textsf{\textbf{« Sellio »}}, une plateforme transactionnelle intelligente de type SaaS multi-boutiques : chaque commerçant y crée sa propre boutique de cosmétiques ou de vêtements, la fait vérifier par l'équipe Sellio (vérification d'identité et de moyen de paiement), choisit l'un des trois gabarits de vitrine (Minimal, Bold, Luxury) et le personnalise, gère son catalogue, ses commandes, ses promotions et son programme de fidélité. La plateforme repose sur une architecture microservices (Spring Boot, passerelle API, PostgreSQL), des interfaces React, le paiement en ligne par Stripe et une cabine d'essayage virtuel en temps réel fondée sur l'API Decart. Elle intègre enfin un volet d'apprentissage automatique : suivi comportemental des visiteurs, moteur de recommandation hybride, segmentation des visiteurs et prédiction de l'attrition, entraînés et évalués sur des jeux de données publics réels.

\vspace{0.5cm}
Notre rapport est composé d'une introduction, de huit chapitres et d'une conclusion générale. Ces huit chapitres sont les suivants :
\begin{itemize}
\item Le \textbf{chapitre 1} présente le contexte général du projet, l'organisme d'accueil, la problématique traitée, l'étude de l'existant et la méthodologie adoptée.
\item Le \textbf{chapitre 2} est consacré à l'analyse et à la spécification des besoins : acteurs, exigences fonctionnelles et non fonctionnelles, architecture retenue, backlog produit et planification des sprints.
\item Le \textbf{chapitre 3} présente la Release 1 : fondations de la plateforme Sellio (authentification, rôles, création de boutique, vérification KYC du commerçant et console d'administration).
\item Le \textbf{chapitre 4} est dédié à la Release 2 : catalogue produit et vitrine personnalisable (attributs de variantes, gabarits Minimal / Bold / Luxury, éditeur de thème).
\item Le \textbf{chapitre 5} présente la Release 3 : parcours transactionnel (panier, tunnel de commande, paiement Stripe, commandes, retours, TVA et livraison, avis).
\item Le \textbf{chapitre 6} est consacré à la Release 4 : marketing et fidélisation (promotions, bannières, e-mailing, programme de fidélité et segments).
\item Le \textbf{chapitre 7} présente la Release 5 : l'essayage virtuel en temps réel par intelligence artificielle générative (API Decart).
\item Le \textbf{chapitre 8} est dédié à la Release 6 : l'intelligence comportementale (suivi des événements, recommandation, segmentation et prédiction de l'attrition par apprentissage automatique).
\end{itemize}

Nous terminons ce rapport par une conclusion générale qui revient sur le travail accompli et propose des perspectives d'évolution pour la plateforme.

\thispagestyle{fancy}
\fancyhead[R]{A.U 2025-2026}
\fancyhead[L]{ chapitre 1: Cadre général du projet }
\chapter{Cadre général du projet}
\section*{Introduction}
Ce premier chapitre pose les bases du projet. Nous présentons d'abord l'organisme d'accueil, puis le contexte qui a conduit à la naissance de Sellio et la problématique à laquelle il répond. Nous étudions ensuite les principales solutions de création de boutiques en ligne afin d'en dégager les limites, avant de décrire la solution proposée. Nous terminons par l'étude des méthodologies de développement qui a conduit au choix de Scrum.

\section{Cadre général}
Ce projet s'inscrit dans le cadre du projet de fin d'études en vue de l'obtention du Diplôme National d'Ingénieur en Informatique, option \textbf{TWIN} (Technologies du Web et de l'Internet), à l'École ESPRIT. Il a été réalisé au sein de l'agence Antigone sur une période de six mois, du \textbf{2 février 2026} au \textbf{2 août 2026}. Le sujet officiel, « \textit{Conception et développement d'une plateforme transactionnelle intelligente} », consiste à concevoir une plateforme de commerce électronique distribuée, capable d'analyser en temps réel les interactions des utilisateurs et d'exploiter ces données par des techniques d'apprentissage automatique afin de personnaliser l'expérience d'achat. Ce sujet s'est concrétisé sous la forme de \textbf{Sellio}, une plateforme SaaS permettant à des commerçants de cosmétiques et de vêtements de créer et d'exploiter leur propre boutique en ligne.

\section{Présentation de l'organisme d'accueil}
\begin{samepage}
\noindent
\begin{minipage}{\textwidth}
Antigone est une agence de stratégie et de marketing digital basée à Radès, fondée en 2021. Elle accompagne ses clients dans la construction de leur stratégie digitale : diagnostic de la situation existante, définition des axes de communication, puis mise en œuvre à travers la création de sites web, d'applications mobiles et de campagnes. Une partie croissante de ses clients sont de petites marques qui souhaitent vendre en ligne, ce qui a naturellement conduit l'agence à s'intéresser au commerce électronique.
\begin{figure}[H]
    \centering
    \img{0.40}{Antigonelogo.jpg}
    \caption{Logo d'Antigone}
    \label{fig:Antigone_logo}
\end{figure}
\end{minipage}
\end{samepage}

\subsection{Fiche d'identité d'Antigone}
Le tableau~\ref{tab:coordonnees} présente la fiche d'identité d'Antigone.
\begin{table}[H]
\centering
\renewcommand{\arraystretch}{1.3}
\begin{tabular}{|l|l|}
\hline
\textbf{Nom de la société} & \textbf{Antigone} \\
\hline
Fondateur et représentant & M. Malek Naouar \\
\hline
Adresse & 12, Rue Habib Thameur, Radès 2040, Ben Arous \\
\hline
Téléphone & (+216) 55 557 193 \\
\hline
Adresse e-mail & antigoneconsulting.info@gmail.com \\
\hline
Site web & https://antigoneagency.com/ \\
\hline
Logo & \logo{1.2cm}{Antigonelogo.jpg} \\
\hline
\end{tabular}
\caption{Fiche d'identité d'Antigone~\cite{Antigone2026}}
\label{tab:coordonnees}
\end{table}

\subsection{Domaines d'activité d'Antigone}
Les activités d'Antigone s'organisent autour de trois axes :
\begin{enumerate}
  \item \textbf{Stratégie de marketing digital :}
  \begin{itemize}
    \item réalisation d'un \textbf{diagnostic} de la présence digitale du client ;
    \item définition des \textbf{axes de communication} les plus pertinents ;
    \item accompagnement dans la \textbf{construction et l'optimisation} de la stratégie marketing.
  \end{itemize}
  \item \textbf{Conseil, études et accompagnement :}
  \begin{itemize}
    \item \textbf{conseils et études} sur mesure ;
    \item démarche pragmatique fondée sur le \textbf{marketing des services} et le \textbf{marketing de l'innovation}.
  \end{itemize}
  \item \textbf{Création et développement digital :}
  \begin{itemize}
    \item conception et développement de \textbf{sites web} et de \textbf{boutiques en ligne} ;
    \item conception et développement d'\textbf{applications mobiles} ;
    \item mise à disposition de solutions adaptées à la réussite des projets clients.
  \end{itemize}
\end{enumerate}

\section{Contexte du projet}
Dans cette section, nous décrivons le problème, examinons les solutions existantes, puis présentons la solution proposée.

\subsection{Problématique}
Le projet est né d'un développement précédent : \textbf{NaturEssence}, une boutique en ligne de cosmétiques naturels réalisée par Antigone pour l'un de ses clients. Au fil des échanges avec d'autres marques accompagnées par l'agence, un constat s'est imposé : une boutique de vêtements ou une autre marque de cosmétiques avait exactement les mêmes besoins, un catalogue avec variantes, une vitrine à l'image de la marque, un tunnel de commande, des promotions, un suivi des clients. Refaire ce travail boutique par boutique coûte cher et produit des applications isolées, chacune à maintenir séparément.\newline

\vspace{6pt}
Du côté des commerçants, la situation n'est pas meilleure. Beaucoup vendent aujourd'hui via Instagram ou WhatsApp : les commandes sont prises en message privé, le stock est tenu dans un carnet ou un tableur, les tailles et les couleurs disponibles sont communiquées au cas par cas, et aucune donnée n'est conservée sur ce que les visiteurs regardent, ajoutent au panier puis abandonnent. Ceux qui passent à une plateforme comme Shopify se heurtent à d'autres obstacles : abonnement en devises, paramétrage de la TVA et des zones de livraison tunisiennes, et fonctions avancées (recommandations, essayage virtuel, analyse prédictive) proposées sous forme d'extensions payantes indépendantes, qui ne partagent pas leurs données.\newline

\vspace{6pt}
Enfin, une plateforme qui héberge plusieurs commerçants pose des questions propres au modèle SaaS : comment s'assurer qu'un commerçant est bien celui qu'il prétend être avant de le laisser encaisser des paiements, comment isoler strictement les données d'une boutique de celles d'une autre, et comment offrir à chacun une vitrine réellement différente sans développer un site par client.\newline

\vspace{6pt}
\textbf{La problématique de ce projet peut être résumée ainsi :}
\vspace{6pt}
\begin{center}
\fbox{
    \begin{minipage}{0.85\textwidth}
        \vspace{8pt}
        \begin{center}
        \textbf{Comment concevoir une plateforme transactionnelle distribuée, fondée sur une architecture microservices, qui permette à des commerçants de cosmétiques et de vêtements de créer, personnaliser et exploiter leur propre boutique en ligne en toute sécurité, tout en analysant en temps réel les interactions de leurs visiteurs grâce à l'apprentissage automatique afin de personnaliser l'expérience d'achat et d'anticiper la perte de clients ?}
        \end{center}
        \vspace{8pt}
    \end{minipage}
}
\end{center}
\vspace{6pt}

\subsection{Étude de l'existant}
Avant de construire une plateforme, il était légitime d'examiner comment les solutions établies répondent aux besoins identifiés. Trois d'entre elles sont représentatives du marché.
\begin{itemize}
\item \textbf{Shopify}~\cite{Shopify2026} est la référence des plateformes SaaS de création de boutiques. Elle offre des thèmes, une gestion de catalogue avec variantes, un paiement intégré et un vaste magasin d'applications. En contrepartie, l'abonnement est facturé en devises, la personnalisation poussée d'un thème demande de modifier son code (Liquid), et les fonctions intelligentes (recommandations avancées, essayage virtuel, prédiction de l'attrition) dépendent d'applications tierces payantes, chacune avec ses propres données.

\item \textbf{WooCommerce}~\cite{WooCommerce2026} est une extension libre de WordPress. Elle est très flexible et gratuite à la base, mais le commerçant doit gérer lui-même l'hébergement, les mises à jour, la sécurité et la compatibilité des extensions. Elle n'est pas multi-boutiques par nature : chaque commerçant exploite sa propre installation.

\item \textbf{Wix eCommerce}~\cite{Wix2026} mise sur la simplicité : un éditeur visuel par glisser-déposer permet de créer une boutique rapidement. Ses capacités de catalogue (variantes complexes, attributs métier) et d'analyse restent toutefois limitées, et la plateforme n'offre ni vérification d'identité des vendeurs ni fonctions prédictives.
\end{itemize}
\Needspace{15cm}
Le tableau~\ref{tab:existing} compare ces trois solutions sur les critères qui comptent le plus pour ce projet.
\begin{table}[H]
\centering
\renewcommand{\arraystretch}{1.4}
\begin{tabular}{|p{4.6cm}|p{2.9cm}|p{2.9cm}|p{2.9cm}|}
\hline
\textbf{Critères} & \textbf{Shopify} & \textbf{WooCommerce} & \textbf{Wix eCommerce} \\
\hline
Multi-boutiques hébergé (SaaS) & Oui & Non & Oui \\
\hline
Variantes métier (taille, couleur, tissu) & Oui & Oui (extensions) & Partiel \\
\hline
Gabarits de vitrine personnalisables sans code & Partiel & Partiel & Oui \\
\hline
Paramétrage TVA / livraison tunisiennes & Manuel & Manuel & Manuel \\
\hline
Vérification d'identité du commerçant (KYC) & Via le prestataire de paiement & Non & Via le prestataire de paiement \\
\hline
Recommandations personnalisées & Extensions & Extensions & Limitées \\
\hline
Prédiction de l'attrition & Non (extensions) & Non & Non \\
\hline
Essayage virtuel & Extensions & Extensions & Non \\
\hline
Programme de fidélité & Extensions & Extensions & Extensions \\
\hline
Modèle économique & Abonnement mensuel en devises & Gratuit + hébergement et extensions & Abonnement mensuel en devises \\
\hline
\end{tabular}
\caption{Comparaison des solutions existantes de création de boutiques en ligne~\cite{Shopify2026, WooCommerce2026, Wix2026}}
\label{tab:existing}
\end{table}

Ces solutions sont matures, mais aucune ne réunit dans un même produit, avec des données partagées, la gestion multi-boutiques, les variantes propres à la mode et aux cosmétiques, la vérification des commerçants, l'essayage virtuel et l'analyse comportementale par apprentissage automatique. Les fonctions intelligentes y sont des modules ajoutés après coup, qui ne voient qu'une partie des données. C'est cet écart qui justifie le développement de Sellio.

\section{Solution proposée}
La solution retenue est une plateforme web baptisée \textbf{Sellio}. Elle se compose de trois espaces : le \textbf{site de la plateforme} (présentation, inscription des commerçants et console d'administration), le \textbf{back-office du commerçant} pour gérer sa boutique, et la \textbf{vitrine} publique de chaque boutique, sur laquelle ses clients achètent.

\subsection{Espace plateforme et intégration des commerçants}
\begin{itemize}
    \item Page d'accueil Sellio, inscription et connexion des commerçants ;
    \item Assistant de création de boutique en cinq étapes : informations, secteur (cosmétiques ou vêtements), identité visuelle, carte bancaire, pièce d'identité et selfie ;
    \item Vérification KYC (\textit{Know Your Customer}) par l'administrateur Sellio, qui approuve ou refuse la boutique avec un motif ;
    \item Console d'administration : statistiques globales, liste paginée des boutiques et de leurs propriétaires, blocage d'un commerçant, suspension d'une boutique.
\end{itemize}

\subsection{Catalogue et vitrine personnalisable}
\begin{itemize}
    \item Gestion des produits, catégories et collections ;
    \item Attributs de variantes sous forme de listes normalisées : tailles, couleurs prédéfinies ou libres (RVB), tissus et autres attributs, extensibles par l'administrateur ;
    \item Trois gabarits de vitrine complets (\textbf{Minimal}, \textbf{Bold}, \textbf{Luxury}) avec menus, pied de page et page d'accueil propres ;
    \item Éditeur de thème : palette de couleurs, boutons, survols, barre de navigation, bandeau d'annonce, typographie, pied de page, appliqués en direct sur la vitrine.
\end{itemize}

\subsection{Parcours transactionnel}
\begin{itemize}
    \item Panier dynamique avec choix obligatoire de la taille pour les articles concernés ;
    \item Tunnel de commande avec zones de livraison, TVA et codes promotionnels ;
    \item Paiement en ligne sécurisé par \textbf{Stripe} ;
    \item Suivi des statuts de commande, historique client, demandes de retour et politique de retour ;
    \item Avis clients modérés par le commerçant.
\end{itemize}

\subsection{Marketing et fidélisation}
\begin{itemize}
    \item Promotions : codes promo et remises automatiques ;
    \item Bannières de page d'accueil ciblées par segment de clientèle ;
    \item Campagnes e-mail (newsletter, contact individuel) via l'API Brevo ;
    \item Programme de fidélité à points avec niveaux (Nouveau, Fidèle, VIP, Inactif) et montée de niveau automatique.
\end{itemize}

\subsection{Essayage virtuel et intelligence comportementale}
\begin{itemize}
    \item \textbf{Cabine d'essayage virtuel en temps réel} : le client active sa caméra et voit le vêtement porté sur lui, généré en direct par le modèle \textit{Lucy VTON} de l'API \textbf{Decart} ;
    \item \textbf{Suivi comportemental} : collecte des événements de navigation (vues, recherches, ajouts au panier, achats) dans une table partitionnée ;
    \item \textbf{Moteur de recommandation hybride} : « produits similaires », « pour vous », « souvent achetés ensemble » et « populaires » ;
    \item \textbf{Segmentation} des visiteurs par K-Means et \textbf{prédiction de l'attrition} des clients par forêt aléatoire ou gradient boosting selon le secteur ;
    \item Tableau de bord « Comportement » destiné au commerçant.
\end{itemize}

\subsection{Apports de la solution}
\begin{itemize}
    \item \textbf{Une plateforme pour plusieurs boutiques} : un seul code, une seule infrastructure, et une boutique opérationnelle en quelques minutes au lieu d'un développement sur mesure ;
    \item \textbf{Une confiance vérifiée} : aucun commerçant ne peut vendre avant la validation de son identité et de son moyen de paiement ; seules les informations non sensibles de la carte sont conservées, conformément à la norme PCI-DSS ;
    \item \textbf{Une isolation stricte des données} : chaque requête est rattachée à la boutique de l'utilisateur connecté, un commerçant ne pouvant consulter que les données de sa propre boutique ;
    \item \textbf{Une vitrine à l'image de la marque} sans écrire de code, grâce aux gabarits et à l'éditeur de thème ;
    \item \textbf{Une adaptation au contexte local} : prix en dinars, TVA et zones de livraison paramétrables ;
    \item \textbf{Une intelligence intégrée} : recommandation, segmentation et prédiction de l'attrition partagent les mêmes données que le reste de la plateforme, et chaque modèle a été évalué sur des données publiques réelles, avec comparaison à une référence simple ;
    \item \textbf{Une expérience d'achat différenciante} grâce à l'essayage virtuel, qui réduit l'incertitude sur le rendu d'un vêtement, principale cause d'hésitation et de retours dans la mode en ligne ;
    \item \textbf{Une architecture évolutive}, où chaque domaine (identité, catalogue, commandes, marketing, analyse) est un microservice qui peut évoluer et être déployé séparément.
\end{itemize}

\section{Étude des méthodologies de développement}
Avant de présenter les différentes approches, il convient de distinguer la \textbf{méthodologie}, qui définit les principes généraux du processus de développement, du \textbf{cadre de développement}, qui fournit une structure concrète (rôles, cérémonies, artefacts) pour appliquer ces principes.

\subsection{Choix de la méthodologie}
Le tableau~\ref{tab:comparaison-methodologies} compare les principales méthodologies de développement logiciel.
\begin{table}[H]
\centering
\renewcommand{\arraystretch}{1.4}
\begin{tabular}{|p{3cm}|p{6.2cm}|p{6.2cm}|}
\hline
\textbf{Méthodologie} & \textbf{Avantages} & \textbf{Inconvénients} \\
\hline
\textbf{Waterfall}
&
\textemdash\ Phases claires et structurées.\newline
\textemdash\ Planification simple du projet.
&
\textemdash\ Difficulté à s'adapter aux changements des besoins.\newline
\textemdash\ Retour d'information tardif.
\\
\hline
\textbf{RUP}
&
\textemdash\ Forte traçabilité.\newline
\textemdash\ Processus orienté vers la gestion des risques.
&
\textemdash\ Gestion complexe.\newline
\textemdash\ Charge documentaire importante.
\\
\hline
\textbf{2TUP}
&
\textemdash\ Approche itérative séparant branche fonctionnelle et branche technique.\newline
\textemdash\ Prise en compte des risques techniques.
&
\textemdash\ Moins flexible lorsque les besoins évoluent.\newline
\textemdash\ Approche davantage orientée processus.
\\
\hline
\textbf{Agile}
&
\textemdash\ Approche itérative et adaptative.\newline
\textemdash\ Retour d'information continu.\newline
\textemdash\ Bonne adaptation aux changements.
&
\textemdash\ Nécessite une communication régulière.\newline
\textemdash\ Dépend de l'implication des parties prenantes.
\\
\hline
\end{tabular}
\caption{Étude comparative des méthodologies de développement}
\label{tab:comparaison-methodologies}
\end{table}

L'\textbf{approche Agile} a été retenue. Le périmètre de Sellio a en effet beaucoup évolué pendant le stage : le projet a commencé comme une boutique unique (NaturEssence), a été découpé en microservices, puis transformé en plateforme multi-boutiques, avant d'intégrer l'essayage virtuel et l'apprentissage automatique. Seule une démarche itérative permettait d'absorber ces évolutions sans remettre en cause le travail déjà livré.

\subsection{Étude comparative des cadres de développement}
Le tableau~\ref{tab:comparaison-cadres} compare les principaux cadres Agile : \textbf{Scrum}, \textbf{Kanban} et \textbf{Extreme Programming (XP)}.
\begin{table}[H]
\centering
\renewcommand{\arraystretch}{1.4}
\begin{tabular}{|p{3cm}|p{6.2cm}|p{6.2cm}|}
\hline
\textbf{Cadre} & \textbf{Avantages} & \textbf{Inconvénients} \\
\hline
\textbf{Scrum}
&
\textemdash\ Organisation du travail en sprints.\newline
\textemdash\ Backlog permettant de prioriser les fonctionnalités.\newline
\textemdash\ Suivi régulier de l'avancement.\newline
\textemdash\ Rôles clairement définis.
&
\textemdash\ Nécessite une organisation et une communication régulières.\newline
\textemdash\ Moins adapté aux tâches très imprévisibles.
\\
\hline
\textbf{Kanban}
&
\textemdash\ Gestion visuelle des tâches.\newline
\textemdash\ Grande flexibilité des priorités.\newline
\textemdash\ Limitation du travail en cours.
&
\textemdash\ Moins structuré que Scrum.\newline
\textemdash\ Pas de rôles ni de sprints fixes.
\\
\hline
\textbf{XP}
&
\textemdash\ Forte orientation vers la qualité du code.\newline
\textemdash\ Tests fréquents et intégration continue.
&
\textemdash\ Nécessite une forte implication de l'équipe.\newline
\textemdash\ Pratiques (programmation en binôme) difficiles à appliquer par une seule développeuse.
\\
\hline
\end{tabular}
\caption{Étude comparative des cadres de développement Agile}
\label{tab:comparaison-cadres}
\end{table}

Nous avons retenu \textbf{Scrum}~\cite{ScrumGuide2020}. Son découpage en sprints de deux semaines correspond au rythme des points de suivi avec l'encadrant professionnel, et le backlog permet de prioriser les user stories d'un périmètre qui s'est élargi au fil du projet. Le projet a ainsi été organisé en six releases et onze sprints, présentés au chapitre suivant.\newline

\Needspace{12cm}
La figure~\ref{fig:fonctionnement-scrum} présente le fonctionnement du cadre Scrum, depuis la planification du sprint jusqu'à la revue et à la rétrospective.
\begin{figure}[H]
    \centering
    \img{0.9}{scrum-framework.png}
    \caption{Fonctionnement du cadre Scrum~\cite{ScrumGuide2020}}
    \label{fig:fonctionnement-scrum}
\end{figure}

\subsection{Les rôles dans la méthode Scrum}
Scrum repose sur trois rôles :\\

\noindent \textbf{- Product Owner :} il définit les besoins, priorise les fonctionnalités et maximise la valeur du produit.\\
\noindent \textbf{- Scrum Master :} facilitateur, il veille à la bonne application de Scrum, anime les cérémonies et lève les obstacles.\\
\noindent \textbf{- Équipe de développement :} elle réalise les tâches planifiées dans chaque sprint, de l'analyse jusqu'à la livraison.\\

Le tableau~\ref{tab:roles_scrum} présente les personnes qui ont occupé ces rôles dans notre projet.
\begin{table}[H]
\centering
\renewcommand{\arraystretch}{1.4}
\begin{tabular}{|p{5cm}|p{8cm}|}
\hline
\textbf{Rôle Scrum} & \textbf{Personne responsable} \\
\hline
Product Owner & M. Malek Naouar \\
\hline
Scrum Master & \textit{[Nom du Scrum Master]} \\
\hline
Équipe de développement & Rania Khedhri \\
\hline
\end{tabular}
\caption{Répartition des rôles Scrum dans notre projet}
\label{tab:roles_scrum}
\end{table}

\section*{Conclusion}
Dans ce chapitre, nous avons présenté l'organisme d'accueil et le contexte du projet, formulé la problématique, étudié les solutions existantes et décrit la solution proposée. Nous avons enfin justifié le choix de Scrum. Le chapitre suivant détaille l'analyse et la spécification des besoins.

\thispagestyle{fancy}
\fancyhead[R]{A.U 2025-2026}
\fancyhead[L]{ chapitre 2: Analyse et spécification des besoins }
\chapter{Analyse et spécification des besoins}
\section*{Introduction}
Dans ce chapitre, nous analysons en détail les besoins de la plateforme Sellio. Nous identifions d'abord les acteurs, puis les besoins fonctionnels et non fonctionnels. Nous présentons ensuite les diagrammes de cas d'utilisation, l'architecture retenue, le backlog produit et la planification des sprints, et enfin l'environnement de travail.

\section{Analyse et spécification des besoins}
La spécification des besoins fonctionnels et non fonctionnels est la première étape du cycle de développement : elle fixe ce que la plateforme doit offrir et les qualités qu'elle doit garantir.

\subsection{Identification des acteurs}
Sellio met en relation une plateforme, des commerçants et les clients de ces commerçants. Chaque acteur y joue un rôle distinct.
\subsubsection*{\underline{\textbf{Acteurs principaux}}}
\begin{itemize}[label=\ding{109}]
\item \textbf{Administrateur Sellio :} gère la plateforme dans son ensemble. Il consulte les statistiques globales, la liste des boutiques et de leurs propriétaires, valide ou refuse les dossiers de vérification (KYC), suspend une boutique ou bloque un commerçant, et enrichit les listes d'attributs partagées (tissus, etc.).
\item \textbf{Commerçant :} propriétaire d'une boutique. Il s'inscrit, crée sa boutique, soumet son dossier de vérification, puis gère depuis le back-office son catalogue, sa vitrine, ses commandes, ses retours, ses promotions, son programme de fidélité, ses campagnes e-mail et consulte l'analyse du comportement de ses visiteurs.
\item \textbf{Collaborateur de boutique :} membre de l'équipe d'un commerçant, doté d'un rôle aux permissions limitées (par exemple gestion des commandes uniquement).
\item \textbf{Visiteur :} internaute non connecté qui parcourt la vitrine d'une boutique, recherche des produits, reçoit des recommandations, remplit son panier et peut utiliser la cabine d'essayage virtuel.
\item \textbf{Client :} visiteur inscrit sur une boutique. En plus des actions du visiteur, il passe commande, paie en ligne, suit ses commandes, demande un retour, publie des avis et cumule des points de fidélité.
\end{itemize}
\subsubsection*{\underline{\textbf{Acteurs secondaires}}}
\begin{itemize}[label=\ding{109}]
\item \textbf{Stripe :} prestataire de paiement. Il traite les paiements des commandes (\textit{PaymentIntent}) et l'enregistrement de la carte du commerçant lors de la vérification (\textit{SetupIntent}).
\item \textbf{Decart :} service d'intelligence artificielle générative qui produit en temps réel la vidéo d'essayage virtuel à partir du flux caméra du client et de l'image du vêtement.
\item \textbf{Brevo~\cite{brevo} :} service d'envoi d'e-mails transactionnels et marketing (confirmations de commande, newsletters, contact client).
\item \textbf{Moteur d'apprentissage automatique :} ensemble de scripts Python (scikit-learn) appelés par le microservice d'analyse pour entraîner les modèles et produire recommandations, segments et scores d'attrition.
\end{itemize}

\subsection{Spécification des besoins}
Nous classons les besoins en deux familles : les besoins fonctionnels, qui décrivent \textbf{ce que} le système doit faire, et les besoins non fonctionnels, qui décrivent \textbf{comment} il doit le faire.
\subsubsection*{\underline{\textbf{Besoins fonctionnels}}}
Les besoins sont regroupés par module, ce qui reflète le découpage de la plateforme en microservices. Pour chaque fonctionnalité, nous précisons le ou les acteurs concernés.

\Needspace{20\baselineskip}
\paragraph{Module Plateforme, authentification et vérification}
Le tableau~\ref{tab:besoins-plateforme} présente les besoins de ce module.
\begin{table}[H]
\centering
\renewcommand{\arraystretch}{1.3}
\begin{tabular}{|p{9.5cm}|p{5.5cm}|}
\hline
\textbf{Fonctionnalité} & \textbf{Acteur(s) concerné(s)} \\
\hline
S'inscrire et s'authentifier de façon sécurisée (JWT et jeton de rafraîchissement) & Commerçant, Client, Administrateur Sellio \\
\hline
Créer une boutique (nom, secteur, identité visuelle) & Commerçant \\
\hline
Soumettre un dossier de vérification : carte bancaire, pièce d'identité (CIN ou passeport) et selfie & Commerçant \\
\hline
Valider ou refuser un dossier de vérification avec motif & Administrateur Sellio \\
\hline
Consulter les statistiques, les boutiques et leurs propriétaires & Administrateur Sellio \\
\hline
Suspendre une boutique, bloquer un commerçant & Administrateur Sellio \\
\hline
Gérer les rôles et permissions de l'équipe de la boutique & Commerçant \\
\hline
\end{tabular}
\caption{Besoins fonctionnels du module Plateforme, authentification et vérification}
\label{tab:besoins-plateforme}
\end{table}

\Needspace{20\baselineskip}
\paragraph{Module Catalogue et vitrine}
Le tableau~\ref{tab:besoins-catalogue} présente les besoins de ce module.
\begin{table}[H]
\centering
\renewcommand{\arraystretch}{1.3}
\begin{tabular}{|p{9.5cm}|p{5.5cm}|}
\hline
\textbf{Fonctionnalité} & \textbf{Acteur(s) concerné(s)} \\
\hline
Gérer les produits, leurs images et leurs variantes & Commerçant, Collaborateur \\
\hline
Choisir tailles, couleurs (prédéfinies ou RVB), tissus et autres attributs dans des listes normalisées & Commerçant, Collaborateur \\
\hline
Étendre les listes d'attributs & Administrateur Sellio, Commerçant \\
\hline
Gérer les catégories et les collections & Commerçant, Collaborateur \\
\hline
Choisir un gabarit de vitrine (Minimal, Bold, Luxury) et personnaliser le thème & Commerçant \\
\hline
Parcourir la vitrine, filtrer et rechercher les produits & Visiteur, Client \\
\hline
\end{tabular}
\caption{Besoins fonctionnels du module Catalogue et vitrine}
\label{tab:besoins-catalogue}
\end{table}

\Needspace{20\baselineskip}
\paragraph{Module Transactions}
Le tableau~\ref{tab:besoins-transactions} présente les besoins de ce module.
\begin{table}[H]
\centering
\renewcommand{\arraystretch}{1.3}
\begin{tabular}{|p{9.5cm}|p{5.5cm}|}
\hline
\textbf{Fonctionnalité} & \textbf{Acteur(s) concerné(s)} \\
\hline
Gérer un panier (taille obligatoire pour les articles concernés) & Visiteur, Client \\
\hline
Passer commande avec livraison, TVA et code promo & Client \\
\hline
Payer en ligne par carte (Stripe) & Client, Stripe \\
\hline
Suivre ses commandes et demander un retour & Client \\
\hline
Gérer les commandes, leurs statuts et les retours & Commerçant, Collaborateur \\
\hline
Paramétrer la TVA, les zones de livraison et la politique de retour & Commerçant \\
\hline
Publier un avis ; modérer et répondre aux avis & Client ; Commerçant \\
\hline
\end{tabular}
\caption{Besoins fonctionnels du module Transactions}
\label{tab:besoins-transactions}
\end{table}

\Needspace{20\baselineskip}
\paragraph{Module Marketing et fidélisation}
Le tableau~\ref{tab:besoins-marketing} présente les besoins de ce module.
\begin{table}[H]
\centering
\renewcommand{\arraystretch}{1.3}
\begin{tabular}{|p{9.5cm}|p{5.5cm}|}
\hline
\textbf{Fonctionnalité} & \textbf{Acteur(s) concerné(s)} \\
\hline
Créer des codes promo et des remises automatiques & Commerçant \\
\hline
Publier des bannières ciblées par segment & Commerçant \\
\hline
Envoyer une newsletter ou contacter un client par e-mail & Commerçant, Brevo \\
\hline
Configurer le programme de fidélité et ses niveaux & Commerçant \\
\hline
Consulter son solde de points et son niveau & Client \\
\hline
\end{tabular}
\caption{Besoins fonctionnels du module Marketing et fidélisation}
\label{tab:besoins-marketing}
\end{table}

\Needspace{20\baselineskip}
\paragraph{Module Intelligence (essayage virtuel et apprentissage automatique)}
Le tableau~\ref{tab:besoins-ia} présente les besoins de ce module.
\begin{table}[H]
\centering
\renewcommand{\arraystretch}{1.3}
\begin{tabular}{|p{9.5cm}|p{5.5cm}|}
\hline
\textbf{Fonctionnalité} & \textbf{Acteur(s) concerné(s)} \\
\hline
Essayer virtuellement un vêtement en temps réel avec sa caméra & Visiteur, Client, Decart \\
\hline
Collecter les événements de navigation & Visiteur, Client (implicite) \\
\hline
Recevoir des recommandations : produits similaires, pour vous, souvent achetés ensemble, populaires & Visiteur, Client, Moteur ML \\
\hline
Consulter l'analyse comportementale : entonnoir, segments de visiteurs, produits consultés & Commerçant \\
\hline
Identifier les clients à risque d'attrition & Commerçant, Moteur ML \\
\hline
Ré-entraîner les modèles de la boutique & Commerçant, Moteur ML \\
\hline
\end{tabular}
\caption{Besoins fonctionnels du module Intelligence}
\label{tab:besoins-ia}
\end{table}

\subsubsection*{\underline{\textbf{Besoins non fonctionnels}}}
Les besoins non fonctionnels définissent les qualités attendues du système. \newline
\Needspace{20\baselineskip}
Le tableau~\ref{tab:exigences-non-fonctionnelles} les présente selon le modèle de qualité ISO/IEC 25010~\cite{ISO25010}.

\begin{longtable}{|p{3.2cm}|p{4cm}|p{8cm}|}
\caption{Exigences non fonctionnelles du système}
\label{tab:exigences-non-fonctionnelles} \\
\hline
\textbf{Caractéristique} & \textbf{Sous-caractéristique} & \textbf{Description} \\
\hline
\endfirsthead
\hline
\textbf{Caractéristique} & \textbf{Sous-caractéristique} & \textbf{Description} \\
\hline
\endhead
\hline \multicolumn{3}{r}{\textit{Suite à la page suivante}} \\
\endfoot
\hline
\endlastfoot

\multirow{4}{*}{Sécurité} & Authentification & Jeton d'accès JWT (24 h) et jeton de rafraîchissement (7 jours) ; mots de passe hachés avec BCrypt. \\
\cline{2-3}
 & Contrôle d'accès & Rôles et permissions vérifiés côté serveur ; un commerçant n'accède qu'aux données de sa propre boutique. \\
\cline{2-3}
 & Protection des paiements & Aucune donnée de carte ne transite par nos serveurs ; seuls la marque, les 4 derniers chiffres, l'expiration et la référence Stripe sont stockés (PCI-DSS). \\
\cline{2-3}
 & Confiance & Une boutique reste « en attente » tant que son dossier KYC n'est pas validé ; un compte bloqué ne peut plus se connecter. \\
\hline
\multirow{2}{*}{Performance} & Temps de réponse & Recommandations servies depuis des modèles pré-entraînés ; les événements sont envoyés par lots pour ne pas ralentir la navigation. \\
\cline{2-3}
 & Volume de données & Table d'événements partitionnée par mois, index BRIN sur la date, agrégats quotidiens et rétention de 13 mois. \\
\hline
\multirow{2}{*}{Fiabilité} & Tolérance aux pannes & La panne d'un microservice (par exemple l'analyse) n'empêche pas les achats ; la vitrine masque simplement les sections de recommandations. \\
\cline{2-3}
 & Intégrité & Transactions ACID pour les commandes et les paiements. \\
\hline
\multirow{2}{*}{Utilisabilité} & Ergonomie & Interfaces responsives ; thème sombre par défaut sur l'espace Sellio avec bascule vers le thème clair. \\
\cline{2-3}
 & Personnalisation & Toute modification du thème est visible immédiatement sur la vitrine, sans écrire de code. \\
\hline
\multirow{2}{*}{Maintenabilité} & Modularité & Un microservice par domaine métier, une bibliothèque partagée pour les entités et DTO. \\
\cline{2-3}
 & Évolutivité & Ajout d'un gabarit de vitrine, d'un attribut ou d'un modèle ML sans toucher aux autres services. \\
\hline
Portabilité & Compatibilité & Navigateurs modernes (Chrome, Firefox, Edge) ; l'essayage virtuel exige une caméra et WebRTC. \\
\hline
Confidentialité & Données personnelles & Le suivi comportemental repose sur un identifiant de visiteur anonyme ; les pièces d'identité ne sont visibles que par l'administrateur Sellio. \\
\hline
\end{longtable}

\subsection{Diagramme de cas d'utilisation général}
Le système est utilisé par cinq acteurs principaux : \textbf{l'Administrateur Sellio}, \textbf{le Commerçant}, \textbf{le Collaborateur de boutique}, \textbf{le Visiteur} et \textbf{le Client}. Les sous-sections suivantes présentent les cas d'utilisation de chacun.

\subsubsection{Cas d'utilisation de l'Administrateur Sellio}
La figure~\ref{fig:cu-admin} présente les cas d'utilisation de l'administrateur de la plateforme.
\begin{figure}[H]
    \centering
    \img{0.9}{cu_admin_sellio.png}
    \caption{Diagramme de cas d'utilisation de l'Administrateur Sellio}
    \label{fig:cu-admin}
\end{figure}

\subsubsection{Cas d'utilisation du Commerçant}
La figure~\ref{fig:cu-commercant} présente les cas d'utilisation du commerçant. Le collaborateur de boutique hérite d'une partie de ces cas, selon les permissions de son rôle.
\begin{figure}[H]
    \centering
    \img{1}{cu_commercant.png}
    \caption{Diagramme de cas d'utilisation du Commerçant}
    \label{fig:cu-commercant}
\end{figure}

\subsubsection{Cas d'utilisation du Visiteur et du Client}
La figure~\ref{fig:cu-client} présente les cas d'utilisation du visiteur et du client. Le client hérite de tous les cas du visiteur.
\begin{figure}[H]
    \centering
    \img{0.9}{cu_client.png}
    \caption{Diagramme de cas d'utilisation du Visiteur et du Client}
    \label{fig:cu-client}
\end{figure}

\section{Analyse globale}
\subsection{Diagramme de classes d'analyse}
La figure~\ref{fig:classes} présente le diagramme de classes d'analyse. La classe centrale est \texttt{Shop} : chaque utilisateur, produit, catégorie, commande, promotion ou événement comportemental est rattaché à une boutique, ce qui matérialise l'isolation multi-boutiques. Une boutique possède un dossier de vérification (\texttt{MerchantVerification}), un statut (\texttt{PENDING}, \texttt{ACTIVE}, \texttt{REJECTED}, \texttt{SUSPENDED}), un gabarit et un thème.
\begin{figure}[H]
    \centering
    \img{1}{diagramme_classes.png}
    \caption{Diagramme de classes d'analyse}
    \label{fig:classes}
\end{figure}

\subsection{Architecture globale}
L'architecture influence directement les performances, la maintenabilité et l'évolutivité de la plateforme. Trois approches ont été comparées : \textbf{l'architecture monolithique}, \textbf{l'architecture N-tiers} et \textbf{l'architecture microservices}.\newline
\Needspace{20\baselineskip}
Le tableau~\ref{tab:comparaison_architectures} présente leurs avantages et leurs limites.

\begin{longtable}{|>{\raggedright\arraybackslash}p{3.2cm}|>{\raggedright\arraybackslash}p{11cm}|}
\hline
\textbf{Architecture} & \textbf{Caractéristiques} \\
\hline
\endfirsthead
\hline
\textbf{Architecture} & \textbf{Caractéristiques} \\
\hline
\endhead
Monolithique &
\begin{itemize}[leftmargin=*, nosep]
    \item Présentation, logique métier et accès aux données dans une seule unité de déploiement.
    \item \textbf{Avantages} : simplicité de développement et de déploiement, débogage facile, cohérence transactionnelle native.
    \item \textbf{Limites} : mise à l'échelle uniquement globale, un défaut dans un module peut affecter tout le système, maintenance difficile quand le code grossit.
\end{itemize} \\
\hline
Client--serveur N-tiers &
\begin{itemize}[leftmargin=*, nosep]
    \item Séparation en couches : présentation, logique métier, données.
    \item \textbf{Avantages} : séparation claire des responsabilités, plusieurs clients peuvent consommer la même couche métier.
    \item \textbf{Limites} : le serveur d'application reste une unité de déploiement unique, avec un couplage interne entre modules.
\end{itemize} \\
\hline
Microservices &
\begin{itemize}[leftmargin=*, nosep]
    \item Décomposition en services indépendants, chacun responsable d'un domaine métier, déployés séparément et communiquant par API REST derrière une passerelle.
    \item \textbf{Avantages} : mise à l'échelle service par service, isolation des pannes, déploiements indépendants, liberté technologique par service.
    \item \textbf{Limites} : complexité opérationnelle, latence réseau entre services, cohérence des données distribuées plus difficile.
\end{itemize} \\
\hline
\caption{Comparaison des architectures logicielles étudiées}
\label{tab:comparaison_architectures}
\end{longtable}

L'\textbf{architecture microservices} a été retenue, pour trois raisons. D'abord, le sujet lui-même demande une plateforme distribuée. Ensuite, les domaines de Sellio ont des charges très différentes : le service d'analyse reçoit un flux continu d'événements et lance des entraînements de modèles coûteux, alors que le service d'identité traite peu de requêtes ; les séparer permet de dimensionner chacun selon sa charge et d'éviter qu'un entraînement ne ralentisse les achats. Enfin, le projet est né d'un monolithe (NaturEssence) dont la croissance rendait les évolutions de plus en plus risquées ; la migration vers des services séparés a rendu chaque domaine plus facile à faire évoluer.

Par pragmatisme, les services partagent une même instance PostgreSQL, chaque service restant propriétaire de ses tables, et une bibliothèque commune (\texttt{shared-lib}) regroupe les entités et les DTO. Ce compromis, adapté à une équipe réduite, est discuté dans les perspectives (base de données par service).

\subsubsection{Choix du système de gestion de base de données}
Le plan de travail initial prévoyait MongoDB. Le tableau~\ref{tab:choix-sgbd} résume la comparaison qui nous a conduits à retenir PostgreSQL.
\begin{table}[H]
\centering
\renewcommand{\arraystretch}{1.4}
\begin{tabular}{|p{4cm}|p{5.3cm}|p{5.3cm}|}
\hline
\textbf{Critère} & \textbf{MongoDB} & \textbf{PostgreSQL} \\
\hline
Transactions (commande + paiement + stock) & Multi-documents possibles mais plus coûteuses & ACID natives \\
\hline
Intégrité référentielle (boutique, client, commande) & À gérer dans l'application & Clés étrangères et contraintes \\
\hline
Données flexibles (attributs, thème) & Natif & Type JSONB \\
\hline
Volume d'événements & Bon & Partitionnement déclaratif et index BRIN \\
\hline
Requêtes analytiques (agrégats, jointures) & Pipeline d'agrégation & SQL standard, fonctions de fenêtre \\
\hline
\end{tabular}
\caption{Comparaison MongoDB / PostgreSQL}
\label{tab:choix-sgbd}
\end{table}
Les données de Sellio sont avant tout relationnelles et transactionnelles (une commande relie un client, des produits, une boutique et un paiement) ; PostgreSQL couvre en outre le besoin de flexibilité (JSONB) et de volume (partitionnement), d'où ce choix~\cite{postgresql2026}.

\subsubsection{Architecture logique de la plateforme}
La figure~\ref{fig:archi_logique} présente l'architecture logique de Sellio et ses couches.
\begin{figure}[H]
\centering
\img{1}{architecture_logique.png}
\caption{Architecture logique de la plateforme Sellio}
\label{fig:archi_logique}
\end{figure}

\subsubsection*{La couche présentation (front-end) :}
Elle se compose de deux applications React :
\begin{itemize}
  \item \textbf{l'application Sellio} (port 3000), qui regroupe le site de la plateforme (accueil, inscription, connexion), la console de l'administrateur Sellio et le back-office de chaque commerçant ;
  \item \textbf{la vitrine} (port 3001), commune à toutes les boutiques : la boutique affichée est déterminée par son identifiant (\textit{slug}) et son gabarit, sa palette et son thème sont chargés dynamiquement.
\end{itemize}

\subsubsection*{La couche passerelle :}
Une passerelle \textbf{Spring Cloud Gateway} (port 8080) est le point d'entrée unique des API. Elle route chaque requête vers le bon microservice selon son chemin (\texttt{/api/v1/auth/**} vers le service d'identité, \texttt{/api/v1/public/products/**} vers le catalogue, etc.).

\subsubsection*{La couche métier (back-end) :}
Elle est composée de cinq microservices Spring Boot :
\begin{itemize}
  \item \textbf{auth-service} (8081) : comptes, JWT, rôles et permissions, boutiques, vérification KYC, console plateforme, fidélité et segments ;
  \item \textbf{catalog-service} (8082) : produits, variantes, catégories, collections, avis, téléversement d'images et jeton d'essayage virtuel ;
  \item \textbf{order-service} (8083) : tunnel de commande, paiement Stripe, commandes, retours, TVA, livraison, promotions et tableau de bord ;
  \item \textbf{marketing-service} (8084) : bannières, apparence de la vitrine, profil de la boutique et e-mailing ;
  \item \textbf{analytics-service} (8085) : collecte des événements, recommandations, analyse comportementale et prédiction de l'attrition.
\end{itemize}

\subsubsection*{La couche intelligence :}
Des scripts \textbf{Python} (pandas, scikit-learn) entraînent et appliquent les modèles. Le service d'analyse les exécute comme sous-processus et stocke les modèles entraînés par boutique.

\subsubsection*{La couche données :}
\textbf{PostgreSQL} stocke l'ensemble des données ; les images de produits sont stockées sur disque et servies par le catalogue.

\subsubsection*{Les services externes :}
\textbf{Stripe} (paiements), \textbf{Decart} (essayage virtuel) et \textbf{Brevo} (e-mails).

\subsubsection{Architecture physique de la plateforme}
La figure~\ref{fig:archi_physique} présente le déploiement physique : le navigateur du client communique en HTTPS avec les serveurs web des applications React et avec la passerelle ; celle-ci relaie les appels vers les microservices, qui accèdent au serveur de base de données. Le navigateur communique directement avec Stripe (saisie de carte) et avec Decart (flux vidéo WebRTC), les jetons nécessaires étant délivrés par nos services.
\begin{figure}[H]
    \centering
    \img{1}{architecture_physique.png}
    \caption{Architecture physique de la plateforme Sellio}
    \label{fig:archi_physique}
\end{figure}

\section{Pilotage du projet avec Scrum}
\subsection{Backlog du produit}
Le tableau~\ref{tab:backlog-produit} présente le backlog du produit Sellio. Les fonctionnalités sont organisées par module et décrites sous forme de user stories.
\small
\begin{longtable}{|p{0.9cm}|p{3cm}|p{1.3cm}|p{6.6cm}|p{1.9cm}|}
\hline
\textbf{ID} & \textbf{Module} & \textbf{ID US} & \textbf{User Story} & \textbf{Complexité} \\
\hline
\endfirsthead
\hline
\textbf{ID} & \textbf{Module} & \textbf{ID US} & \textbf{User Story} & \textbf{Complexité} \\
\hline
\endhead
\hline \multicolumn{5}{r}{\textit{Suite à la page suivante}} \\
\endfoot
\hline
\endlastfoot

\multirow{5}{*}{M1} & \multirow{5}{3cm}{Authentification} & M1.1 & En tant qu'utilisateur, je souhaite créer un compte. & Moyenne \\
\cline{3-5}
 & & M1.2 & En tant qu'utilisateur, je souhaite me connecter. & Moyenne \\
\cline{3-5}
 & & M1.3 & En tant qu'utilisateur, je souhaite rester connecté grâce au rafraîchissement de mon jeton. & Moyenne \\
\cline{3-5}
 & & M1.4 & En tant qu'utilisateur, je souhaite me déconnecter. & Faible \\
\cline{3-5}
 & & M1.5 & En tant que client, je souhaite consulter et modifier mon profil. & Faible \\
\hline
\multirow{3}{*}{M2} & \multirow{3}{3cm}{Rôles et utilisateurs} & M2.1 & En tant que commerçant, je souhaite créer des rôles et leur associer des permissions. & Moyenne \\
\cline{3-5}
 & & M2.2 & En tant que commerçant, je souhaite gérer les comptes de mon équipe et de mes clients. & Moyenne \\
\cline{3-5}
 & & M2.3 & En tant que système, je veux restreindre chaque endpoint selon les permissions. & Élevée \\
\hline
\multirow{4}{*}{M3} & \multirow{4}{3cm}{Plateforme et boutiques} & M3.1 & En tant que visiteur, je souhaite découvrir Sellio sur une page d'accueil. & Faible \\
\cline{3-5}
 & & M3.2 & En tant que commerçant, je souhaite m'inscrire sur Sellio. & Moyenne \\
\cline{3-5}
 & & M3.3 & En tant que commerçant, je souhaite créer ma boutique avec un assistant. & Élevée \\
\cline{3-5}
 & & M3.4 & En tant que commerçant, je souhaite ouvrir ma vitrine depuis le back-office. & Faible \\
\hline
\multirow{5}{*}{M4} & \multirow{5}{3cm}{Vérification KYC et console} & M4.1 & En tant que commerçant, je souhaite enregistrer ma carte bancaire. & Élevée \\
\cline{3-5}
 & & M4.2 & En tant que commerçant, je souhaite soumettre ma pièce d'identité et un selfie. & Moyenne \\
\cline{3-5}
 & & M4.3 & En tant qu'administrateur Sellio, je souhaite valider ou refuser un dossier. & Moyenne \\
\cline{3-5}
 & & M4.4 & En tant qu'administrateur Sellio, je souhaite consulter les boutiques et leurs propriétaires. & Moyenne \\
\cline{3-5}
 & & M4.5 & En tant qu'administrateur Sellio, je souhaite suspendre une boutique ou bloquer un commerçant. & Moyenne \\
\hline
\multirow{3}{*}{M5} & \multirow{3}{3cm}{Catalogue} & M5.1 & En tant que commerçant, je souhaite gérer mes produits. & Élevée \\
\cline{3-5}
 & & M5.2 & En tant que commerçant, je souhaite gérer mes catégories. & Moyenne \\
\cline{3-5}
 & & M5.3 & En tant que commerçant, je souhaite gérer mes collections. & Moyenne \\
\hline
\multirow{3}{*}{M6} & \multirow{3}{3cm}{Attributs de variantes} & M6.1 & En tant que commerçant, je souhaite choisir les tailles dans une liste. & Moyenne \\
\cline{3-5}
 & & M6.2 & En tant que commerçant, je souhaite choisir des couleurs prédéfinies ou libres (RVB). & Moyenne \\
\cline{3-5}
 & & M6.3 & En tant que commerçant, je souhaite choisir le tissu et d'autres attributs, et en ajouter. & Moyenne \\
\hline
\multirow{3}{*}{M7} & \multirow{3}{3cm}{Vitrine} & M7.1 & En tant que commerçant, je souhaite choisir un gabarit Minimal, Bold ou Luxury. & Élevée \\
\cline{3-5}
 & & M7.2 & En tant que visiteur, je souhaite parcourir la vitrine, ses menus et ses catégories. & Élevée \\
\cline{3-5}
 & & M7.3 & En tant que visiteur, je souhaite consulter la fiche d'un produit. & Moyenne \\
\hline
\multirow{2}{*}{M8} & \multirow{2}{3cm}{Personnalisation} & M8.1 & En tant que commerçant, je souhaite choisir la palette de couleurs de ma boutique. & Moyenne \\
\cline{3-5}
 & & M8.2 & En tant que commerçant, je souhaite personnaliser boutons, survols, navigation, textes et pied de page. & Élevée \\
\hline
\multirow{3}{*}{M9} & \multirow{3}{3cm}{Panier et commande} & M9.1 & En tant que visiteur, je souhaite gérer mon panier en choisissant la taille. & Moyenne \\
\cline{3-5}
 & & M9.2 & En tant que client, je souhaite passer commande avec livraison et code promo. & Élevée \\
\cline{3-5}
 & & M9.3 & En tant que client, je souhaite retrouver mon panier sur un autre appareil. & Faible \\
\hline
M10 & Paiement & M10.1 & En tant que client, je souhaite payer ma commande par carte. & Élevée \\
\hline
\multirow{3}{*}{M11} & \multirow{3}{3cm}{Commandes et tableau de bord} & M11.1 & En tant que commerçant, je souhaite gérer mes commandes et leurs statuts. & Moyenne \\
\cline{3-5}
 & & M11.2 & En tant que client, je souhaite suivre mes commandes. & Faible \\
\cline{3-5}
 & & M11.3 & En tant que commerçant, je souhaite consulter un tableau de bord des ventes. & Moyenne \\
\hline
\multirow{3}{*}{M12} & \multirow{3}{3cm}{Retours, TVA et livraison} & M12.1 & En tant que client, je souhaite demander un retour. & Moyenne \\
\cline{3-5}
 & & M12.2 & En tant que commerçant, je souhaite traiter les retours et définir ma politique de retour. & Moyenne \\
\cline{3-5}
 & & M12.3 & En tant que commerçant, je souhaite paramétrer la TVA et les zones de livraison. & Moyenne \\
\hline
\multirow{2}{*}{M13} & \multirow{2}{3cm}{Avis} & M13.1 & En tant que client, je souhaite publier un avis sur un produit acheté. & Faible \\
\cline{3-5}
 & & M13.2 & En tant que commerçant, je souhaite modérer les avis et y répondre. & Faible \\
\hline
\multirow{2}{*}{M14} & \multirow{2}{3cm}{Promotions} & M14.1 & En tant que commerçant, je souhaite créer des codes promo. & Moyenne \\
\cline{3-5}
 & & M14.2 & En tant que commerçant, je souhaite créer des remises automatiques. & Moyenne \\
\hline
\multirow{2}{*}{M15} & \multirow{2}{3cm}{Bannières et e-mailing} & M15.1 & En tant que commerçant, je souhaite publier des bannières ciblées. & Moyenne \\
\cline{3-5}
 & & M15.2 & En tant que commerçant, je souhaite envoyer des e-mails à mes clients. & Moyenne \\
\hline
\multirow{3}{*}{M16} & \multirow{3}{3cm}{Fidélité et segments} & M16.1 & En tant que commerçant, je souhaite configurer le gain de points. & Moyenne \\
\cline{3-5}
 & & M16.2 & En tant que commerçant, je souhaite gérer les niveaux et ajuster les points d'un client. & Moyenne \\
\cline{3-5}
 & & M16.3 & En tant que client, je souhaite consulter mes points et mon niveau. & Faible \\
\hline
\multirow{2}{*}{M17} & \multirow{2}{3cm}{Essayage virtuel} & M17.1 & En tant que visiteur, je souhaite essayer un vêtement en direct avec ma caméra. & Élevée \\
\cline{3-5}
 & & M17.2 & En tant que système, je veux délivrer des jetons Decart à durée limitée sans exposer la clé. & Moyenne \\
\hline
M18 & Suivi comportemental & M18.1 & En tant que système, je veux collecter les événements de navigation. & Élevée \\
\hline
\multirow{2}{*}{M19} & \multirow{2}{3cm}{Recommandation} & M19.1 & En tant que visiteur, je souhaite voir des produits similaires et « pour vous ». & Élevée \\
\cline{3-5}
 & & M19.2 & En tant que visiteur, je souhaite voir les produits souvent achetés ensemble. & Moyenne \\
\hline
\multirow{3}{*}{M20} & \multirow{3}{3cm}{Analyse comportementale et attrition} & M20.1 & En tant que commerçant, je souhaite voir les segments de mes visiteurs. & Élevée \\
\cline{3-5}
 & & M20.2 & En tant que commerçant, je souhaite identifier les clients à risque d'attrition. & Élevée \\
\cline{3-5}
 & & M20.3 & En tant que commerçant, je souhaite consulter le tableau de bord « Comportement » et ré-entraîner les modèles. & Moyenne \\
\hline
\caption{Backlog du produit Sellio}
\label{tab:backlog-produit}
\end{longtable}
\normalsize

\subsection{Planification des sprints}
Après l'élaboration du backlog, le projet a été réparti en \textbf{six releases} et \textbf{onze sprints} de deux semaines. Les deux premières semaines du stage (du 2 au 13 février 2026) ont été consacrées à l'étude du projet NaturEssence existant, à l'étude de l'existant et à la préparation du backlog ; les dernières semaines, à la consolidation, aux tests et à la rédaction de ce rapport.\newline
Le tableau~\ref{tab:planification-sprints} présente la planification globale des sprints.
\begin{longtable}{|p{1.6cm}|p{4.2cm}|p{6.5cm}|p{1.6cm}|}
\hline
\textbf{Sprint} & \textbf{Modules développés} & \textbf{Objectif du sprint} & \textbf{Charge (pts)} \\
\hline
\endfirsthead
\hline
\textbf{Sprint} & \textbf{Modules développés} & \textbf{Objectif du sprint} & \textbf{Charge (pts)} \\
\hline
\endhead
\hline \multicolumn{4}{r}{\textit{Suite à la page suivante}} \\
\endfoot
\hline
\endlastfoot
\multicolumn{4}{|c|}{\textbf{Release 1 : Fondations de la plateforme}} \\
\hline
\textbf{Sprint 1} & M1 -- Authentification \newline M2 -- Rôles et utilisateurs & Découper le projet en microservices derrière une passerelle et sécuriser l'accès (JWT, rôles, permissions). & 22 \\
\hline
\textbf{Sprint 2} & M3 -- Plateforme et boutiques \newline M4 -- Vérification KYC et console & Transformer la boutique unique en plateforme multi-boutiques : inscription des commerçants, création de boutique, vérification KYC, console d'administration. & 25 \\
\hline
\multicolumn{4}{|c|}{\textbf{Release 2 : Catalogue et vitrine}} \\
\hline
\textbf{Sprint 3} & M5 -- Catalogue \newline M6 -- Attributs de variantes & Offrir un catalogue adapté aux cosmétiques et aux vêtements, avec des attributs normalisés. & 19 \\
\hline
\textbf{Sprint 4} & M7 -- Vitrine \newline M8 -- Personnalisation & Livrer trois gabarits de vitrine complets et un éditeur de thème appliqué en direct. & 18 \\
\hline
\multicolumn{4}{|c|}{\textbf{Release 3 : Parcours transactionnel}} \\
\hline
\textbf{Sprint 5} & M9 -- Panier et commande \newline M10 -- Paiement & Permettre d'acheter : panier, tunnel de commande et paiement Stripe. & 17 \\
\hline
\textbf{Sprint 6} & M11 -- Commandes et tableau de bord \newline M12 -- Retours, TVA et livraison \newline M13 -- Avis & Gérer l'après-vente et le paramétrage fiscal et logistique. & 17 \\
\hline
\multicolumn{4}{|c|}{\textbf{Release 4 : Marketing et fidélisation}} \\
\hline
\textbf{Sprint 7} & M14 -- Promotions \newline M15 -- Bannières et e-mailing & Outiller le commerçant pour attirer et relancer ses clients. & 15 \\
\hline
\textbf{Sprint 8} & M16 -- Fidélité et segments & Fidéliser les clients par un programme à points et des niveaux. & 11 \\
\hline
\multicolumn{4}{|c|}{\textbf{Release 5 : Essayage virtuel}} \\
\hline
\textbf{Sprint 9} & M17 -- Essayage virtuel & Intégrer la cabine d'essayage en temps réel fondée sur l'API Decart. & 12 \\
\hline
\multicolumn{4}{|c|}{\textbf{Release 6 : Intelligence comportementale}} \\
\hline
\textbf{Sprint 10} & M18 -- Suivi comportemental \newline M19 -- Recommandation & Collecter les événements et servir des recommandations hybrides. & 15 \\
\hline
\textbf{Sprint 11} & M20 -- Analyse comportementale et attrition & Segmenter les visiteurs, prédire l'attrition et restituer l'analyse au commerçant. & 16 \\
\hline
\multicolumn{3}{|r|}{\textbf{Total}} & \textbf{187} \\
\hline
\caption{Planification des sprints -- Sellio}
\label{tab:planification-sprints}
\end{longtable}

\section{Environnement de travail}
\subsection{Environnement matériel de développement}
Le développement a été réalisé sur un ordinateur portable présentant les caractéristiques suivantes :
\begin{itemize}
    \item \textbf{Modèle :} \textit{[modèle de l'ordinateur]}
    \item \textbf{Système d'exploitation :} Windows 11 Professionnel (64 bits)
    \item \textbf{Processeur :} \textit{[processeur]}
    \item \textbf{Mémoire RAM :} \textit{[mémoire]}
    \item \textbf{Stockage :} \textit{[stockage]}
\end{itemize}

\subsection{Environnement logiciel de développement}
\Needspace{20\baselineskip}
Le tableau~\ref{tab:outils} présente les principaux outils utilisés.
\vspace{0.5em}
\renewcommand{\arraystretch}{1.3}
\begin{center}
\begin{tabular}{|l|p{9cm}|c|}
\hline
\textbf{Outil} & \textbf{Description} & \textbf{Logo} \\
\hline
Visual Studio Code~\cite{vscode} & Éditeur principal pour le front-end React, le back-end Spring Boot et les scripts Python. & \logo{25pt}{vscode.jpeg} \\
\hline
Git et GitHub~\cite{github} & Gestion de versions et hébergement du dépôt du projet. & \logo{25pt}{github.png} \\
\hline
PlantUML~\cite{plantuml} & Génération des diagrammes UML à partir de descriptions textuelles. & \logo{25pt}{plantuml.png} \\
\hline
PostgreSQL~\cite{postgresql2026} & Système de gestion de base de données relationnelle. & \logo{25pt}{postgresql.png} \\
\hline
Maven~\cite{maven} & Construction et gestion des dépendances des microservices Java. & \logo{25pt}{maven.png} \\
\hline
Stripe Dashboard~\cite{stripe} & Suivi des paiements et des cartes en mode test. & \logo{25pt}{stripe.png} \\
\hline
\end{tabular}
\captionof{table}{Outils utilisés pour la modélisation, le développement et le suivi}
\label{tab:outils}
\end{center}

\subsubsection{Langages de programmation}
\Needspace{16\baselineskip}
Le tableau~\ref{tab:langages} présente les langages utilisés.
\vspace{0.5em}
\begin{center}
\begin{tabular}{|l|p{9cm}|c|}
\hline
\textbf{Langage} & \textbf{Usage} & \textbf{Logo} \\
\hline
Java 17~\cite{java17} & Microservices (API REST, logique métier, sécurité). & \logo{25pt}{java.png} \\
\hline
JavaScript (JSX)~\cite{javascript} & Applications React : Sellio, back-office et vitrine. & \logo{25pt}{javascript.png} \\
\hline
Python 3~\cite{python} & Nettoyage des données, entraînement et évaluation des modèles. & \logo{25pt}{python.png} \\
\hline
SQL & Requêtes, partitionnement et agrégats PostgreSQL. & \logo{25pt}{sql.png} \\
\hline
\end{tabular}
\captionof{table}{Langages de programmation utilisés}
\label{tab:langages}
\end{center}

\subsubsection{Frameworks et bibliothèques}
\Needspace{24\baselineskip}
Le tableau~\ref{tab:frameworks} présente les frameworks et bibliothèques utilisés.
\vspace{0.5em}
\begin{center}
\begin{tabular}{|l|p{9cm}|c|}
\hline
\textbf{Framework} & \textbf{Rôle} & \textbf{Logo} \\
\hline
Spring Boot 3.2~\cite{springboot} & Socle des microservices (Spring Web, Spring Data JPA, Spring Security). & \logo{25pt}{springboot.png} \\
\hline
Spring Cloud Gateway~\cite{springcloudgateway} & Passerelle API : point d'entrée unique et routage vers les services. & \logo{25pt}{springcloud.png} \\
\hline
React 18~\cite{react} et Vite~\cite{vite} & Interfaces web à base de composants ; serveur de développement et compilation. & \logo{25pt}{react.png} \\
\hline
Tailwind CSS~\cite{tailwind} & Mise en forme utilitaire ; variables CSS pour les thèmes des vitrines. & \logo{25pt}{tailwind.png} \\
\hline
scikit-learn~\cite{scikitlearn}, pandas, NumPy & Modèles d'apprentissage automatique et préparation des données. & \logo{25pt}{scikitlearn.png} \\
\hline
Stripe.js~\cite{stripe} & Saisie sécurisée de la carte et confirmation des paiements. & \logo{25pt}{stripe.png} \\
\hline
SDK Decart~\cite{decart} & Connexion WebRTC au modèle d'essayage virtuel en temps réel. & \logo{25pt}{decart.png} \\
\hline
\end{tabular}
\captionof{table}{Frameworks et bibliothèques utilisés}
\label{tab:frameworks}
\end{center}

\section*{Conclusion}
Dans ce chapitre, nous avons identifié les acteurs, spécifié les besoins fonctionnels et non fonctionnels, justifié l'architecture microservices et le choix de PostgreSQL, puis présenté le backlog produit et la planification en six releases et onze sprints. Le chapitre suivant présente la première release.

\thispagestyle{fancy}
\fancyhead[R]{A.U 2025-2026}
\fancyhead[L]{ chapitre 3: Release 1 : Fondations de la plateforme }
\chapter{Release 1 : Fondations de la plateforme Sellio}
\section*{Introduction}
Ce chapitre présente la première release, qui pose les fondations de Sellio. Le premier sprint découpe l'application existante en microservices derrière une passerelle et sécurise l'accès (JWT, rôles et permissions). Le second transforme la boutique unique en plateforme multi-boutiques : inscription des commerçants, assistant de création de boutique, vérification d'identité (KYC) et console d'administration de la plateforme.

%==================================================================
\section{Backlog du Sprint 1}
\Needspace{18\baselineskip}
Le tableau~\ref{tab:backlog-sprint1} présente le \textbf{backlog du Sprint 1}, consacré à l'architecture microservices, à l'authentification et au contrôle d'accès.

\renewcommand{\arraystretch}{1.05}
\setlength{\tabcolsep}{3pt}
\begin{longtable}{|>{\raggedright\arraybackslash}p{4cm}|>{\raggedright\arraybackslash}p{9cm}|>{\centering\arraybackslash}p{1.2cm}|}
\hline
\textbf{User Story} & \textbf{Tâches} & \textbf{Estimation} \\
\hline
\endfirsthead
\hline
\textbf{User Story} & \textbf{Tâches} & \textbf{Estimation} \\
\hline
\endhead
\hline
\multicolumn{3}{r}{\textit{Suite à la page suivante}} \\
\endfoot
\hline
\endlastfoot

\textbf{En tant qu'utilisateur, je souhaite créer un compte}
& Découper l'application NaturEssence en microservices (auth, catalog, order, marketing, analytics) et une bibliothèque partagée. & 1 \\
\cline{2-3}
& Mettre en place la passerelle Spring Cloud Gateway et ses règles de routage. & 1 \\
\cline{2-3}
& Développer l'API d'inscription avec hachage BCrypt et contrôle d'unicité de l'e-mail. & 1 \\
\cline{2-3}
& Développer l'interface d'inscription. & 1 \\
\cline{2-3}
& Tester la fonctionnalité. & 1 \\
\hline

\textbf{En tant qu'utilisateur, je souhaite me connecter}
& Développer l'API de connexion et la génération du jeton JWT signé. & 1 \\
\cline{2-3}
& Développer le filtre JWT partagé par tous les microservices. & 1 \\
\cline{2-3}
& Développer l'interface de connexion et la gestion des erreurs. & 1 \\
\cline{2-3}
& Tester la fonctionnalité. & 1 \\
\hline

\textbf{En tant qu'utilisateur, je souhaite rester connecté}
& Développer le jeton de rafraîchissement persistant et l'API \texttt{/refresh}. & 1 \\
\cline{2-3}
& Renouveler automatiquement le jeton côté client (intercepteur HTTP). & 1 \\
\hline

\textbf{En tant qu'utilisateur, je souhaite me déconnecter}
& Développer l'API de déconnexion (révocation du jeton de rafraîchissement) et le bouton associé. & 1 \\
\hline

\textbf{En tant que client, je souhaite consulter et modifier mon profil}
& Développer l'API et l'interface « Mon profil ». & 1 \\
\cline{2-3}
& Tester la fonctionnalité. & 1 \\
\hline

\textbf{En tant que commerçant, je souhaite créer des rôles et leur associer des permissions}
& Développer l'API CRUD des rôles et l'association des permissions. & 1 \\
\cline{2-3}
& Développer l'interface « Rôles et permissions » (matrice des permissions). & 1 \\
\cline{2-3}
& Tester la fonctionnalité. & 1 \\
\hline

\textbf{En tant que commerçant, je souhaite gérer les comptes de mon équipe et de mes clients}
& Développer l'API de gestion des utilisateurs (création, recherche, filtres, statut, statistiques). & 1 \\
\cline{2-3}
& Développer les interfaces « Clients », « Détail client » et « Ajouter un compte ». & 1 \\
\cline{2-3}
& Tester la fonctionnalité. & 1 \\
\hline

\textbf{En tant que système, je veux restreindre chaque endpoint selon les permissions}
& Configurer Spring Security dans chaque microservice (endpoints publics, profil, administration). & 1 \\
\cline{2-3}
& Protéger les routes du back-office côté interface (garde d'authentification). & 1 \\
\hline

\multicolumn{2}{|r|}{\textbf{Total}} & \textbf{22} \\
\hline
\caption{Backlog du Sprint 1}
\label{tab:backlog-sprint1}
\end{longtable}

\section{Diagramme de cas d'utilisation du Sprint 1}
\Needspace{20\baselineskip}
La figure~\ref{fig:cu-sprint1} présente le diagramme de cas d'utilisation du premier sprint.
\begin{figure}[H]
    \centering
    \img{0.8}{cu_sprint1.png}
    \caption{Diagramme de cas d'utilisation du Sprint 1}
    \label{fig:cu-sprint1}
\end{figure}

\section{Services web du Sprint 1}
Les services web permettent aux applications de communiquer via des API REST. Toutes les requêtes passent par la passerelle (\texttt{http://localhost:8080}), qui les route vers le service compétent. Le tableau~\ref{tab:ws-sprint1} présente les services web du Sprint 1.

\renewcommand{\arraystretch}{1.15}
\setlength{\tabcolsep}{4pt}
\begin{longtable}{|p{1.5cm}|p{6.8cm}|p{6.2cm}|}
\hline
\textbf{Méthode} & \textbf{URL} & \textbf{Description} \\
\hline
\endfirsthead
\hline
\textbf{Méthode} & \textbf{URL} & \textbf{Description} \\
\hline
\endhead
\multicolumn{3}{|c|}{\textbf{Authentification --- /api/v1/auth (auth-service)}} \\
\hline
POST & /api/v1/auth/register & Créer un compte client et retourner les jetons. \\
\hline
POST & /api/v1/auth/login & Authentifier un utilisateur et retourner le jeton d'accès et le jeton de rafraîchissement. \\
\hline
POST & /api/v1/auth/refresh & Obtenir un nouveau jeton d'accès. \\
\hline
POST & /api/v1/auth/logout & Révoquer le jeton de rafraîchissement. \\
\hline
GET & /api/v1/auth/me & Retourner l'utilisateur connecté, son rôle et sa boutique. \\
\hline
\multicolumn{3}{|c|}{\textbf{Profil --- /api/v1/profile}} \\
\hline
GET & /api/v1/profile & Consulter son profil. \\
\hline
PUT & /api/v1/profile & Modifier son profil. \\
\hline
\multicolumn{3}{|c|}{\textbf{Rôles et permissions --- /api/v1/admin/roles}} \\
\hline
GET & /api/v1/admin/roles & Lister les rôles. \\
\hline
POST & /api/v1/admin/roles & Créer un rôle. \\
\hline
PUT & /api/v1/admin/roles/\{id\} & Modifier un rôle. \\
\hline
DELETE & /api/v1/admin/roles/\{id\} & Supprimer un rôle. \\
\hline
PUT & /api/v1/admin/roles/\{id\}/permissions & Associer des permissions à un rôle. \\
\hline
GET & /api/v1/admin/roles/permissions-matrix & Obtenir la matrice rôles / permissions. \\
\hline
\multicolumn{3}{|c|}{\textbf{Utilisateurs --- /api/v1/admin/users}} \\
\hline
GET & /api/v1/admin/users & Lister les utilisateurs de la boutique. \\
\hline
POST & /api/v1/admin/users & Créer un compte. \\
\hline
PUT & /api/v1/admin/users/\{id\} & Modifier un compte. \\
\hline
PATCH & /api/v1/admin/users/\{id\}/status & Activer ou désactiver un compte. \\
\hline
GET & /api/v1/admin/users/search & Rechercher un utilisateur. \\
\hline
GET & /api/v1/admin/users/stats & Statistiques des comptes. \\
\hline
\caption{Services web du Sprint 1}
\label{tab:ws-sprint1}
\end{longtable}

\section{Les cas d'utilisation du Sprint 1}
\subsection{Cas d'utilisation « S'authentifier »}
\subsubsection{Description textuelle}
Le tableau~\ref{tab:uc-login} présente la description textuelle du cas d'utilisation « S'authentifier ».
\renewcommand{\arraystretch}{1.5}
\begin{longtable}{|p{4cm}|p{11cm}|}
\hline
\textbf{Acteurs} & Commerçant, Collaborateur, Client, Administrateur Sellio \\
\hline
\textbf{Objectif} & Accéder à la plateforme de façon sécurisée et obtenir un jeton reflétant son rôle, ses permissions et sa boutique. \\
\hline
\textbf{Pré-condition} & L'utilisateur possède un compte. \\
\hline
\textbf{Scénario principal} &
1. L'utilisateur saisit son e-mail et son mot de passe.\newline
2. L'interface envoie la requête à la passerelle, qui la route vers le service d'identité.\newline
3. Le système recherche le compte et compare le mot de passe à son empreinte BCrypt.\newline
4. Le système vérifie que le compte n'est pas bloqué.\newline
5. Le système génère un jeton d'accès JWT (24 heures) contenant l'identité, le rôle et la boutique, ainsi qu'un jeton de rafraîchissement (7 jours).\newline
6. Le système retourne les jetons et le profil de l'utilisateur.\newline
7. L'interface stocke les jetons et redirige l'utilisateur selon son rôle : console Sellio pour l'administrateur, back-office pour le commerçant, vitrine pour le client. \\
\hline
\textbf{Scénario alternatif} &
\textbf{Identifiants incorrects :} à l'étape 3, le système renvoie le message générique « Email ou mot de passe incorrect », sans préciser le champ erroné, pour ne pas faciliter l'énumération des comptes.\newline
\textbf{Compte bloqué :} à l'étape 4, le système répond 401 avec « Votre compte a été bloqué. Contactez l'équipe Sellio. »\newline
\textbf{Jeton expiré :} lors d'un appel ultérieur, l'intercepteur utilise le jeton de rafraîchissement pour obtenir un nouveau jeton sans redemander le mot de passe. \\
\hline
\caption{Description textuelle du cas d'utilisation « S'authentifier »}
\label{tab:uc-login}
\end{longtable}

\subsubsection{Diagramme de séquence système}
\Needspace{18\baselineskip}
La figure~\ref{fig:dss-login} présente le diagramme de séquence système de ce cas d'utilisation.
\begin{figure}[H]
    \centering
    \img{0.6}{dss_login.png}
    \caption{Diagramme de séquence système du cas « S'authentifier »}
    \label{fig:dss-login}
\end{figure}

\subsubsection{Diagramme de séquence objet}
\Needspace{20\baselineskip}
La figure~\ref{fig:dso-login} présente le diagramme de séquence objet : la requête traverse la passerelle, le contrôleur \texttt{AuthController}, le service \texttt{AuthService}, le gestionnaire d'authentification Spring Security, puis le dépôt \texttt{UserRepository}.
\begin{figure}[H]
    \centering
    \img{0.9}{dso_login.png}
    \caption{Diagramme de séquence objet du cas « S'authentifier »}
    \label{fig:dso-login}
\end{figure}

\section{Réalisation du Sprint 1}
\Needspace{20\baselineskip}
\textbf{Interface de connexion :} la figure~\ref{fig:ui-login} présente la page de connexion de Sellio, affichée par défaut en thème sombre avec une bascule vers le thème clair.
\begin{figure}[H]
    \centering
    \img{0.8}{ui_login.png}
    \caption{Interface de connexion à Sellio}
    \label{fig:ui-login}
\end{figure}

\Needspace{20\baselineskip}
\textbf{Interface de gestion des rôles et permissions :} la figure~\ref{fig:ui-roles} présente la matrice des permissions : le commerçant crée un rôle (par exemple « Préparateur de commandes ») et coche les modules auxquels il donne accès.
\begin{figure}[H]
    \centering
    \img{0.8}{ui_roles.png}
    \caption{Gestion des rôles et permissions}
    \label{fig:ui-roles}
\end{figure}

\Needspace{20\baselineskip}
\textbf{Interface de gestion des clients :} la figure~\ref{fig:ui-clients} présente la liste des clients de la boutique, avec recherche, filtres par statut et par segment, et accès à la fiche détaillée.
\begin{figure}[H]
    \centering
    \img{0.8}{ui_clients.png}
    \caption{Liste des clients de la boutique}
    \label{fig:ui-clients}
\end{figure}

%==================================================================
\section{Backlog du Sprint 2}
\Needspace{18\baselineskip}
Le tableau~\ref{tab:backlog-sprint2} présente le \textbf{backlog du Sprint 2}, consacré à la transformation en plateforme multi-boutiques.

\renewcommand{\arraystretch}{1.05}
\setlength{\tabcolsep}{3pt}
\begin{longtable}{|>{\raggedright\arraybackslash}p{4cm}|>{\raggedright\arraybackslash}p{9cm}|>{\centering\arraybackslash}p{1.2cm}|}
\hline
\textbf{User Story} & \textbf{Tâches} & \textbf{Estimation} \\
\hline
\endfirsthead
\hline
\textbf{User Story} & \textbf{Tâches} & \textbf{Estimation} \\
\hline
\endhead
\hline
\multicolumn{3}{r}{\textit{Suite à la page suivante}} \\
\endfoot
\hline
\endlastfoot

\textbf{En tant que visiteur, je souhaite découvrir Sellio sur une page d'accueil}
& Concevoir et développer la page d'accueil Sellio (thème sombre premium). & 1 \\
\cline{2-3}
& Développer la bascule thème sombre / thème clair. & 1 \\
\hline

\textbf{En tant que commerçant, je souhaite m'inscrire sur Sellio}
& Développer l'API d'inscription commerçant (\texttt{/register-merchant}). & 1 \\
\cline{2-3}
& Développer les interfaces d'inscription et de connexion Sellio. & 1 \\
\hline

\textbf{En tant que commerçant, je souhaite créer ma boutique avec un assistant}
& Ajouter l'entité \texttt{Shop} et rattacher utilisateurs, produits, commandes et événements à une boutique. & 1 \\
\cline{2-3}
& Filtrer toutes les requêtes par la boutique de l'utilisateur connecté (isolation multi-boutiques). & 1 \\
\cline{2-3}
& Développer l'API de création de boutique (nom, \textit{slug}, secteur, logo, couleurs). & 1 \\
\cline{2-3}
& Développer l'assistant en cinq étapes avec retour en arrière. & 1 \\
\cline{2-3}
& Tester la fonctionnalité. & 1 \\
\hline

\textbf{En tant que commerçant, je souhaite ouvrir ma vitrine depuis le back-office}
& Développer le lien « Voir le site » vers la vitrine de la boutique. & 1 \\
\hline

\textbf{En tant que commerçant, je souhaite enregistrer ma carte bancaire}
& Développer l'API Stripe de création d'un \textit{SetupIntent}. & 1 \\
\cline{2-3}
& Intégrer Stripe Elements et confirmer la carte côté navigateur. & 1 \\
\cline{2-3}
& Stocker uniquement la marque, les 4 derniers chiffres, l'expiration et la référence \texttt{pm\_}. & 1 \\
\hline

\textbf{En tant que commerçant, je souhaite soumettre ma pièce d'identité et un selfie}
& Développer l'entité \texttt{MerchantVerification} et l'API de soumission avec contrôles (CIN à 8 chiffres, passeport, taille des images). & 1 \\
\cline{2-3}
& Développer l'étape « Identité » (CIN recto-verso ou passeport, selfie). & 1 \\
\cline{2-3}
& Développer la page de suivi du dossier et la nouvelle soumission après refus. & 1 \\
\hline

\textbf{En tant qu'administrateur Sellio, je souhaite valider ou refuser un dossier}
& Développer les API d'approbation et de refus (motif obligatoire). & 1 \\
\cline{2-3}
& Développer l'écran de revue du dossier dans la console. & 1 \\
\cline{2-3}
& Tester la fonctionnalité. & 1 \\
\hline

\textbf{En tant qu'administrateur Sellio, je souhaite consulter les boutiques et leurs propriétaires}
& Créer le compte administrateur Sellio au démarrage (\textit{seed}). & 1 \\
\cline{2-3}
& Développer les API de statistiques, de liste et de détail des boutiques. & 1 \\
\cline{2-3}
& Développer la console avec pagination. & 1 \\
\hline

\textbf{En tant qu'administrateur Sellio, je souhaite suspendre une boutique ou bloquer un commerçant}
& Développer les API de changement de statut d'une boutique et d'un utilisateur. & 1 \\
\cline{2-3}
& Refuser la connexion d'un compte bloqué et afficher un message de boutique fermée sur la vitrine. & 1 \\
\cline{2-3}
& Tester la fonctionnalité. & 1 \\
\hline

\multicolumn{2}{|r|}{\textbf{Total}} & \textbf{25} \\
\hline
\caption{Backlog du Sprint 2}
\label{tab:backlog-sprint2}
\end{longtable}

\section{Diagramme de cas d'utilisation du Sprint 2}
\Needspace{20\baselineskip}
La figure~\ref{fig:cu-sprint2} présente le diagramme de cas d'utilisation du deuxième sprint.
\begin{figure}[H]
    \centering
    \img{0.85}{cu_sprint2.png}
    \caption{Diagramme de cas d'utilisation du Sprint 2}
    \label{fig:cu-sprint2}
\end{figure}

\section{Services web du Sprint 2}
Le tableau~\ref{tab:ws-sprint2} présente les services web du Sprint 2.
\renewcommand{\arraystretch}{1.15}
\setlength{\tabcolsep}{4pt}
\begin{longtable}{|p{1.5cm}|p{6.8cm}|p{6.2cm}|}
\hline
\textbf{Méthode} & \textbf{URL} & \textbf{Description} \\
\hline
\endfirsthead
\hline
\textbf{Méthode} & \textbf{URL} & \textbf{Description} \\
\hline
\endhead
\multicolumn{3}{|c|}{\textbf{Commerçant et boutique --- /api/v1/auth}} \\
\hline
POST & /api/v1/auth/register-merchant & Créer un compte commerçant. \\
\hline
POST & /api/v1/auth/my-shop & Créer sa boutique avec son dossier de vérification. \\
\hline
GET & /api/v1/auth/my-shop & Consulter sa boutique. \\
\hline
PATCH & /api/v1/auth/my-shop & Modifier les informations de sa boutique. \\
\hline
GET & /api/v1/auth/my-shop/verification & Consulter l'état de son dossier. \\
\hline
PUT & /api/v1/auth/my-shop/verification & Soumettre à nouveau un dossier refusé. \\
\hline
GET & /api/v1/public/shops/\{slug\} & Informations publiques d'une boutique (statut, gabarit, thème). \\
\hline
POST & /api/v1/public/stripe/setup-intent & Créer un \textit{SetupIntent} Stripe pour enregistrer la carte. \\
\hline
\multicolumn{3}{|c|}{\textbf{Console plateforme --- /api/v1/admin/platform}} \\
\hline
GET & /api/v1/admin/platform/stats & Statistiques globales (boutiques, commerçants, dossiers en attente). \\
\hline
GET & /api/v1/admin/platform/shops & Lister les boutiques et leurs propriétaires (paginé). \\
\hline
GET & /api/v1/admin/platform/shops/\{id\} & Détail d'une boutique. \\
\hline
PATCH & /api/v1/admin/platform/shops/\{id\}/status & Suspendre ou réactiver une boutique. \\
\hline
PATCH & /api/v1/admin/platform/users/\{id\}/status & Bloquer ou débloquer un commerçant. \\
\hline
GET & /api/v1/admin/platform/users & Lister les utilisateurs de la plateforme. \\
\hline
GET & /api/v1/admin/platform/verifications & Lister les dossiers de vérification. \\
\hline
GET & /api/v1/admin/platform/verifications/\{id\} & Consulter un dossier (pièces et selfie). \\
\hline
POST & /api/v1/admin/platform/verifications/\{id\}/approve & Valider un dossier : la boutique passe à \texttt{ACTIVE}. \\
\hline
POST & /api/v1/admin/platform/verifications/\{id\}/reject & Refuser un dossier avec un motif : la boutique passe à \texttt{REJECTED}. \\
\hline
\caption{Services web du Sprint 2}
\label{tab:ws-sprint2}
\end{longtable}

\section{Les cas d'utilisation du Sprint 2}
\subsection{Cas d'utilisation « Créer une boutique et faire vérifier son identité »}
\subsubsection{Description textuelle}
Le tableau~\ref{tab:uc-kyc} présente la description textuelle de ce cas d'utilisation.
\renewcommand{\arraystretch}{1.5}
\begin{longtable}{|p{4cm}|p{11cm}|}
\hline
\textbf{Acteurs} & Commerçant, Administrateur Sellio, Stripe \\
\hline
\textbf{Objectif} & Ouvrir une boutique qui ne pourra vendre qu'après vérification de l'identité et du moyen de paiement de son propriétaire. \\
\hline
\textbf{Pré-condition} & Le commerçant est inscrit et connecté ; il ne possède pas encore de boutique. \\
\hline
\textbf{Scénario principal} &
1. Le commerçant saisit le nom de sa boutique ; le système propose un identifiant (\textit{slug}).\newline
2. Il choisit le secteur (cosmétiques ou vêtements), puis son logo et sa palette.\newline
3. À l'étape « Carte », le système crée un \textit{SetupIntent} Stripe ; le commerçant saisit sa carte dans le champ sécurisé Stripe, qui la confirme sans que le numéro ne transite par Sellio.\newline
4. À l'étape « Identité », il choisit CIN ou passeport, photographie sa pièce et prend un selfie.\newline
5. Le système contrôle le dossier (format du numéro, validité de la carte, taille des images), crée la boutique au statut \texttt{PENDING} et le dossier au statut en attente.\newline
6. Le back-office affiche « Dossier en cours de vérification » ; la vitrine n'est pas encore ouverte au public.\newline
7. L'administrateur Sellio consulte le dossier, compare le selfie à la pièce d'identité et le valide.\newline
8. La boutique passe au statut \texttt{ACTIVE} ; le commerçant accède à tout le back-office. \\
\hline
\textbf{Scénario alternatif} &
\textbf{Carte refusée par Stripe :} à l'étape 3, le message de Stripe est affiché et le commerçant peut saisir une autre carte.\newline
\textbf{Dossier incomplet :} à l'étape 5, le système indique la pièce manquante (par exemple « Le numéro de CIN doit contenir 8 chiffres »).\newline
\textbf{Dossier refusé :} à l'étape 7, l'administrateur saisit obligatoirement un motif ; la boutique passe à \texttt{REJECTED}, le motif est affiché au commerçant qui peut soumettre un nouveau dossier. \\
\hline
\caption{Description textuelle du cas « Créer une boutique et faire vérifier son identité »}
\label{tab:uc-kyc}
\end{longtable}

\subsubsection{Diagramme de séquence système}
\Needspace{18\baselineskip}
La figure~\ref{fig:dss-kyc} présente le diagramme de séquence système de ce cas.
\begin{figure}[H]
    \centering
    \img{0.8}{dss_kyc.png}
    \caption{Diagramme de séquence système du cas « Créer une boutique et faire vérifier son identité »}
    \label{fig:dss-kyc}
\end{figure}

\subsubsection{Diagramme de séquence objet}
\Needspace{20\baselineskip}
La figure~\ref{fig:dso-kyc} présente le diagramme de séquence objet de ce cas.
\begin{figure}[H]
    \centering
    \img{0.95}{dso_kyc.png}
    \caption{Diagramme de séquence objet du cas « Créer une boutique et faire vérifier son identité »}
    \label{fig:dso-kyc}
\end{figure}

\section{Réalisation du Sprint 2}
\Needspace{20\baselineskip}
\textbf{Page d'accueil Sellio :} la figure~\ref{fig:ui-sellio-home} présente la page d'accueil de la plateforme, qui présente l'offre aux commerçants et donne accès à l'inscription et à la connexion.
\begin{figure}[H]
    \centering
    \img{0.85}{ui_sellio_home.png}
    \caption{Page d'accueil de Sellio}
    \label{fig:ui-sellio-home}
\end{figure}

\Needspace{20\baselineskip}
\textbf{Assistant de création de boutique :} les figures~\ref{fig:ui-wizard-1} et~\ref{fig:ui-wizard-2} présentent l'assistant : identité de la boutique et secteur, puis enregistrement de la carte et de la pièce d'identité.
\begin{figure}[H]
    \centering
    \img{0.8}{ui_nouvelle_boutique.png}
    \caption{Assistant de création de boutique : informations et secteur}
    \label{fig:ui-wizard-1}
\end{figure}
\begin{figure}[H]
    \centering
    \img{0.8}{ui_kyc_carte_identite.png}
    \caption{Assistant de création de boutique : carte bancaire et pièce d'identité}
    \label{fig:ui-wizard-2}
\end{figure}

\Needspace{20\baselineskip}
\textbf{Console d'administration Sellio :} la figure~\ref{fig:ui-console} présente la console : indicateurs globaux, liste paginée des boutiques avec leur propriétaire et leur statut, et actions de suspension ou de blocage. La figure~\ref{fig:ui-verif} présente la revue d'un dossier de vérification.
\begin{figure}[H]
    \centering
    \img{0.85}{ui_console_sellio.png}
    \caption{Console d'administration Sellio}
    \label{fig:ui-console}
\end{figure}
\begin{figure}[H]
    \centering
    \img{0.85}{ui_revue_dossier.png}
    \caption{Revue d'un dossier de vérification par l'administrateur}
    \label{fig:ui-verif}
\end{figure}

\subsection*{Conclusion}
Au terme de cette release, Sellio dispose d'une architecture microservices sécurisée et d'un véritable modèle multi-boutiques : un commerçant peut s'inscrire, créer sa boutique et la faire vérifier, et l'administrateur de la plateforme garde le contrôle sur les boutiques ouvertes. La release suivante porte sur le catalogue et la vitrine.

\thispagestyle{fancy}
\fancyhead[R]{A.U 2025-2026}
\fancyhead[L]{ chapitre 4: Release 2 : Catalogue et vitrine }
\chapter{Release 2 : Catalogue produit et vitrine personnalisable}
\section*{Introduction}
Ce chapitre présente la deuxième release, qui donne à chaque boutique son contenu et son apparence. Le Sprint 3 met en place un catalogue adapté aux deux secteurs visés, cosmétiques et vêtements, avec des attributs de variantes normalisés. Le Sprint 4 livre trois gabarits de vitrine complets et un éditeur de thème dont chaque réglage s'applique en direct à la vitrine.

%==================================================================
\section{Backlog du Sprint 3}
\Needspace{18\baselineskip}
Le tableau~\ref{tab:backlog-sprint3} présente le \textbf{backlog du Sprint 3}.

\renewcommand{\arraystretch}{1.05}
\setlength{\tabcolsep}{3pt}
\begin{longtable}{|>{\raggedright\arraybackslash}p{4cm}|>{\raggedright\arraybackslash}p{9cm}|>{\centering\arraybackslash}p{1.2cm}|}
\hline
\textbf{User Story} & \textbf{Tâches} & \textbf{Estimation} \\
\hline
\endfirsthead
\hline
\textbf{User Story} & \textbf{Tâches} & \textbf{Estimation} \\
\hline
\endhead
\hline
\multicolumn{3}{r}{\textit{Suite à la page suivante}} \\
\endfoot
\hline
\endlastfoot

\textbf{En tant que commerçant, je souhaite gérer mes produits}
& Développer les entités \texttt{Product} et \texttt{ProductVariant} (prix, stock, statut, images). & 1 \\
\cline{2-3}
& Développer l'API CRUD des produits, avec archivage, désactivation et statistiques. & 1 \\
\cline{2-3}
& Développer l'API de téléversement d'images. & 1 \\
\cline{2-3}
& Développer les interfaces « Produits », « Ajouter un produit » et « Modifier un produit ». & 1 \\
\cline{2-3}
& Adapter le formulaire au secteur de la boutique (champs cosmétiques ou vêtements). & 1 \\
\cline{2-3}
& Tester la fonctionnalité. & 1 \\
\hline

\textbf{En tant que commerçant, je souhaite gérer mes catégories}
& Développer l'API des catégories hiérarchiques (parent, enfants, ordre d'affichage). & 1 \\
\cline{2-3}
& Développer les interfaces de gestion et de réordonnancement des catégories. & 1 \\
\cline{2-3}
& Tester la fonctionnalité. & 1 \\
\hline

\textbf{En tant que commerçant, je souhaite gérer mes collections}
& Développer l'API des collections et l'association des produits. & 1 \\
\cline{2-3}
& Développer les interfaces « Collections » et « Gérer une collection ». & 1 \\
\hline

\textbf{En tant que commerçant, je souhaite choisir les tailles dans une liste}
& Définir les groupes de tailles (lettres, numériques, jeans, enfants, chaussures, taille unique). & 1 \\
\cline{2-3}
& Développer le sélecteur de tailles à choix multiple. & 1 \\
\hline

\textbf{En tant que commerçant, je souhaite choisir des couleurs prédéfinies ou libres}
& Définir la palette de 28 couleurs de base nommées. & 1 \\
\cline{2-3}
& Développer le sélecteur de couleurs avec ajout d'une couleur libre (sélecteur RVB et nom). & 1 \\
\hline

\textbf{En tant que commerçant, je souhaite choisir le tissu et d'autres attributs, et en ajouter}
& Définir la liste des tissus (fibres naturelles, artificielles, synthétiques) et des autres attributs (coupe, saison, entretien). & 1 \\
\cline{2-3}
& Stocker les valeurs ajoutées par la boutique (\texttt{customOptions}) et les fusionner avec les listes par défaut. & 1 \\
\cline{2-3}
& Développer les sélecteurs extensibles et le hook \texttt{useShopOptions}. & 1 \\
\cline{2-3}
& Tester la fonctionnalité. & 1 \\
\hline

\multicolumn{2}{|r|}{\textbf{Total}} & \textbf{19} \\
\hline
\caption{Backlog du Sprint 3}
\label{tab:backlog-sprint3}
\end{longtable}

\section{Diagramme de cas d'utilisation du Sprint 3}
\Needspace{20\baselineskip}
La figure~\ref{fig:cu-sprint3} présente le diagramme de cas d'utilisation du troisième sprint.
\begin{figure}[H]
    \centering
    \img{0.8}{cu_sprint3.png}
    \caption{Diagramme de cas d'utilisation du Sprint 3}
    \label{fig:cu-sprint3}
\end{figure}

\section{Services web du Sprint 3}
Le tableau~\ref{tab:ws-sprint3} présente les services web du Sprint 3, exposés par le microservice \texttt{catalog-service}.
\renewcommand{\arraystretch}{1.15}
\setlength{\tabcolsep}{4pt}
\begin{longtable}{|p{1.5cm}|p{6.8cm}|p{6.2cm}|}
\hline
\textbf{Méthode} & \textbf{URL} & \textbf{Description} \\
\hline
\endfirsthead
\hline
\textbf{Méthode} & \textbf{URL} & \textbf{Description} \\
\hline
\endhead
\multicolumn{3}{|c|}{\textbf{Produits (administration) --- /api/v1/admin/products}} \\
\hline
GET & /api/v1/admin/products & Lister les produits de la boutique. \\
\hline
GET & /api/v1/admin/products/stats & Statistiques du catalogue. \\
\hline
POST & /api/v1/admin/products & Créer un produit et ses variantes. \\
\hline
PUT & /api/v1/admin/products/\{id\} & Modifier un produit. \\
\hline
PATCH & /api/v1/admin/products/\{id\}/archive & Archiver un produit. \\
\hline
PATCH & /api/v1/admin/products/\{id\}/deactivate & Désactiver un produit. \\
\hline
DELETE & /api/v1/admin/products/\{id\} & Supprimer un produit. \\
\hline
POST & /api/v1/admin/upload & Téléverser une image. \\
\hline
\multicolumn{3}{|c|}{\textbf{Catégories et collections (administration)}} \\
\hline
GET, POST & /api/v1/admin/categories & Lister, créer une catégorie. \\
\hline
PUT, DELETE & /api/v1/admin/categories/\{id\} & Modifier, supprimer une catégorie. \\
\hline
PUT & /api/v1/admin/categories/reorder & Réordonner les catégories. \\
\hline
GET, POST & /api/v1/admin/collections & Lister, créer une collection. \\
\hline
PUT, DELETE & /api/v1/admin/collections/\{id\} & Modifier, supprimer une collection. \\
\hline
\multicolumn{3}{|c|}{\textbf{Catalogue public (vitrine)}} \\
\hline
GET & /api/v1/public/products & Lister et filtrer les produits publiés. \\
\hline
GET & /api/v1/public/products/\{slug\} & Fiche d'un produit. \\
\hline
GET & /api/v1/public/products/category/\{categoryId\} & Produits d'une catégorie. \\
\hline
GET & /api/v1/public/products/collection/\{collectionSlug\} & Produits d'une collection. \\
\hline
GET & /api/v1/public/products/by-ids & Produits à partir d'une liste d'identifiants (recommandations). \\
\hline
GET & /api/v1/public/categories/menu & Arborescence du menu de la vitrine. \\
\hline
GET & /api/v1/public/categories/footer & Catégories du pied de page. \\
\hline
GET & /api/v1/public/collections/homepage & Collections mises en avant. \\
\hline
\caption{Services web du Sprint 3}
\label{tab:ws-sprint3}
\end{longtable}

\section{Les cas d'utilisation du Sprint 3}
\subsection{Cas d'utilisation « Ajouter un produit avec ses variantes »}
\subsubsection{Description textuelle}
Le tableau~\ref{tab:uc-produit} présente la description textuelle de ce cas.
\renewcommand{\arraystretch}{1.5}
\begin{longtable}{|p{4cm}|p{11cm}|}
\hline
\textbf{Acteurs} & Commerçant, Collaborateur disposant de la permission « Produits » \\
\hline
\textbf{Objectif} & Publier un produit avec des attributs normalisés, exploitables par les filtres de la vitrine et par le moteur de recommandation. \\
\hline
\textbf{Pré-condition} & L'utilisateur est connecté et sa boutique est active. \\
\hline
\textbf{Scénario principal} &
1. L'utilisateur ouvre « Ajouter un produit » ; le formulaire s'adapte au secteur de la boutique.\newline
2. Il saisit le nom, la description, le prix, la catégorie et téléverse les images.\newline
3. Pour un vêtement, il choisit les tailles dans les groupes proposés, les couleurs dans la palette, et le tissu dans la liste.\newline
4. Le système génère une variante par combinaison taille / couleur, dont l'utilisateur ajuste le stock.\newline
5. L'utilisateur enregistre ; le système valide les données et crée le produit rattaché à sa boutique.\newline
6. Le produit apparaît dans la liste et, s'il est publié, sur la vitrine. \\
\hline
\textbf{Scénario alternatif} &
\textbf{Valeur absente de la liste :} à l'étape 3, l'utilisateur ajoute une couleur libre (sélecteur RVB et nom) ou un nouveau tissu ; la valeur est enregistrée dans les options de la boutique et proposée pour les produits suivants.\newline
\textbf{Données invalides :} à l'étape 5, le système signale les champs erronés (prix négatif, nom manquant) sans perdre la saisie. \\
\hline
\caption{Description textuelle du cas « Ajouter un produit avec ses variantes »}
\label{tab:uc-produit}
\end{longtable}

\subsubsection{Diagramme de séquence système}
\Needspace{18\baselineskip}
\begin{figure}[H]
    \centering
    \img{0.65}{dss_produit.png}
    \caption{Diagramme de séquence système du cas « Ajouter un produit »}
    \label{fig:dss-produit}
\end{figure}

\subsubsection{Diagramme de séquence objet}
\Needspace{20\baselineskip}
\begin{figure}[H]
    \centering
    \img{0.9}{dso_produit.png}
    \caption{Diagramme de séquence objet du cas « Ajouter un produit »}
    \label{fig:dso-produit}
\end{figure}

\section{Réalisation du Sprint 3}
\Needspace{20\baselineskip}
\textbf{Liste des produits :} la figure~\ref{fig:ui-produits} présente la liste des produits avec leurs statistiques, filtres et actions.
\begin{figure}[H]
    \centering
    \img{0.85}{ui_produits.png}
    \caption{Liste des produits de la boutique}
    \label{fig:ui-produits}
\end{figure}

\Needspace{20\baselineskip}
\textbf{Sélecteurs d'attributs :} la figure~\ref{fig:ui-attributs} présente le formulaire d'un vêtement : groupes de tailles, palette de couleurs avec ajout d'une couleur libre, et liste des tissus extensible. Remplacer la saisie libre par ces listes évite les doublons (« bleu marine », « Marine », « navy ») qui fausseraient les filtres et les recommandations.
\begin{figure}[H]
    \centering
    \img{0.85}{ui_attributs.png}
    \caption{Sélecteurs de tailles, couleurs et tissus}
    \label{fig:ui-attributs}
\end{figure}

\Needspace{20\baselineskip}
\textbf{Catégories et collections :} la figure~\ref{fig:ui-categories} présente la gestion des catégories hiérarchiques.
\begin{figure}[H]
    \centering
    \img{0.85}{ui_categories.png}
    \caption{Gestion des catégories}
    \label{fig:ui-categories}
\end{figure}

%==================================================================
\section{Backlog du Sprint 4}
\Needspace{18\baselineskip}
Le tableau~\ref{tab:backlog-sprint4} présente le \textbf{backlog du Sprint 4}, consacré à la vitrine et à sa personnalisation.

\renewcommand{\arraystretch}{1.05}
\setlength{\tabcolsep}{3pt}
\begin{longtable}{|>{\raggedright\arraybackslash}p{4cm}|>{\raggedright\arraybackslash}p{9cm}|>{\centering\arraybackslash}p{1.2cm}|}
\hline
\textbf{User Story} & \textbf{Tâches} & \textbf{Estimation} \\
\hline
\endfirsthead
\hline
\textbf{User Story} & \textbf{Tâches} & \textbf{Estimation} \\
\hline
\endhead
\hline
\multicolumn{3}{r}{\textit{Suite à la page suivante}} \\
\endfoot
\hline
\endlastfoot

\textbf{En tant que commerçant, je souhaite choisir un gabarit Minimal, Bold ou Luxury}
& Concevoir l'architecture des gabarits : en-tête, accueil et pied de page par gabarit, briques partagées (navigation, données d'accueil, catégories du pied de page). & 1 \\
\cline{2-3}
& Développer le gabarit Minimal (épuré, typographie fine, grands espaces). & 1 \\
\cline{2-3}
& Développer le gabarit Bold (contrastes forts, en-tête transparent sur bannière plein écran). & 1 \\
\cline{2-3}
& Développer le gabarit Luxury (typographie à empattements, tons sourds, mise en page éditoriale). & 1 \\
\cline{2-3}
& Faire correspondre les anciens thèmes aux trois gabarits et permettre la prévisualisation (\texttt{?preview=}). & 1 \\
\cline{2-3}
& Supprimer la section codée en dur propre à NaturEssence (atelier « DIY »). & 1 \\
\hline

\textbf{En tant que visiteur, je souhaite parcourir la vitrine, ses menus et ses catégories}
& Développer les menus déroulants alimentés par l'arborescence des catégories. & 1 \\
\cline{2-3}
& Développer la page de catégorie avec filtres (prix, taille, couleur). & 1 \\
\cline{2-3}
& Afficher un message lorsque la boutique est en attente, fermée ou suspendue. & 1 \\
\hline

\textbf{En tant que visiteur, je souhaite consulter la fiche d'un produit}
& Développer la fiche produit (galerie, variantes, stock, avis). & 1 \\
\cline{2-3}
& Adapter la carte produit à chaque gabarit, avec choix rapide de la taille. & 1 \\
\hline

\textbf{En tant que commerçant, je souhaite choisir la palette de couleurs de ma boutique}
& Exposer la palette sous forme de variables CSS (\texttt{--rgb-primary}, \texttt{--rgb-accent}, \dots) consommées par Tailwind. & 1 \\
\cline{2-3}
& Développer le choix de palette dans « Votre boutique » (couleurs prédéfinies et personnalisées). & 1 \\
\hline

\textbf{En tant que commerçant, je souhaite personnaliser boutons, survols, navigation, textes et pied de page}
& Définir le modèle de thème (14 variables : annonce, navigation, survol, titres, boutons, prix, badges, pied de page). & 1 \\
\cline{2-3}
& Développer l'éditeur de thème avec aperçu et réinitialisation. & 1 \\
\cline{2-3}
& Appliquer le thème sur la vitrine via des variables \texttt{--t-*} et des classes d'ancrage (\texttt{t-nav}, \texttt{t-btn}, \texttt{t-price}\dots). & 1 \\
\cline{2-3}
& Valider le thème côté serveur (couleurs \#RRGGBB, taille et type du logo). & 1 \\
\cline{2-3}
& Tester sur les trois gabarits. & 1 \\
\hline

\multicolumn{2}{|r|}{\textbf{Total}} & \textbf{18} \\
\hline
\caption{Backlog du Sprint 4}
\label{tab:backlog-sprint4}
\end{longtable}

\section{Diagramme de cas d'utilisation du Sprint 4}
\Needspace{20\baselineskip}
\begin{figure}[H]
    \centering
    \img{0.8}{cu_sprint4.png}
    \caption{Diagramme de cas d'utilisation du Sprint 4}
    \label{fig:cu-sprint4}
\end{figure}

\section{Services web du Sprint 4}
Le tableau~\ref{tab:ws-sprint4} présente les services web du Sprint 4.
\renewcommand{\arraystretch}{1.15}
\setlength{\tabcolsep}{4pt}
\begin{longtable}{|p{1.5cm}|p{6.8cm}|p{6.2cm}|}
\hline
\textbf{Méthode} & \textbf{URL} & \textbf{Description} \\
\hline
\endfirsthead
\hline
\textbf{Méthode} & \textbf{URL} & \textbf{Description} \\
\hline
\endhead
\multicolumn{3}{|c|}{\textbf{Boutique et thème (auth-service)}} \\
\hline
PATCH & /api/v1/auth/my-shop & Enregistrer le gabarit, la palette, le thème et le logo. \\
\hline
GET & /api/v1/public/shops/\{slug\} & Charger le gabarit, la palette et le thème d'une boutique. \\
\hline
\multicolumn{3}{|c|}{\textbf{Apparence et profil de la boutique (marketing-service)}} \\
\hline
GET & /api/v1/admin/appearance/\{scope\} & Consulter les réglages d'apparence d'une zone (accueil, en-tête, pied de page). \\
\hline
PUT & /api/v1/admin/appearance/\{scope\} & Modifier les réglages d'apparence. \\
\hline
POST & /api/v1/admin/appearance/\{scope\}/reset & Réinitialiser les réglages. \\
\hline
GET & /api/v1/public/appearance/\{scope\} & Réglages d'apparence publics. \\
\hline
GET, PUT & /api/v1/admin/store & Consulter, modifier le profil de la boutique (contact, réseaux sociaux). \\
\hline
GET & /api/v1/public/store & Profil public de la boutique. \\
\hline
\caption{Services web du Sprint 4}
\label{tab:ws-sprint4}
\end{longtable}

\section{Les cas d'utilisation du Sprint 4}
\subsection{Cas d'utilisation « Personnaliser la vitrine »}
\subsubsection{Description textuelle}
Le tableau~\ref{tab:uc-theme} présente la description textuelle de ce cas.
\renewcommand{\arraystretch}{1.5}
\begin{longtable}{|p{4cm}|p{11cm}|}
\hline
\textbf{Acteurs} & Commerçant \\
\hline
\textbf{Objectif} & Donner à la vitrine l'apparence de la marque sans écrire de code. \\
\hline
\textbf{Pré-condition} & Le commerçant est connecté et possède une boutique. \\
\hline
\textbf{Scénario principal} &
1. Le commerçant ouvre « Votre boutique ».\newline
2. Il choisit un gabarit (Minimal, Bold ou Luxury) et peut le prévisualiser sans l'enregistrer.\newline
3. Il choisit sa palette : couleur de marque, accent, fond, texte.\newline
4. Dans l'éditeur de thème, il règle le bandeau d'annonce, la navigation et son survol, les boutons et leur survol, les prix, les badges et le pied de page.\newline
5. Il enregistre ; le système valide chaque couleur (format \#RRGGBB) et stocke le thème.\newline
6. À l'ouverture suivante, la vitrine charge le thème et l'applique sous forme de variables CSS ; seuls les réglages renseignés remplacent le style par défaut du gabarit. \\
\hline
\textbf{Scénario alternatif} &
\textbf{Couleur invalide :} à l'étape 5, le système refuse le thème avec « Le thème ne peut contenir que des couleurs \#RRGGBB ».\newline
\textbf{Retour au style du gabarit :} le commerçant réinitialise un réglage ; la variable correspondante est retirée et le style d'origine du gabarit s'applique de nouveau. \\
\hline
\caption{Description textuelle du cas « Personnaliser la vitrine »}
\label{tab:uc-theme}
\end{longtable}

\subsubsection{Diagramme de séquence système}
\Needspace{18\baselineskip}
\begin{figure}[H]
    \centering
    \img{0.65}{dss_theme.png}
    \caption{Diagramme de séquence système du cas « Personnaliser la vitrine »}
    \label{fig:dss-theme}
\end{figure}

\subsubsection{Diagramme de séquence objet}
\Needspace{20\baselineskip}
\begin{figure}[H]
    \centering
    \img{0.9}{dso_theme.png}
    \caption{Diagramme de séquence objet du cas « Personnaliser la vitrine »}
    \label{fig:dso-theme}
\end{figure}

\section{Réalisation du Sprint 4}
\Needspace{20\baselineskip}
\textbf{Les trois gabarits :} les figures~\ref{fig:ui-minimal}, \ref{fig:ui-bold} et~\ref{fig:ui-luxury} présentent la même boutique affichée avec chacun des trois gabarits. Le contenu (produits, catégories, bannières) est identique ; seuls la structure et le style changent.
\begin{figure}[H]
    \centering
    \img{0.85}{ui_vitrine_minimal.png}
    \caption{Vitrine avec le gabarit Minimal}
    \label{fig:ui-minimal}
\end{figure}
\begin{figure}[H]
    \centering
    \img{0.85}{ui_vitrine_bold.png}
    \caption{Vitrine avec le gabarit Bold}
    \label{fig:ui-bold}
\end{figure}
\begin{figure}[H]
    \centering
    \img{0.85}{ui_vitrine_luxury.png}
    \caption{Vitrine avec le gabarit Luxury}
    \label{fig:ui-luxury}
\end{figure}

\Needspace{20\baselineskip}
\textbf{Éditeur de thème :} la figure~\ref{fig:ui-theme} présente l'éditeur de thème, organisé en groupes (couleurs de base, bandeau d'annonce, navigation, boutons, prix et badges, pied de page).
\begin{figure}[H]
    \centering
    \img{0.85}{ui_editeur_theme.png}
    \caption{Éditeur de thème de la vitrine}
    \label{fig:ui-theme}
\end{figure}

\Needspace{20\baselineskip}
\textbf{Fiche produit :} la figure~\ref{fig:ui-fiche} présente la fiche d'un vêtement, avec le choix de la taille et de la couleur.
\begin{figure}[H]
    \centering
    \img{0.85}{ui_fiche_produit.png}
    \caption{Fiche produit sur la vitrine}
    \label{fig:ui-fiche}
\end{figure}

\subsection*{Conclusion}
Cette release a doté chaque boutique d'un catalogue structuré et d'une vitrine qui lui ressemble. Les attributs normalisés ne servent pas qu'à l'affichage : ils alimentent aussi la partie « contenu » du moteur de recommandation présenté au chapitre 8. La release suivante porte sur le parcours d'achat.

\thispagestyle{fancy}
\fancyhead[R]{A.U 2025-2026}
\fancyhead[L]{ chapitre 5: Release 3 : Parcours transactionnel }
\chapter{Release 3 : Parcours transactionnel}
\section*{Introduction}
Ce chapitre présente la troisième release, cœur transactionnel de la plateforme. Le Sprint 5 couvre l'acte d'achat : panier dynamique, tunnel de commande et paiement par carte avec Stripe. Le Sprint 6 couvre ce qui suit la commande : gestion des statuts par le commerçant, tableau de bord des ventes, retours, paramétrage de la TVA et des zones de livraison, et avis clients.

%==================================================================
\section{Backlog du Sprint 5}
\Needspace{18\baselineskip}
Le tableau~\ref{tab:backlog-sprint5} présente le \textbf{backlog du Sprint 5}.

\renewcommand{\arraystretch}{1.05}
\setlength{\tabcolsep}{3pt}
\begin{longtable}{|>{\raggedright\arraybackslash}p{4cm}|>{\raggedright\arraybackslash}p{9cm}|>{\centering\arraybackslash}p{1.2cm}|}
\hline
\textbf{User Story} & \textbf{Tâches} & \textbf{Estimation} \\
\hline
\endfirsthead
\hline
\textbf{User Story} & \textbf{Tâches} & \textbf{Estimation} \\
\hline
\endhead
\hline
\multicolumn{3}{r}{\textit{Suite à la page suivante}} \\
\endfoot
\hline
\endlastfoot

\textbf{En tant que visiteur, je souhaite gérer mon panier en choisissant la taille}
& Développer le contexte panier (ajout, quantité, suppression, total) persistant dans le navigateur. & 1 \\
\cline{2-3}
& Distinguer les lignes par produit, taille et couleur (clé de ligne). & 1 \\
\cline{2-3}
& Bloquer « Ajouter au panier » tant que la taille n'est pas choisie, sur la fiche et sur la carte produit. & 1 \\
\cline{2-3}
& Développer le tiroir de panier. & 1 \\
\cline{2-3}
& Tester la fonctionnalité. & 1 \\
\hline

\textbf{En tant que client, je souhaite passer commande avec livraison et code promo}
& Développer l'API de création de commande (sous-total, remise, livraison, TVA incluse, décrément du stock). & 1 \\
\cline{2-3}
& Calculer les frais de livraison selon la zone et la TVA selon la configuration. & 1 \\
\cline{2-3}
& Développer l'API de validation d'un code promo (dates, montant minimum, catégories, produits). & 1 \\
\cline{2-3}
& Développer le tunnel de commande (adresse, livraison, paiement, récapitulatif) et la page de confirmation. & 1 \\
\cline{2-3}
& Envoyer l'e-mail de confirmation de commande. & 1 \\
\cline{2-3}
& Tester la fonctionnalité. & 1 \\
\hline

\textbf{En tant que client, je souhaite retrouver mon panier sur un autre appareil}
& Développer l'API de sauvegarde du panier du client connecté et sa restauration à la connexion. & 1 \\
\hline

\textbf{En tant que client, je souhaite payer ma commande par carte}
& Configurer le SDK Stripe côté serveur (clé secrète en variable d'environnement). & 1 \\
\cline{2-3}
& Développer l'API de création d'un \textit{PaymentIntent}. & 1 \\
\cline{2-3}
& Intégrer Stripe Elements et confirmer le paiement dans le navigateur. & 1 \\
\cline{2-3}
& Conserver le paiement à la livraison comme alternative. & 1 \\
\cline{2-3}
& Tester avec les cartes de test Stripe (succès, refus, authentification 3-D Secure). & 1 \\
\hline

\multicolumn{2}{|r|}{\textbf{Total}} & \textbf{17} \\
\hline
\caption{Backlog du Sprint 5}
\label{tab:backlog-sprint5}
\end{longtable}

\section{Diagramme de cas d'utilisation du Sprint 5}
\Needspace{20\baselineskip}
\begin{figure}[H]
    \centering
    \img{0.8}{cu_sprint5.png}
    \caption{Diagramme de cas d'utilisation du Sprint 5}
    \label{fig:cu-sprint5}
\end{figure}

\section{Services web du Sprint 5}
Le tableau~\ref{tab:ws-sprint5} présente les services web du Sprint 5.
\renewcommand{\arraystretch}{1.15}
\setlength{\tabcolsep}{4pt}
\begin{longtable}{|p{1.5cm}|p{6.8cm}|p{6.2cm}|}
\hline
\textbf{Méthode} & \textbf{URL} & \textbf{Description} \\
\hline
\endfirsthead
\hline
\textbf{Méthode} & \textbf{URL} & \textbf{Description} \\
\hline
\endhead
\multicolumn{3}{|c|}{\textbf{Tunnel de commande --- /api/v1/public/checkout (order-service)}} \\
\hline
GET & /api/v1/public/checkout/shipping-zones & Zones de livraison et tarifs. \\
\hline
GET & /api/v1/public/checkout/tva-config & Configuration de la TVA. \\
\hline
POST & /api/v1/public/checkout/orders & Créer une commande. \\
\hline
GET & /api/v1/public/coupons/validate & Valider un code promo pour un panier. \\
\hline
GET & /api/v1/public/coupons/announcement & Promotion à afficher dans le bandeau d'annonce. \\
\hline
\multicolumn{3}{|c|}{\textbf{Paiement --- /api/v1/public/stripe}} \\
\hline
POST & /api/v1/public/stripe/payment-intent & Créer un \textit{PaymentIntent} et retourner son \textit{client secret}. \\
\hline
\multicolumn{3}{|c|}{\textbf{Panier du client --- /api/v1/profile (auth-service)}} \\
\hline
GET & /api/v1/profile/cart & Récupérer le panier sauvegardé. \\
\hline
PUT & /api/v1/profile/cart & Sauvegarder le panier. \\
\hline
\caption{Services web du Sprint 5}
\label{tab:ws-sprint5}
\end{longtable}

\section{Les cas d'utilisation du Sprint 5}
\subsection{Cas d'utilisation « Passer commande et payer par carte »}
\subsubsection{Description textuelle}
Le tableau~\ref{tab:uc-commande} présente la description textuelle de ce cas.
\renewcommand{\arraystretch}{1.5}
\begin{longtable}{|p{4cm}|p{11cm}|}
\hline
\textbf{Acteurs} & Client, Stripe \\
\hline
\textbf{Objectif} & Transformer un panier en commande payée. \\
\hline
\textbf{Pré-condition} & Le panier contient au moins un article ; la taille de chaque vêtement est choisie ; la boutique est active. \\
\hline
\textbf{Scénario principal} &
1. Le client ouvre le tunnel de commande et saisit ou choisit son adresse.\newline
2. Il choisit sa ville ; le système affiche les frais de la zone de livraison correspondante et la TVA.\newline
3. Il saisit un code promo ; le système le valide et affiche la remise.\newline
4. Il choisit « Carte bancaire » ; le système crée un \textit{PaymentIntent} Stripe et renvoie son \textit{client secret}.\newline
5. Le client saisit sa carte dans le champ sécurisé Stripe, qui confirme le paiement (avec authentification 3-D Secure si la banque l'exige).\newline
6. Après confirmation, l'interface envoie la commande ; le système calcule le sous-total, applique la remise, la livraison (gratuite au-delà du seuil) et la TVA incluse dans les prix, enregistre la commande au statut \texttt{EN\_ATTENTE} avec le mode de paiement \texttt{CARTE}, puis décrémente le stock.\newline
7. Le système envoie l'e-mail de confirmation ; le client est redirigé vers la page de confirmation et un événement d'achat est enregistré pour l'analyse comportementale. \\
\hline
\textbf{Scénario alternatif} &
\textbf{Code promo invalide :} à l'étape 3, le système en donne la raison (expiré, montant minimum non atteint, réservé à d'autres catégories).\newline
\textbf{Paiement refusé :} à l'étape 5, le message de Stripe est affiché ; aucune commande n'est créée et le client peut réessayer.\newline
\textbf{Paiement à la livraison :} à l'étape 4, si le client choisit ce mode, la commande est créée directement, sans passer par Stripe. \\
\hline
\caption{Description textuelle du cas « Passer commande et payer par carte »}
\label{tab:uc-commande}
\end{longtable}

\subsubsection{Diagramme de séquence système}
\Needspace{18\baselineskip}
\begin{figure}[H]
    \centering
    \img{0.7}{dss_commande.png}
    \caption{Diagramme de séquence système du cas « Passer commande et payer par carte »}
    \label{fig:dss-commande}
\end{figure}

\subsubsection{Diagramme de séquence objet}
\Needspace{20\baselineskip}
\begin{figure}[H]
    \centering
    \img{0.95}{dso_commande.png}
    \caption{Diagramme de séquence objet du cas « Passer commande et payer par carte »}
    \label{fig:dso-commande}
\end{figure}

\section{Réalisation du Sprint 5}
\Needspace{20\baselineskip}
\textbf{Panier :} la figure~\ref{fig:ui-panier} présente le tiroir de panier. Chaque ligne correspond à une combinaison produit / taille / couleur ; la section « Souvent achetés ensemble », ajoutée au Sprint 10, y propose des articles complémentaires.
\begin{figure}[H]
    \centering
    \img{0.8}{ui_panier.png}
    \caption{Tiroir de panier}
    \label{fig:ui-panier}
\end{figure}

\Needspace{20\baselineskip}
\textbf{Tunnel de commande et paiement :} les figures~\ref{fig:ui-checkout} et~\ref{fig:ui-stripe} présentent le tunnel de commande puis le formulaire de paiement Stripe.
\begin{figure}[H]
    \centering
    \img{0.85}{ui_checkout.png}
    \caption{Tunnel de commande : adresse, livraison et code promo}
    \label{fig:ui-checkout}
\end{figure}
\begin{figure}[H]
    \centering
    \img{0.8}{ui_paiement_stripe.png}
    \caption{Paiement par carte avec Stripe}
    \label{fig:ui-stripe}
\end{figure}

\Needspace{10\baselineskip}
\textbf{Remarque sur la devise :} Stripe ne prend pas en charge le dinar tunisien ; en mode test, les montants sont donc débités en euros. Une mise en production en Tunisie passerait par un prestataire de paiement local ou par la conversion des montants, et la confirmation de la commande devrait alors s'appuyer sur un \textit{webhook} du prestataire plutôt que sur le seul retour du navigateur. De même, les prix unitaires transmis par le panier devraient être revérifiés côté serveur à partir du catalogue, et le stock contrôlé avant la validation. Ces points figurent dans les perspectives.

%==================================================================
\section{Backlog du Sprint 6}
\Needspace{18\baselineskip}
Le tableau~\ref{tab:backlog-sprint6} présente le \textbf{backlog du Sprint 6}.

\renewcommand{\arraystretch}{1.05}
\setlength{\tabcolsep}{3pt}
\begin{longtable}{|>{\raggedright\arraybackslash}p{4cm}|>{\raggedright\arraybackslash}p{9cm}|>{\centering\arraybackslash}p{1.2cm}|}
\hline
\textbf{User Story} & \textbf{Tâches} & \textbf{Estimation} \\
\hline
\endfirsthead
\hline
\textbf{User Story} & \textbf{Tâches} & \textbf{Estimation} \\
\hline
\endhead
\hline
\multicolumn{3}{r}{\textit{Suite à la page suivante}} \\
\endfoot
\hline
\endlastfoot

\textbf{En tant que commerçant, je souhaite gérer mes commandes et leurs statuts}
& Développer l'API de liste, de détail et de recherche par référence des commandes. & 1 \\
\cline{2-3}
& Développer le changement de statut (en attente, confirmée, en préparation, expédiée, livrée, annulée, remboursée), avec remise en stock en cas d'annulation et e-mail au client à la livraison. & 1 \\
\cline{2-3}
& Développer les interfaces « Commandes » et « Détail commande ». & 1 \\
\cline{2-3}
& Tester la fonctionnalité. & 1 \\
\hline

\textbf{En tant que client, je souhaite suivre mes commandes}
& Développer l'API et la page « Mes commandes ». & 1 \\
\hline

\textbf{En tant que commerçant, je souhaite consulter un tableau de bord des ventes}
& Développer l'API d'agrégats (chiffre d'affaires, commandes, panier moyen, meilleurs produits). & 1 \\
\cline{2-3}
& Développer le tableau de bord avec graphiques. & 1 \\
\hline

\textbf{En tant que client, je souhaite demander un retour}
& Développer l'API de demande de retour (articles, motif) dans le délai de la politique. & 1 \\
\cline{2-3}
& Développer la page « Mes retours ». & 1 \\
\hline

\textbf{En tant que commerçant, je souhaite traiter les retours et définir ma politique de retour}
& Développer l'API de traitement des retours et de la politique de retour. & 1 \\
\cline{2-3}
& Développer les interfaces « Retours », « Détail retour » et « Historique des remboursements ». & 1 \\
\cline{2-3}
& Tester la fonctionnalité. & 1 \\
\hline

\textbf{En tant que commerçant, je souhaite paramétrer la TVA et les zones de livraison}
& Développer l'API de configuration de la TVA et des taux. & 1 \\
\cline{2-3}
& Développer l'API des zones de livraison (villes, tarif, seuil de gratuité, délai). & 1 \\
\cline{2-3}
& Développer l'interface « TVA et livraison ». & 1 \\
\hline

\textbf{En tant que client, je souhaite publier un avis sur un produit acheté}
& Développer l'API de publication (un avis par produit et par commande) et l'affichage sur la fiche produit. & 1 \\
\hline

\textbf{En tant que commerçant, je souhaite modérer les avis et y répondre}
& Développer l'API et l'interface de modération (statut, réponse, suppression). & 1 \\
\hline

\multicolumn{2}{|r|}{\textbf{Total}} & \textbf{17} \\
\hline
\caption{Backlog du Sprint 6}
\label{tab:backlog-sprint6}
\end{longtable}

\section{Diagramme de cas d'utilisation du Sprint 6}
\Needspace{20\baselineskip}
\begin{figure}[H]
    \centering
    \img{0.85}{cu_sprint6.png}
    \caption{Diagramme de cas d'utilisation du Sprint 6}
    \label{fig:cu-sprint6}
\end{figure}

\section{Services web du Sprint 6}
Le tableau~\ref{tab:ws-sprint6} présente les services web du Sprint 6.
\renewcommand{\arraystretch}{1.15}
\setlength{\tabcolsep}{4pt}
\begin{longtable}{|p{1.5cm}|p{6.8cm}|p{6.2cm}|}
\hline
\textbf{Méthode} & \textbf{URL} & \textbf{Description} \\
\hline
\endfirsthead
\hline
\textbf{Méthode} & \textbf{URL} & \textbf{Description} \\
\hline
\endhead
\multicolumn{3}{|c|}{\textbf{Commandes --- order-service}} \\
\hline
GET & /api/v1/admin/orders & Lister les commandes de la boutique. \\
\hline
GET & /api/v1/admin/orders/\{id\} & Détail d'une commande. \\
\hline
GET & /api/v1/admin/orders/reference/\{reference\} & Rechercher une commande par référence. \\
\hline
GET & /api/v1/admin/orders/user/\{userId\} & Commandes d'un client. \\
\hline
PATCH & /api/v1/admin/orders/\{id\}/status & Changer le statut d'une commande. \\
\hline
GET & /api/v1/profile/orders & Commandes du client connecté. \\
\hline
GET & /api/v1/admin/dashboard & Indicateurs du tableau de bord. \\
\hline
\multicolumn{3}{|c|}{\textbf{Retours}} \\
\hline
GET, POST & /api/v1/profile/returns & Lister, créer ses demandes de retour. \\
\hline
GET & /api/v1/admin/returns & Lister les retours. \\
\hline
PATCH & /api/v1/admin/returns/\{id\}/status & Accepter, refuser ou rembourser un retour. \\
\hline
GET, PUT & /api/v1/admin/returns/policy & Consulter, modifier la politique de retour. \\
\hline
GET & /api/v1/public/return-policy & Politique de retour publique. \\
\hline
\multicolumn{3}{|c|}{\textbf{TVA et livraison --- /api/v1/admin/tva-shipping}} \\
\hline
GET, PUT & /api/v1/admin/tva-shipping/config & Configuration générale de la TVA. \\
\hline
GET, POST & /api/v1/admin/tva-shipping/rates & Lister, créer un taux de TVA. \\
\hline
PUT, DELETE & /api/v1/admin/tva-shipping/rates/\{id\} & Modifier, supprimer un taux. \\
\hline
GET, POST & /api/v1/admin/tva-shipping/zones & Lister, créer une zone de livraison. \\
\hline
PUT, DELETE & /api/v1/admin/tva-shipping/zones/\{id\} & Modifier, supprimer une zone. \\
\hline
\multicolumn{3}{|c|}{\textbf{Avis --- catalog-service}} \\
\hline
GET & /api/v1/public/reviews/product/\{productId\} & Avis publiés d'un produit. \\
\hline
POST & /api/v1/profile/reviews & Publier un avis. \\
\hline
GET & /api/v1/reviews/has-reviewed & Savoir si le client a déjà noté le produit. \\
\hline
GET & /api/v1/admin/reviews & Lister les avis à modérer. \\
\hline
PATCH & /api/v1/admin/reviews/\{id\}/statut & Publier ou masquer un avis. \\
\hline
PATCH & /api/v1/admin/reviews/\{id\}/reponse & Répondre à un avis. \\
\hline
\caption{Services web du Sprint 6}
\label{tab:ws-sprint6}
\end{longtable}

\section{Les cas d'utilisation du Sprint 6}
\subsection{Cas d'utilisation « Traiter une demande de retour »}
\subsubsection{Description textuelle}
Le tableau~\ref{tab:uc-retour} présente la description textuelle de ce cas.
\renewcommand{\arraystretch}{1.5}
\begin{longtable}{|p{4cm}|p{11cm}|}
\hline
\textbf{Acteurs} & Client, Commerçant \\
\hline
\textbf{Objectif} & Permettre au client de retourner un article et au commerçant de traiter la demande jusqu'au remboursement. \\
\hline
\textbf{Pré-condition} & La commande est livrée et le délai de retour de la politique de la boutique n'est pas dépassé. \\
\hline
\textbf{Scénario principal} &
1. Le client ouvre « Mes commandes », choisit une commande livrée et clique sur « Retourner ».\newline
2. Il sélectionne les articles, indique le motif et valide.\newline
3. Le système vérifie le délai et crée la demande au statut « en attente ».\newline
4. Le commerçant consulte la demande dans « Retours » et l'accepte.\newline
5. À la réception des articles, il marque le retour comme remboursé ; le système met à jour la commande et l'historique des remboursements.\newline
6. Le client suit l'avancement de sa demande dans « Mes retours ». \\
\hline
\textbf{Scénario alternatif} &
\textbf{Délai dépassé :} à l'étape 3, le système refuse la demande et rappelle la politique de retour.\newline
\textbf{Retour refusé :} à l'étape 4, le commerçant refuse la demande en indiquant la raison, communiquée au client. \\
\hline
\caption{Description textuelle du cas « Traiter une demande de retour »}
\label{tab:uc-retour}
\end{longtable}

\subsubsection{Diagramme de séquence système}
\Needspace{18\baselineskip}
\begin{figure}[H]
    \centering
    \img{0.7}{dss_retour.png}
    \caption{Diagramme de séquence système du cas « Traiter une demande de retour »}
    \label{fig:dss-retour}
\end{figure}

\subsubsection{Diagramme de séquence objet}
\Needspace{20\baselineskip}
\begin{figure}[H]
    \centering
    \img{0.9}{dso_retour.png}
    \caption{Diagramme de séquence objet du cas « Traiter une demande de retour »}
    \label{fig:dso-retour}
\end{figure}

\section{Réalisation du Sprint 6}
\Needspace{20\baselineskip}
\textbf{Gestion des commandes :} la figure~\ref{fig:ui-commandes} présente la liste des commandes et la figure~\ref{fig:ui-dashboard} le tableau de bord des ventes.
\begin{figure}[H]
    \centering
    \img{0.85}{ui_commandes.png}
    \caption{Gestion des commandes}
    \label{fig:ui-commandes}
\end{figure}
\begin{figure}[H]
    \centering
    \img{0.85}{ui_dashboard.png}
    \caption{Tableau de bord des ventes}
    \label{fig:ui-dashboard}
\end{figure}

\Needspace{20\baselineskip}
\textbf{Retours, TVA et livraison :} les figures~\ref{fig:ui-retours} et~\ref{fig:ui-tva} présentent le traitement des retours et le paramétrage de la TVA et des zones de livraison.
\begin{figure}[H]
    \centering
    \img{0.85}{ui_retours.png}
    \caption{Traitement des retours}
    \label{fig:ui-retours}
\end{figure}
\begin{figure}[H]
    \centering
    \img{0.85}{ui_tva_livraison.png}
    \caption{Paramétrage de la TVA et des zones de livraison}
    \label{fig:ui-tva}
\end{figure}

\Needspace{20\baselineskip}
\textbf{Avis clients :} la figure~\ref{fig:ui-avis} présente la modération des avis.
\begin{figure}[H]
    \centering
    \img{0.85}{ui_avis.png}
    \caption{Modération des avis clients}
    \label{fig:ui-avis}
\end{figure}

\subsection*{Conclusion}
Cette release a rendu la plateforme pleinement transactionnelle : un client peut acheter, payer et retourner un article, et le commerçant pilote ses ventes et son paramétrage fiscal et logistique. La release suivante apporte les outils marketing et de fidélisation.

\thispagestyle{fancy}
\fancyhead[R]{A.U 2025-2026}
\fancyhead[L]{ chapitre 6: Release 4 : Marketing et fidélisation }
\chapter{Release 4 : Marketing et fidélisation}
\section*{Introduction}
Vendre ne suffit pas : un commerçant doit aussi attirer de nouveaux clients et faire revenir les anciens. Ce chapitre présente la quatrième release, qui lui en donne les moyens. Le Sprint 7 couvre les promotions, les bannières ciblées et l'e-mailing ; le Sprint 8 couvre le programme de fidélité à points et ses niveaux.

%==================================================================
\section{Backlog du Sprint 7}
\Needspace{18\baselineskip}
Le tableau~\ref{tab:backlog-sprint7} présente le \textbf{backlog du Sprint 7}.

\renewcommand{\arraystretch}{1.05}
\setlength{\tabcolsep}{3pt}
\begin{longtable}{|>{\raggedright\arraybackslash}p{4cm}|>{\raggedright\arraybackslash}p{9cm}|>{\centering\arraybackslash}p{1.2cm}|}
\hline
\textbf{User Story} & \textbf{Tâches} & \textbf{Estimation} \\
\hline
\endfirsthead
\hline
\textbf{User Story} & \textbf{Tâches} & \textbf{Estimation} \\
\hline
\endhead
\hline
\multicolumn{3}{r}{\textit{Suite à la page suivante}} \\
\endfoot
\hline
\endlastfoot

\textbf{En tant que commerçant, je souhaite créer des codes promo}
& Développer l'entité \texttt{Coupon} : type (pourcentage, montant fixe, livraison offerte, cadeau, « un acheté, un offert »), valeur, montant minimum, période et plage horaire, limites globale et par client, segment, catégories et produits ciblés. & 1 \\
\cline{2-3}
& Développer l'API CRUD, l'activation / désactivation et le suivi des utilisations. & 1 \\
\cline{2-3}
& Calculer les statistiques par code (commandes, revenus, taux de conversion). & 1 \\
\cline{2-3}
& Développer l'interface « Promotions ». & 1 \\
\cline{2-3}
& Tester la fonctionnalité. & 1 \\
\hline

\textbf{En tant que commerçant, je souhaite créer des remises automatiques}
& Développer l'API des remises sur un produit ou une catégorie, avec période de validité. & 1 \\
\cline{2-3}
& Afficher le prix barré et le badge de remise sur la vitrine. & 1 \\
\cline{2-3}
& Tester la fonctionnalité. & 1 \\
\hline

\textbf{En tant que commerçant, je souhaite publier des bannières ciblées}
& Développer l'API des bannières : titre, images ordinateur et mobile, vidéo, bouton d'action, alignement, animation, durée, planning. & 1 \\
\cline{2-3}
& Cibler une bannière par segment de clientèle et par type d'appareil. & 1 \\
\cline{2-3}
& Développer le formulaire de bannière avec aperçu ordinateur et mobile. & 1 \\
\cline{2-3}
& Afficher les bannières dans le carrousel d'accueil de chaque gabarit. & 1 \\
\hline

\textbf{En tant que commerçant, je souhaite envoyer des e-mails à mes clients}
& Intégrer l'API d'envoi transactionnel de Brevo (clé en variable d'environnement). & 1 \\
\cline{2-3}
& Développer l'envoi d'une newsletter à un segment et le contact individuel d'un client. & 1 \\
\cline{2-3}
& Développer l'interface « E-mail marketing » et ses statistiques. & 1 \\
\hline

\multicolumn{2}{|r|}{\textbf{Total}} & \textbf{15} \\
\hline
\caption{Backlog du Sprint 7}
\label{tab:backlog-sprint7}
\end{longtable}

\section{Diagramme de cas d'utilisation du Sprint 7}
\Needspace{20\baselineskip}
\begin{figure}[H]
    \centering
    \img{0.8}{cu_sprint7.png}
    \caption{Diagramme de cas d'utilisation du Sprint 7}
    \label{fig:cu-sprint7}
\end{figure}

\section{Services web du Sprint 7}
Le tableau~\ref{tab:ws-sprint7} présente les services web du Sprint 7.
\renewcommand{\arraystretch}{1.15}
\setlength{\tabcolsep}{4pt}
\begin{longtable}{|p{1.5cm}|p{6.8cm}|p{6.2cm}|}
\hline
\textbf{Méthode} & \textbf{URL} & \textbf{Description} \\
\hline
\endfirsthead
\hline
\textbf{Méthode} & \textbf{URL} & \textbf{Description} \\
\hline
\endhead
\multicolumn{3}{|c|}{\textbf{Promotions --- /api/v1/admin/promotions (order-service)}} \\
\hline
GET & /api/v1/admin/promotions/stats & Statistiques des promotions. \\
\hline
GET, POST & /api/v1/admin/promotions/coupons & Lister, créer un code promo. \\
\hline
PUT, DELETE & /api/v1/admin/promotions/coupons/\{id\} & Modifier, supprimer un code promo. \\
\hline
PATCH & /api/v1/admin/promotions/coupons/\{id\}/toggle & Activer ou désactiver un code. \\
\hline
GET, POST & /api/v1/admin/promotions/discounts & Lister, créer une remise. \\
\hline
PUT, DELETE & /api/v1/admin/promotions/discounts/\{id\} & Modifier, supprimer une remise. \\
\hline
PATCH & /api/v1/admin/promotions/discounts/\{id\}/toggle & Activer ou désactiver une remise. \\
\hline
\multicolumn{3}{|c|}{\textbf{Bannières --- marketing-service}} \\
\hline
GET, POST & /api/v1/admin/banners & Lister, créer une bannière. \\
\hline
PUT, DELETE & /api/v1/admin/banners/\{id\} & Modifier, supprimer une bannière. \\
\hline
PATCH & /api/v1/admin/banners/\{id\}/toggle & Activer ou désactiver une bannière. \\
\hline
GET & /api/v1/public/banners?position=\&segment= & Bannières actives pour une position et un segment. \\
\hline
\multicolumn{3}{|c|}{\textbf{E-mail marketing --- /api/v1/admin/email}} \\
\hline
GET & /api/v1/admin/email/stats & Statistiques d'envoi. \\
\hline
POST & /api/v1/admin/email/newsletter & Envoyer une newsletter. \\
\hline
POST & /api/v1/admin/email/contact/\{userId\} & Contacter un client. \\
\hline
\caption{Services web du Sprint 7}
\label{tab:ws-sprint7}
\end{longtable}

\section{Les cas d'utilisation du Sprint 7}
\subsection{Cas d'utilisation « Créer un code promo ciblé »}
\subsubsection{Description textuelle}
Le tableau~\ref{tab:uc-coupon} présente la description textuelle de ce cas.
\renewcommand{\arraystretch}{1.5}
\begin{longtable}{|p{4cm}|p{11cm}|}
\hline
\textbf{Acteurs} & Commerçant \\
\hline
\textbf{Objectif} & Créer une offre réservée à une cible et en mesurer l'efficacité. \\
\hline
\textbf{Pré-condition} & Le commerçant est connecté et sa boutique est active. \\
\hline
\textbf{Scénario principal} &
1. Le commerçant ouvre « Promotions » puis « Nouveau code ».\newline
2. Il saisit le code (par exemple \texttt{VIP20}), le type et la valeur de la remise, et le montant minimum de commande.\newline
3. Il définit la période, éventuellement une plage horaire, et les limites d'utilisation.\newline
4. Il restreint le code à un segment (par exemple VIP) et, s'il le souhaite, à des catégories ou des produits.\newline
5. Le système vérifie que le code n'existe pas déjà, calcule son statut (programmé, actif ou expiré) à partir des dates, et l'enregistre.\newline
6. Lorsqu'un client l'utilise, le système enregistre l'utilisation et met à jour le nombre de commandes, les revenus et le taux de conversion du code. \\
\hline
\textbf{Scénario alternatif} &
\textbf{Code déjà existant :} à l'étape 5, le système demande un autre code.\newline
\textbf{Limite atteinte :} à l'étape 6, si le client a déjà utilisé le code le nombre maximal de fois autorisé, le système refuse la remise avec un message explicite. \\
\hline
\caption{Description textuelle du cas « Créer un code promo ciblé »}
\label{tab:uc-coupon}
\end{longtable}

\subsubsection{Diagramme de séquence système}
\Needspace{18\baselineskip}
\begin{figure}[H]
    \centering
    \img{0.65}{dss_coupon.png}
    \caption{Diagramme de séquence système du cas « Créer un code promo ciblé »}
    \label{fig:dss-coupon}
\end{figure}

\subsubsection{Diagramme de séquence objet}
\Needspace{20\baselineskip}
\begin{figure}[H]
    \centering
    \img{0.9}{dso_coupon.png}
    \caption{Diagramme de séquence objet du cas « Créer un code promo ciblé »}
    \label{fig:dso-coupon}
\end{figure}

\section{Réalisation du Sprint 7}
\Needspace{20\baselineskip}
\textbf{Promotions :} la figure~\ref{fig:ui-promotions} présente la gestion des codes promo et des remises, avec leurs indicateurs.
\begin{figure}[H]
    \centering
    \img{0.85}{ui_promotions.png}
    \caption{Gestion des promotions}
    \label{fig:ui-promotions}
\end{figure}

\Needspace{20\baselineskip}
\textbf{Bannières :} la figure~\ref{fig:ui-bannieres} présente le formulaire de bannière, avec aperçu sur ordinateur et sur mobile.
\begin{figure}[H]
    \centering
    \img{0.85}{ui_bannieres.png}
    \caption{Création d'une bannière ciblée}
    \label{fig:ui-bannieres}
\end{figure}

\Needspace{20\baselineskip}
\textbf{E-mail marketing :} la figure~\ref{fig:ui-email} présente l'envoi d'une newsletter.
\begin{figure}[H]
    \centering
    \img{0.85}{ui_email_marketing.png}
    \caption{E-mail marketing}
    \label{fig:ui-email}
\end{figure}

%==================================================================
\section{Backlog du Sprint 8}
\Needspace{18\baselineskip}
Le tableau~\ref{tab:backlog-sprint8} présente le \textbf{backlog du Sprint 8}, consacré au programme de fidélité.

\renewcommand{\arraystretch}{1.05}
\setlength{\tabcolsep}{3pt}
\begin{longtable}{|>{\raggedright\arraybackslash}p{4cm}|>{\raggedright\arraybackslash}p{9cm}|>{\centering\arraybackslash}p{1.2cm}|}
\hline
\textbf{User Story} & \textbf{Tâches} & \textbf{Estimation} \\
\hline
\endfirsthead
\hline
\textbf{User Story} & \textbf{Tâches} & \textbf{Estimation} \\
\hline
\endhead
\hline
\multicolumn{3}{r}{\textit{Suite à la page suivante}} \\
\endfoot
\hline
\endlastfoot

\textbf{En tant que commerçant, je souhaite configurer le gain de points}
& Développer l'entité \texttt{LoyaltyConfig} (points par dinar, bienvenue, avis, anniversaire, promotion automatique). & 1 \\
\cline{2-3}
& Attribuer les points à la livraison d'une commande, à l'inscription et à la publication d'un avis, avec le multiplicateur du niveau. & 1 \\
\cline{2-3}
& Journaliser chaque mouvement de points (\texttt{PointsTransaction}). & 1 \\
\cline{2-3}
& Développer l'onglet « Configuration ». & 1 \\
\hline

\textbf{En tant que commerçant, je souhaite gérer les niveaux et ajuster les points d'un client}
& Développer l'API CRUD des segments (seuil de points, multiplicateur, avantages). & 1 \\
\cline{2-3}
& Développer la montée de niveau automatique après chaque attribution de points. & 1 \\
\cline{2-3}
& Développer l'ajustement manuel (le solde ne peut pas devenir négatif) et le classement des clients. & 1 \\
\cline{2-3}
& Développer les onglets « Avantages / Niveaux », « Classement » et « Ajustement manuel ». & 1 \\
\cline{2-3}
& Tester la fonctionnalité. & 1 \\
\hline

\textbf{En tant que client, je souhaite consulter mes points et mon niveau}
& Développer l'API \texttt{/profile/loyalty} (solde, niveau, historique, progression vers le niveau suivant). & 1 \\
\cline{2-3}
& Afficher la carte de fidélité dans « Mon profil ». & 1 \\
\hline

\multicolumn{2}{|r|}{\textbf{Total}} & \textbf{11} \\
\hline
\caption{Backlog du Sprint 8}
\label{tab:backlog-sprint8}
\end{longtable}

\section{Diagramme de cas d'utilisation du Sprint 8}
\Needspace{20\baselineskip}
\begin{figure}[H]
    \centering
    \img{0.8}{cu_sprint8.png}
    \caption{Diagramme de cas d'utilisation du Sprint 8}
    \label{fig:cu-sprint8}
\end{figure}

\section{Services web du Sprint 8}
Le tableau~\ref{tab:ws-sprint8} présente les services web du Sprint 8.
\renewcommand{\arraystretch}{1.15}
\setlength{\tabcolsep}{4pt}
\begin{longtable}{|p{1.5cm}|p{6.8cm}|p{6.2cm}|}
\hline
\textbf{Méthode} & \textbf{URL} & \textbf{Description} \\
\hline
\endfirsthead
\hline
\textbf{Méthode} & \textbf{URL} & \textbf{Description} \\
\hline
\endhead
\multicolumn{3}{|c|}{\textbf{Fidélité --- /api/v1/admin/loyalty (auth-service)}} \\
\hline
GET & /api/v1/admin/loyalty/config & Consulter la configuration du programme. \\
\hline
PUT & /api/v1/admin/loyalty/config & Modifier la configuration. \\
\hline
GET & /api/v1/admin/loyalty/leaderboard & Classement des clients par points. \\
\hline
POST & /api/v1/admin/loyalty/users/\{userId\}/adjust-points & Ajuster les points d'un client (positif ou négatif). \\
\hline
\multicolumn{3}{|c|}{\textbf{Segments --- /api/v1/admin/segments}} \\
\hline
GET, POST & /api/v1/admin/segments & Lister, créer un niveau. \\
\hline
PUT, DELETE & /api/v1/admin/segments/\{id\} & Modifier, supprimer un niveau. \\
\hline
GET & /api/v1/admin/users/by-segment/\{segmentName\} & Clients d'un niveau. \\
\hline
\multicolumn{3}{|c|}{\textbf{Espace client}} \\
\hline
GET & /api/v1/profile/loyalty & Solde, niveau et historique du client connecté. \\
\hline
\caption{Services web du Sprint 8}
\label{tab:ws-sprint8}
\end{longtable}

\section{Les cas d'utilisation du Sprint 8}
\subsection{Cas d'utilisation « Gagner des points et changer de niveau »}
\subsubsection{Description textuelle}
Le tableau~\ref{tab:uc-fidelite} présente la description textuelle de ce cas.
\renewcommand{\arraystretch}{1.5}
\begin{longtable}{|p{4cm}|p{11cm}|}
\hline
\textbf{Acteurs} & Client, Commerçant \\
\hline
\textbf{Objectif} & Récompenser automatiquement les achats et faire progresser le client dans les niveaux du programme. \\
\hline
\textbf{Pré-condition} & Le client est inscrit ; le programme de fidélité est configuré. \\
\hline
\textbf{Scénario principal} &
1. Le commerçant passe une commande du client au statut « livrée ».\newline
2. Le système calcule les points : $\text{points} = \text{arrondi}(\text{montant} \times \text{points par dinar} \times \text{multiplicateur du niveau})$. Par exemple, 120 DT avec 1 point par dinar et un multiplicateur de 1,5 donnent 180 points.\newline
3. Le système crédite le solde et journalise le mouvement.\newline
4. Si la promotion automatique est active, le système cherche le niveau le plus élevé dont le seuil est atteint et y fait passer le client.\newline
5. Le client consulte son solde, son niveau et sa progression dans « Mon profil ». \\
\hline
\textbf{Scénario alternatif} &
\textbf{Ajustement manuel :} le commerçant retire des points ; le solde est borné à zéro.\newline
\textbf{Pas de rétrogradation automatique :} une baisse du solde ne fait pas redescendre le client ; le niveau « Inactif » n'est attribué que manuellement. \\
\hline
\caption{Description textuelle du cas « Gagner des points et changer de niveau »}
\label{tab:uc-fidelite}
\end{longtable}

\subsubsection{Diagramme de séquence système}
\Needspace{18\baselineskip}
\begin{figure}[H]
    \centering
    \img{0.65}{dss_fidelite.png}
    \caption{Diagramme de séquence système du cas « Gagner des points et changer de niveau »}
    \label{fig:dss-fidelite}
\end{figure}

\subsubsection{Diagramme de séquence objet}
\Needspace{20\baselineskip}
\begin{figure}[H]
    \centering
    \img{0.9}{dso_fidelite.png}
    \caption{Diagramme de séquence objet du cas « Gagner des points et changer de niveau »}
    \label{fig:dso-fidelite}
\end{figure}

\section{Réalisation du Sprint 8}
\Needspace{20\baselineskip}
\textbf{Programme de fidélité :} les figures~\ref{fig:ui-fidelite-config} et~\ref{fig:ui-fidelite-niveaux} présentent la configuration du programme et la gestion des niveaux (Nouveau, Fidèle, VIP, Inactif).
\begin{figure}[H]
    \centering
    \img{0.85}{ui_fidelite_config.png}
    \caption{Configuration du programme de fidélité}
    \label{fig:ui-fidelite-config}
\end{figure}
\begin{figure}[H]
    \centering
    \img{0.85}{ui_fidelite_niveaux.png}
    \caption{Niveaux et avantages du programme de fidélité}
    \label{fig:ui-fidelite-niveaux}
\end{figure}

\Needspace{20\baselineskip}
\textbf{Espace client :} la figure~\ref{fig:ui-fidelite-client} présente la carte de fidélité du client dans son profil.
\begin{figure}[H]
    \centering
    \img{0.75}{ui_fidelite_client.png}
    \caption{Carte de fidélité dans l'espace client}
    \label{fig:ui-fidelite-client}
\end{figure}

\subsection*{Conclusion}
Cette release a donné au commerçant les leviers classiques du marketing en ligne : promotions ciblées, bannières, e-mails et fidélité. Ces outils reposent toutefois sur des segments définis à la main (Nouveau, Fidèle, VIP). Les deux dernières releases vont plus loin grâce à l'intelligence artificielle : l'essayage virtuel d'abord, puis l'analyse automatique du comportement des visiteurs.

\thispagestyle{fancy}
\fancyhead[R]{A.U 2025-2026}
\fancyhead[L]{ chapitre 7: Release 5 : Essayage virtuel }
\chapter{Release 5 : Essayage virtuel en temps réel par IA générative}
\section*{Introduction}
Dans la mode en ligne, la première source d'hésitation est simple : le client ne sait pas à quoi ressemblera le vêtement sur lui. Cette incertitude freine l'achat et multiplie les retours. Ce chapitre présente la cinquième release, qui répond à ce problème par une cabine d'essayage virtuel : depuis la fiche d'un vêtement, le client active sa caméra et se voit, en direct, portant l'article. Nous présentons d'abord l'architecture de la fonctionnalité, puis le choix du modèle et sa configuration, avant de détailler le Sprint 9.

\section{Architecture de l'essayage virtuel}
\subsection{Flux d'architecture}
La figure~\ref{fig:archi-tryon} présente le flux de l'essayage virtuel. Il fait intervenir trois parties : le navigateur du client, notre microservice catalogue et l'API de Decart.
\begin{figure}[H]
    \centering
    \img{0.95}{architecture_tryon.png}
    \caption{Architecture de l'essayage virtuel}
    \label{fig:archi-tryon}
\end{figure}

Le flux se déroule en quatre temps :
\begin{enumerate}
  \item \textbf{Obtention d'un jeton :} le navigateur demande un jeton à \texttt{catalog-service} (\texttt{POST /api/v1/public/tryon/token}). Le service, seul détenteur de la clé secrète Decart, demande à Decart un \textbf{jeton client éphémère} valable dix minutes et le renvoie au navigateur. La clé secrète n'est donc jamais exposée côté client.
  \item \textbf{Préparation des entrées :} le navigateur ouvre la caméra (1280 × 720) et télécharge l'image du vêtement, qu'il redimensionne (768 pixels au plus sur le grand côté) et compresse en JPEG pour limiter la latence.
  \item \textbf{Connexion temps réel :} le SDK Decart ouvre une connexion \textbf{WebRTC} vers le modèle \textit{Lucy VTON}. Le flux caméra part vers Decart ; l'image du vêtement et une consigne textuelle sont envoyées sur la même session.
  \item \textbf{Restitution :} Decart renvoie un flux vidéo dans lequel le haut porté par le client est remplacé par le vêtement ; le navigateur l'affiche à côté de l'image de la caméra.
\end{enumerate}

\subsection{Couche de présentation}
Le composant React \texttt{TryOnModal} gère l'ensemble de la session : états successifs (« Jeton Decart… », « Caméra… », « Connexion au modèle… », « Essayage en direct »), affichage des deux vidéos, compte à rebours de la session et messages d'erreur compréhensibles. Le bouton « Essayer virtuellement » n'apparaît que sur les fiches des boutiques du secteur vêtements.

\subsection{Couche d'intermédiation}
Le service \texttt{DecartTryOnService} du microservice catalogue appelle \texttt{POST https://api.decart.ai/v1/client/tokens} avec la clé secrète lue dans une variable d'environnement, et renvoie au navigateur le jeton, sa date d'expiration et la durée de session autorisée. Si la clé n'est pas configurée, il répond \texttt{503}, et en cas d'échec de Decart, \texttt{502}, ce qui permet à l'interface d'afficher un message adapté.

\subsection{Couche de génération}
La génération est entièrement réalisée par Decart. Le modèle \textit{Lucy VTON} est un modèle vidéo génératif temps réel : il transforme chaque image du flux caméra en conservant la personne, sa posture et l'arrière-plan, et en remplaçant le vêtement ciblé par celui de l'image de référence.

\section{Mise en œuvre de l'essayage virtuel}
\subsection{Critères de choix}
Le choix de la technologie a été guidé par cinq critères :
\begin{itemize}
  \item \textbf{Réalisme :} le rendu doit être crédible (plis, éclairage, morphologie), faute de quoi il dessert la vente ;
  \item \textbf{Interactivité :} un essayage en direct, où le client bouge et se tourne, est bien plus convaincant qu'une photo retouchée ;
  \item \textbf{Infrastructure :} la plateforme ne dispose pas de serveur GPU ;
  \item \textbf{Intégration :} la solution doit s'intégrer à une application web, sans installation côté client ;
  \item \textbf{Coût et sécurité :} coût maîtrisé par session et clé secrète jamais exposée.
\end{itemize}

\subsection{Comparaison des solutions candidates}
Le tableau~\ref{tab:comparaison-tryon} compare les quatre familles de solutions étudiées.
\renewcommand{\arraystretch}{1.4}
\begin{longtable}{|p{3.2cm}|p{4.2cm}|p{4.2cm}|p{2.6cm}|}
\hline
\textbf{Solution} & \textbf{Avantages} & \textbf{Contraintes} & \textbf{Décision} \\
\hline
\endfirsthead
\hline
\textbf{Solution} & \textbf{Avantages} & \textbf{Contraintes} & \textbf{Décision} \\
\hline
\endhead
Superposition 2D par détection de pose (MediaPipe~\cite{mediapipe}) &
Gratuit, exécuté dans le navigateur, temps réel. &
Rendu peu réaliste : le vêtement est « collé » sur l'image, sans plis ni occlusions ; image détourée à préparer pour chaque produit. &
Écartée (réalisme insuffisant). \\
\hline
Modèles de diffusion open source auto-hébergés (IDM-VTON~\cite{idmvton}, OOTDiffusion~\cite{ootdiffusion}) &
Très bon réalisme sur photo ; aucun coût de licence ; données maîtrisées. &
Nécessitent un GPU puissant ; plusieurs secondes par image ; photo statique, pas de vidéo. &
Écartée (pas de GPU, pas de temps réel). \\
\hline
API d'essayage image par image hébergées (par exemple FASHN~\cite{fashn}) &
Bon réalisme ; aucune infrastructure ; intégration simple par API REST. &
Photo statique : le client doit téléverser une photo et attendre le résultat ; facturation à l'image. &
Écartée (expérience moins interactive). \\
\hline
Decart \textit{Lucy VTON} temps réel~\cite{decart} &
Vidéo en direct via WebRTC ; aucune infrastructure ; SDK JavaScript ; jetons clients éphémères. &
Facturation à la seconde de génération ; dépendance à un fournisseur ; WebRTC bloqué par certains VPN ou réseaux. &
\textbf{Retenue.} \\
\hline
\caption{Comparaison des solutions d'essayage virtuel}
\label{tab:comparaison-tryon}
\end{longtable}

\subsection{Solution retenue}
Nous avons retenu \textbf{Decart} et son modèle \textbf{\textit{Lucy VTON}}. C'est la seule option qui combine un rendu génératif réaliste et un essayage \textbf{en direct}, sans serveur GPU de notre côté. Le mécanisme de jetons clients éphémères répond en outre au critère de sécurité : la clé secrète reste sur notre serveur. Ses contraintes, coût à la seconde et dépendance au fournisseur, ont été encadrées par la configuration décrite ci-dessous.

\subsection{Configuration du modèle}
Le tableau~\ref{tab:config-tryon} résume la configuration retenue et sa justification.
\renewcommand{\arraystretch}{1.3}
\begin{longtable}{|p{4.2cm}|p{3.6cm}|p{6.6cm}|}
\hline
\textbf{Paramètre} & \textbf{Valeur} & \textbf{Justification} \\
\hline
\endfirsthead
\hline
\textbf{Paramètre} & \textbf{Valeur} & \textbf{Justification} \\
\hline
\endhead
Modèle & \texttt{lucy-vton-latest} & Modèle d'essayage virtuel temps réel de Decart. \\
\hline
Résolution de la caméra & 1280 × 720 & Résolution attendue par le modèle ; lue dans sa description pour s'y adapter. \\
\hline
Image du vêtement & JPEG, 768 px max., qualité 0,85 & Réduit le volume envoyé sans perte visible de détails. \\
\hline
Consigne (\textit{prompt}) & « Substitute the current top with a \textit{<nom du produit>}, matching the color, knit, round neck and long sleeves of the reference garment image » & Indique au modèle quelle pièce remplacer et quelles caractéristiques respecter ; le nom du produit est inséré automatiquement. \\
\hline
Amélioration du prompt & désactivée & Garder une consigne prévisible et identique d'un produit à l'autre. \\
\hline
Durée d'une session & 45 secondes & Plafonne le coût par essai ; le compte à rebours ne démarre qu'à la première image générée. \\
\hline
Validité du jeton client & 600 secondes & Suffisant pour une session, inutilisable au-delà. \\
\hline
Délai de détection d'absence de vidéo & 20 secondes & Au-delà, un message explique la cause probable (VPN bloquant WebRTC, crédits épuisés). \\
\hline
\caption{Configuration de l'essayage virtuel}
\label{tab:config-tryon}
\end{longtable}

\subsection{Gestion des erreurs}
Une fonctionnalité qui dépend d'une caméra, d'un réseau temps réel et d'un service externe peut échouer de nombreuses manières. Chaque cause est traduite en un message clair : caméra refusée, crédits Decart épuisés, image refusée par la modération, session coupée pour violation de politique, ou absence de vidéo (souvent due à un VPN, WebRTC ayant besoin d'UDP). Un problème rencontré pendant le développement mérite d'être signalé : un module interne du SDK, chargé seulement de mesurer la latence, faisait échouer le flux vidéo avec le serveur de développement Vite ; il a été neutralisé par un module de substitution (\textit{stub}) sans effet sur la génération.

%==================================================================
\section{Backlog du Sprint 9}
\Needspace{18\baselineskip}
Le tableau~\ref{tab:backlog-sprint9} présente le \textbf{backlog du Sprint 9}.

\renewcommand{\arraystretch}{1.05}
\setlength{\tabcolsep}{3pt}
\begin{longtable}{|>{\raggedright\arraybackslash}p{4cm}|>{\raggedright\arraybackslash}p{9cm}|>{\centering\arraybackslash}p{1.2cm}|}
\hline
\textbf{User Story} & \textbf{Tâches} & \textbf{Estimation} \\
\hline
\endfirsthead
\hline
\textbf{User Story} & \textbf{Tâches} & \textbf{Estimation} \\
\hline
\endhead
\hline
\multicolumn{3}{r}{\textit{Suite à la page suivante}} \\
\endfoot
\hline
\endlastfoot

\textbf{En tant que visiteur, je souhaite essayer un vêtement en direct avec ma caméra}
& Étudier et comparer les solutions d'essayage virtuel. & 1 \\
\cline{2-3}
& Ouvrir la caméra à la résolution du modèle et gérer le refus d'autorisation. & 1 \\
\cline{2-3}
& Préparer l'image du vêtement (téléchargement, redimensionnement, compression). & 1 \\
\cline{2-3}
& Connecter le SDK Decart en WebRTC et afficher le flux généré. & 1 \\
\cline{2-3}
& Construire la consigne à partir du nom du produit. & 1 \\
\cline{2-3}
& Limiter la session à 45 secondes avec compte à rebours et fermeture automatique. & 1 \\
\cline{2-3}
& Traduire les erreurs (caméra, crédits, modération, réseau) en messages clairs. & 1 \\
\cline{2-3}
& Afficher le bouton « Essayer virtuellement » sur les fiches des boutiques de vêtements. & 1 \\
\cline{2-3}
& Tester sur plusieurs vêtements et conditions d'éclairage. & 1 \\
\hline

\textbf{En tant que système, je veux délivrer des jetons Decart à durée limitée sans exposer la clé}
& Développer \texttt{DecartTryOnService} et l'API \texttt{/api/v1/public/tryon/token}. & 1 \\
\cline{2-3}
& Ajouter la route à la passerelle et un intermédiaire de secours dans le serveur de développement. & 1 \\
\cline{2-3}
& Tester la fonctionnalité (clé absente, erreur Decart). & 1 \\
\hline

\multicolumn{2}{|r|}{\textbf{Total}} & \textbf{12} \\
\hline
\caption{Backlog du Sprint 9}
\label{tab:backlog-sprint9}
\end{longtable}

\section{Diagramme de cas d'utilisation du Sprint 9}
\Needspace{20\baselineskip}
\begin{figure}[H]
    \centering
    \img{0.75}{cu_sprint9.png}
    \caption{Diagramme de cas d'utilisation du Sprint 9}
    \label{fig:cu-sprint9}
\end{figure}

\section{Services web du Sprint 9}
Le tableau~\ref{tab:ws-sprint9} présente les services web du Sprint 9, y compris les appels à l'API externe de Decart.
\renewcommand{\arraystretch}{1.15}
\setlength{\tabcolsep}{4pt}
\begin{longtable}{|p{1.7cm}|p{6.6cm}|p{6.2cm}|}
\hline
\textbf{Méthode} & \textbf{URL} & \textbf{Description} \\
\hline
\endfirsthead
\hline
\textbf{Méthode} & \textbf{URL} & \textbf{Description} \\
\hline
\endhead
\multicolumn{3}{|c|}{\textbf{API Sellio --- catalog-service}} \\
\hline
POST & /api/v1/public/tryon/token & Obtenir un jeton client Decart éphémère, sa date d'expiration et la durée de session. \\
\hline
\multicolumn{3}{|c|}{\textbf{API externe Decart}} \\
\hline
POST & https://api.decart.ai/v1/client/tokens & Créer un jeton client (appelé par le serveur avec la clé secrète). \\
\hline
WebRTC & \texttt{client.realtime.connect(stream, \{model\})} (SDK) & Ouvrir la session temps réel et envoyer le flux caméra. \\
\hline
SDK & \texttt{realtime.set(\{prompt, image\})} & Envoyer la consigne et l'image du vêtement. \\
\hline
\caption{Services web du Sprint 9}
\label{tab:ws-sprint9}
\end{longtable}

\section{Les cas d'utilisation du Sprint 9}
\subsection{Cas d'utilisation « Essayer un vêtement virtuellement »}
\subsubsection{Description textuelle}
Le tableau~\ref{tab:uc-tryon} présente la description textuelle de ce cas.
\renewcommand{\arraystretch}{1.5}
\begin{longtable}{|p{4cm}|p{11cm}|}
\hline
\textbf{Acteurs} & Visiteur ou Client, Decart \\
\hline
\textbf{Objectif} & Voir en direct le vêtement porté sur soi avant de l'acheter. \\
\hline
\textbf{Pré-condition} & La boutique est du secteur vêtements ; le produit a une image ; l'appareil dispose d'une caméra. \\
\hline
\textbf{Scénario principal} &
1. Le visiteur clique sur « Essayer virtuellement » depuis la fiche produit.\newline
2. Le système demande un jeton éphémère au service catalogue, qui l'obtient auprès de Decart.\newline
3. Le navigateur demande l'accès à la caméra ; le visiteur l'autorise.\newline
4. Le système prépare l'image du vêtement et ouvre la session WebRTC avec le modèle \textit{Lucy VTON}.\newline
5. Le système envoie la consigne et l'image du vêtement.\newline
6. Decart renvoie la vidéo générée ; le système l'affiche et lance le compte à rebours de 45 secondes.\newline
7. Le visiteur bouge et se tourne pour observer le rendu.\newline
8. À la fin du compte à rebours ou à la fermeture, le système coupe la session et libère la caméra. \\
\hline
\textbf{Scénario alternatif} &
\textbf{Caméra refusée :} à l'étape 3, le système affiche « Caméra refusée. Autorisez-la dans le navigateur, puis réessayez. »\newline
\textbf{Service indisponible :} à l'étape 2, si la clé n'est pas configurée ou si Decart échoue, le système affiche que l'essayage est indisponible.\newline
\textbf{Aucune vidéo après 20 secondes :} à l'étape 6, le système suggère de désactiver le VPN ou signale des crédits épuisés. \\
\hline
\caption{Description textuelle du cas « Essayer un vêtement virtuellement »}
\label{tab:uc-tryon}
\end{longtable}

\subsubsection{Diagramme de séquence système}
\Needspace{18\baselineskip}
\begin{figure}[H]
    \centering
    \img{0.7}{dss_tryon.png}
    \caption{Diagramme de séquence système du cas « Essayer un vêtement virtuellement »}
    \label{fig:dss-tryon}
\end{figure}

\subsubsection{Diagramme de séquence objet}
\Needspace{20\baselineskip}
\begin{figure}[H]
    \centering
    \img{0.95}{dso_tryon.png}
    \caption{Diagramme de séquence objet du cas « Essayer un vêtement virtuellement »}
    \label{fig:dso-tryon}
\end{figure}

\section{Réalisation du Sprint 9}
\Needspace{20\baselineskip}
\textbf{Bouton d'essayage :} la figure~\ref{fig:ui-tryon-btn} présente la fiche d'un vêtement avec le bouton « Essayer virtuellement ».
\begin{figure}[H]
    \centering
    \img{0.8}{ui_tryon_bouton.png}
    \caption{Accès à l'essayage virtuel depuis la fiche produit}
    \label{fig:ui-tryon-btn}
\end{figure}

\Needspace{20\baselineskip}
\textbf{Cabine d'essayage :} la figure~\ref{fig:ui-tryon} présente la cabine en fonctionnement : à gauche l'image de la caméra, à droite la vidéo générée par Decart, avec le compte à rebours.
\begin{figure}[H]
    \centering
    \img{0.9}{ui_tryon_session.png}
    \caption{Essayage virtuel en direct}
    \label{fig:ui-tryon}
\end{figure}

\section*{Conclusion}
Cette release a ajouté à Sellio une fonctionnalité rarement proposée par défaut par les plateformes de boutiques en ligne : un essayage virtuel en direct, fondé sur un modèle génératif temps réel et intégré de façon sécurisée grâce aux jetons éphémères. Le chapitre suivant présente l'autre volet intelligent de la plateforme : l'analyse du comportement des visiteurs par apprentissage automatique.

\thispagestyle{fancy}
\fancyhead[R]{A.U 2025-2026}
\fancyhead[L]{ chapitre 8: Release 6 : Intelligence comportementale }
\chapter{Release 6 : Intelligence comportementale et apprentissage automatique}
\section*{Introduction}
Ce dernier chapitre répond au cœur du sujet : analyser les interactions des utilisateurs et exploiter ces données par apprentissage automatique pour personnaliser l'expérience d'achat. Nous présentons d'abord l'architecture de ce volet, puis les données utilisées pour entraîner et évaluer les modèles et le protocole d'évaluation. Nous détaillons ensuite les trois modèles, recommandation, segmentation et prédiction de l'attrition, avec le choix de chaque algorithme, sa configuration et ses résultats. Nous terminons par les deux sprints de la release.

\section{Architecture de l'intelligence comportementale}
\subsection{Flux d'architecture}
La figure~\ref{fig:archi-ml} présente le flux complet, de la collecte d'un clic jusqu'à l'affichage d'une recommandation ou d'un score d'attrition.
\begin{figure}[H]
    \centering
    \img{1}{architecture_ml.png}
    \caption{Architecture de l'intelligence comportementale}
    \label{fig:archi-ml}
\end{figure}

\subsection{Couche de collecte}
Un module JavaScript de la vitrine (\texttt{tracker.js}) enregistre douze types d'événements : pages vues, catégories et produits consultés, clics, recherches, ajouts et retraits du panier, ajouts aux favoris, début de commande, achats et clics sur une recommandation. Chaque événement porte l'identifiant de la boutique, un \textbf{identifiant de visiteur anonyme} conservé dans le navigateur, un identifiant de session (une nouvelle session commence après 30 minutes d'inactivité) et, si le client est connecté, son identifiant. Les événements sont regroupés en lots et envoyés à \texttt{POST /api/v1/analytics/events}, pour ne pas ralentir la navigation. Le serveur n'accepte que les types connus, limite la taille des lots (50 événements) et le débit par visiteur, afin qu'un script ne puisse pas fausser les données.

\subsection{Couche de stockage}
Le volume d'événements croît beaucoup plus vite que celui des commandes. La table \texttt{user\_events} est donc conçue pour ce volume :
\begin{itemize}
  \item \textbf{partitionnement mensuel} par date : les requêtes sur une période ne lisent que les mois concernés, et la suppression d'un mois ancien est instantanée ;
  \item \textbf{index BRIN} sur la date : quelques kilo-octets suffisent à indexer des millions de lignes insérées dans l'ordre chronologique ;
  \item \textbf{index B-tree partiels} sur les seuls événements utiles aux modèles (produits consultés, achats) ;
  \item \textbf{agrégats quotidiens} pour le tableau de bord et \textbf{rétention de 13 mois}, suffisante pour comparer une saison à celle de l'année précédente.
\end{itemize}

\subsection{Couche d'entraînement}
Chaque nuit à 3 h 40, et à la demande du commerçant, le service \texttt{MlTrainingService} extrait les données de chaque boutique, les transmet à des scripts Python (\texttt{ml\_train.py}) exécutés en sous-processus, puis enregistre les résultats : voisins de chaque produit et règles d'association dans \texttt{product\_recommendations}, segment de chaque visiteur dans \texttt{behavior\_segments}, et métriques de chaque entraînement dans \texttt{ml\_model\_runs}. Les modèles d'attrition, entraînés hors ligne sur des données réelles, sont chargés par \texttt{predict.py}.

\subsection{Couche de service et de restitution}
Les recommandations sont lues dans des tables pré-calculées : servir une recommandation revient à une simple requête indexée, sans calcul au moment de l'affichage. La vitrine les affiche sur la fiche produit, sur l'accueil et dans le panier ; le back-office restitue segments, entonnoir de conversion et risques d'attrition dans la page « Comportement ».

\section{Données d'entraînement et d'évaluation}
\subsection{Pourquoi des jeux de données publics réels}
Une plateforme nouvelle n'a pas d'historique : ses premières boutiques ne comptent que quelques centaines d'événements, trop peu pour mesurer la qualité d'un modèle. Générer des données artificielles aurait été une impasse : un modèle évalué sur des données fabriquées retrouve surtout les règles qui ont servi à les fabriquer, et ses scores ne disent rien de son comportement face à de vrais clients. Nous avons donc entraîné et évalué les modèles sur \textbf{trois jeux de données publics réels}, choisis pour correspondre aux deux secteurs de Sellio, et nettoyés par des scripts Python reproductibles. Le tableau~\ref{tab:datasets} les présente.

\renewcommand{\arraystretch}{1.35}
\begin{longtable}{|p{3.1cm}|p{4.3cm}|p{3.4cm}|p{3.6cm}|}
\hline
\textbf{Jeu de données} & \textbf{Contenu} & \textbf{Volume brut} & \textbf{Usage dans Sellio} \\
\hline
\endfirsthead
\hline
\textbf{Jeu de données} & \textbf{Contenu} & \textbf{Volume brut} & \textbf{Usage dans Sellio} \\
\hline
\endhead
Online Retail II (UCI, licence CC BY 4.0)~\cite{onlineretail} & Transactions d'un e-commerçant britannique, déc. 2009 -- déc. 2011. & 1 067 371 lignes de factures & Attrition (modèle général), recommandation, « souvent achetés ensemble ». \\
\hline
eCommerce Events History in Cosmetics Shop (Kaggle, REES46)~\cite{cosmetics} & Événements de navigation d'une boutique de cosmétiques en ligne (vue, panier, retrait, achat). & 3 533 286 événements (décembre 2019) & Segmentation des visiteurs, recommandation (secteur cosmétiques). \\
\hline
H\&M Personalized Fashion Recommendations (Kaggle)~\cite{hm} & Achats, articles et clients d'une enseigne de mode, sept. 2018 -- sept. 2020. & 31,8 millions d'achats, 105 542 articles & Attrition (modèle vêtements), recommandation (secteur vêtements). \\
\hline
\caption{Jeux de données publics utilisés}
\label{tab:datasets}
\end{longtable}

\subsection{Nettoyage des données}
Aucun de ces jeux n'est directement exploitable : chacun contient des doublons, des valeurs manquantes, des valeurs aberrantes ou du trafic non humain. Chaque script de nettoyage produit un rapport qui détaille, étape par étape, le nombre de lignes avant et après. Le tableau~\ref{tab:nettoyage} en résume les principales étapes.

\renewcommand{\arraystretch}{1.3}
\begin{longtable}{|p{2.6cm}|p{8.6cm}|p{3.4cm}|}
\hline
\textbf{Jeu} & \textbf{Étape (lignes supprimées)} & \textbf{Résultat} \\
\hline
\endfirsthead
\hline
\textbf{Jeu} & \textbf{Étape (lignes supprimées)} & \textbf{Résultat} \\
\hline
\endhead
Online Retail II &
Doublons exacts (34 335) ; factures annulées « C » et ajustements (19 110) ; quantité ou prix $\le 0$ (6 014) ; codes qui ne sont pas des produits : frais de port, remises, frais bancaires (4 698) ; client inconnu (226 637) ; annotations internes « damaged », « wet » (479) ; valeurs au-delà du 99,9\textsuperscript{e} centile (1 507). &
774 591 lignes de vente \\
\hline
Cosmetics &
Doublons exacts, par exemple double clic (183 860) ; prix $\le 0$ (7 537) ; session manquante (714) ; \textbf{robots} : plus de 500 événements par jour ou 60 par minute, 136 comptes (78 998) ; prix aberrant par rapport à la médiane du produit (696) ; échantillon aléatoire de 60 000 visiteurs pour tenir en mémoire. &
525 253 événements, 60 000 visiteurs \\
\hline
H\&M &
Échantillon d'un client sur seize ; achats en magasin retirés, Sellio étant une boutique en ligne (588 289) ; prix au-delà du 99,9\textsuperscript{e} centile (1 394) ; achats identiques du même jour regroupés en une ligne avec quantité (167 798). Traduction des articles en fiches Sellio (genre, couleur, tissu) et mise à l'échelle des prix anonymisés (article médian $\approx$ 50 DT). &
1 224 099 achats, 68 874 clients, 71 735 articles \\
\hline
\caption{Principales étapes de nettoyage}
\label{tab:nettoyage}
\end{longtable}

La traduction du jeu H\&M vers le modèle de données de Sellio a aussi servi de test d'adéquation : la couleur a pu être renseignée pour 98\,\% des articles et le tissu pour 80\,\%, ce qui confirme que les attributs définis au Sprint 3 couvrent les données d'une vraie enseigne de mode.

\subsection{Protocole d'évaluation}
Deux règles ont guidé toutes les évaluations :
\begin{itemize}
  \item \textbf{Découpage temporel :} un modèle est toujours entraîné sur le passé et évalué sur la période qui suit, comme en production. Pour l'attrition, trois dates d'observation T1, T2 et T3 sont utilisées : entraînement à T1, choix du modèle et des réglages à T2, puis évaluation finale une seule fois à T3, le modèle retenu étant ré-entraîné sur T1 et T2. Un découpage aléatoire aurait laissé le modèle « voir le futur » et gonflé ses scores.
  \item \textbf{Comparaison à une référence simple :} chaque modèle est comparé à une règle naïve (produits les plus populaires, ou ancienneté du dernier achat pour l'attrition). Un modèle qui ne bat pas sa référence n'apporte rien et n'est pas retenu.
\end{itemize}

Le tableau~\ref{tab:metriques} définit les métriques utilisées.
\renewcommand{\arraystretch}{1.3}
\begin{longtable}{|p{3.2cm}|p{11.4cm}|}
\hline
\textbf{Métrique} & \textbf{Signification} \\
\hline
\endfirsthead
\hline
\textbf{Métrique} & \textbf{Signification} \\
\hline
\endhead
ROC-AUC & Probabilité qu'un client parti reçoive un score plus élevé qu'un client resté (0,5 = hasard, 1 = parfait). \\
\hline
Précision & Parmi les clients signalés « à risque », part de ceux qui partent vraiment. \\
\hline
Rappel & Parmi les clients qui partent vraiment, part de ceux que le modèle a signalés. \\
\hline
F1 & Moyenne harmonique de la précision et du rappel. \\
\hline
HR@10 (\textit{hit rate}) & Part des cas où au moins un produit réellement acheté ensuite figure dans les 10 produits recommandés. \\
\hline
NDCG@10 & Comme HR@10, mais récompense davantage un bon produit placé en haut de la liste. \\
\hline
Couverture & Part du catalogue qui apparaît au moins une fois dans les recommandations. \\
\hline
Silhouette & Qualité d'un regroupement : proximité d'un visiteur avec son groupe comparée au groupe voisin (de $-1$ à $1$). \\
\hline
Davies-Bouldin & Rapport entre dispersion interne et séparation des groupes (plus bas = mieux). \\
\hline
ARI & Accord entre deux regroupements ; utilisé pour mesurer la stabilité des segments d'un entraînement à l'autre (1 = identiques). \\
\hline
\caption{Métriques d'évaluation}
\label{tab:metriques}
\end{longtable}

\section{Moteur de recommandation}
\subsection{Choix du modèle}
Le moteur doit fonctionner pour des boutiques de toutes tailles, y compris une boutique qui vient d'ouvrir (problème du « démarrage à froid »), et produire des recommandations rapides à servir. Le tableau~\ref{tab:comparaison-reco} compare les approches candidates.
\renewcommand{\arraystretch}{1.4}
\begin{longtable}{|p{3.1cm}|p{4.3cm}|p{4.3cm}|p{2.5cm}|}
\hline
\textbf{Approche} & \textbf{Avantages} & \textbf{Contraintes} & \textbf{Décision} \\
\hline
\endfirsthead
\hline
\textbf{Approche} & \textbf{Avantages} & \textbf{Contraintes} & \textbf{Décision} \\
\hline
\endhead
Popularité & Simple, robuste, forte sur les produits saisonniers. & Identique pour tous : aucune personnalisation. & Référence, et composante du modèle retenu. \\
\hline
Filtrage par contenu (TF-IDF) & Fonctionne dès le premier jour, à partir des fiches produit. & Recommande des produits trop semblables ; ignore les goûts réels. & Composante du modèle retenu. \\
\hline
Filtrage collaboratif (factorisation SVD) & Apprend des comportements réels ; découvre des associations inattendues. & Inutilisable pour un produit sans interactions. & Composante du modèle retenu. \\
\hline
Réseaux de neurones (modèles séquentiels) & Très performants sur de grands volumes. & Exigent beaucoup de données et un GPU ; peu explicables. & Écartée. \\
\hline
\textbf{Hybride pondéré} & Combine les forces des trois premières ; s'adapte au volume de données de chaque produit. & Deux paramètres à régler par boutique. & \textbf{Retenue.} \\
\hline
\caption{Comparaison des approches de recommandation}
\label{tab:comparaison-reco}
\end{longtable}

\subsection{Modèle retenu et configuration}
Le modèle hybride se construit en quatre étapes :
\begin{enumerate}
  \item \textbf{Signal implicite pondéré :} chaque événement reçoit un poids selon l'intérêt qu'il traduit (vue 1, clic depuis une recherche 1,5, favori 3, ajout au panier et début de commande 4, achat 8), atténué avec le temps (demi-vie de 30 jours).
  \item \textbf{Similarité collaborative :} la matrice visiteurs × produits est pondérée à la manière de l'IDF en recherche documentaire, pour éviter que les produits vus par tout le monde ne dominent, puis factorisée par \textbf{SVD tronquée} (64 à 128 facteurs latents). La similarité cosinus entre les vecteurs des produits donne la similarité collaborative.
  \item \textbf{Similarité de contenu :} les fiches produit (nom, catégorie, attributs) sont vectorisées par \textbf{TF-IDF}.
  \item \textbf{Mélange adaptatif :} pour deux produits $i$ et $j$, avec $n$ le plus petit nombre d'interactions des deux,
  \[ \text{sim}(i,j) = \beta \cdot \text{sim}_{\text{CF}}(i,j) + (1-\beta)\cdot \text{sim}_{\text{contenu}}(i,j), \qquad \beta = \frac{n}{n+\lambda}. \]
  Un produit riche en historique est recommandé d'après les comportements ; un produit nouveau, d'après sa description.
\end{enumerate}
Au moment de servir, le score personnel d'un produit est la somme de ses similarités avec les produits que le visiteur a consultés, pondérées par l'intérêt ; il est ensuite mélangé à la popularité :
\[ \text{score} = (1-\alpha)\cdot \text{personnel} + \alpha \cdot \text{popularité}. \]
Les paramètres $\lambda \in \{10, 30, 100, 300, 1000\}$ et $\alpha \in \{0; 0{,}2; 0{,}4; 0{,}6; 0{,}8\}$ sont \textbf{réglés automatiquement pour chaque boutique} en cachant, pour chaque visiteur, son dernier produit fortement signalé (ajout au panier ou achat) et en retenant les valeurs qui maximisent le NDCG@10 ; sans au moins 30 cas, les valeurs par défaut ($\alpha = 0{,}4$) s'appliquent.

Pour la section « souvent achetés ensemble », nous utilisons des \textbf{règles d'association} $A \rightarrow B$ extraites des paniers, classées par \textit{lift} $= P(B \mid A) / P(B)$ et ne conservant que les paires observées au moins deux fois.

\subsection{Résultats}
La tâche évaluée est la plus exigeante : prédire les \textbf{nouveaux} produits qu'un client achètera (Online Retail II, H\&M) ou le prochain produit mis au panier (Cosmetics), parmi plusieurs milliers. Le tableau~\ref{tab:resultats-reco} présente les résultats sur la période de test.

\renewcommand{\arraystretch}{1.3}
\begin{longtable}{|p{5.2cm}|p{2.9cm}|p{2.9cm}|p{2.9cm}|}
\hline
\textbf{Modèle (HR@10 / NDCG@10)} & \textbf{Online Retail II} & \textbf{Cosmetics} & \textbf{H\&M} \\
\hline
\endfirsthead
\hline
\textbf{Modèle (HR@10 / NDCG@10)} & \textbf{Online Retail II} & \textbf{Cosmetics} & \textbf{H\&M} \\
\hline
\endhead
Popularité (référence) & 0,308 / 0,052 & 0,159 / 0,044 & 0,053 / 0,015 \\
\hline
Contenu seul (TF-IDF) & 0,259 / 0,044 & 0,069 / 0,019 & 0,038 / 0,014 \\
\hline
Filtrage collaboratif seul (SVD) & 0,240 / 0,044 & 0,084 / 0,020 & 0,020 / 0,008 \\
\hline
Hybride & 0,337 / 0,062 & 0,079 / 0,022 & 0,039 / 0,015 \\
\hline
\textbf{Hybride + popularité (retenu)} & \textbf{0,357 / 0,069} ($\alpha=0{,}4$) & \textbf{0,178 / 0,052} ($\alpha=0{,}8$) & \textbf{0,059 / 0,022} ($\alpha=0{,}4$) \\
\hline
Cas évalués & 1 500 & 1 373 & 2 000 \\
\hline
\caption{Résultats du moteur de recommandation sur la période de test}
\label{tab:resultats-reco}
\end{longtable}

Trois constats se dégagent :
\begin{itemize}
  \item \textbf{Le modèle retenu bat la référence sur les trois jeux}, de +16\,\% (Online Retail II), +12\,\% (Cosmetics) et +10\,\% (H\&M) en HR@10, et de +34\,\%, +17\,\% et +41\,\% en NDCG@10 : il place les bons produits plus haut dans la liste.
  \item \textbf{La popularité est une référence difficile à battre.} Sur Cosmetics, les produits consommables rachetés par tous pèsent lourd, d'où un $\alpha$ élevé (0,8) choisi automatiquement ; c'est précisément pour cela qu'$\alpha$ est réglé par boutique plutôt que fixé.
  \item \textbf{Les scores absolus restent modestes}, notamment sur H\&M (5,9\,\%), où il faut deviner un article parmi 5 000 dans un catalogue de mode très renouvelé. Ces valeurs sont cohérentes avec la difficulté de la tâche ; elles ne doivent pas être comparées aux scores de 90\,\% et plus que l'on obtient sur des données artificielles.
\end{itemize}
Enfin, pour « souvent achetés ensemble », les règles d'association placent un produit réellement acheté dans le même panier parmi les 5 suggestions dans \textbf{47,1\,\%} des cas, contre 43,0\,\% pour la popularité, et couvrent 90\,\% des produits.

\section{Segmentation des visiteurs}
\subsection{Choix du modèle}
Les niveaux de fidélité du chapitre 6 sont définis à la main à partir des points. La segmentation comportementale cherche au contraire des groupes dans les données, sans étiquette préalable : c'est un problème d'apprentissage \textbf{non supervisé}. Le tableau~\ref{tab:comparaison-seg} compare les algorithmes candidats.
\renewcommand{\arraystretch}{1.4}
\begin{longtable}{|p{3.1cm}|p{4.3cm}|p{4.3cm}|p{2.5cm}|}
\hline
\textbf{Algorithme} & \textbf{Avantages} & \textbf{Contraintes} & \textbf{Décision} \\
\hline
\endfirsthead
\hline
\textbf{Algorithme} & \textbf{Avantages} & \textbf{Contraintes} & \textbf{Décision} \\
\hline
\endhead
\textbf{K-Means} & Rapide sur des dizaines de milliers de visiteurs ; centres de groupes directement interprétables. & Nombre de groupes à choisir ; sensible aux valeurs extrêmes. & \textbf{Retenu.} \\
\hline
DBSCAN & Trouve des groupes de forme libre et isole le bruit. & Paramètres délicats en dimension 9 ; beaucoup de visiteurs classés « bruit ». & Écarté. \\
\hline
Mélange gaussien (GMM) & Appartenance probabiliste. & Plus lent, moins stable ; interprétation moins directe. & Écarté. \\
\hline
Classification hiérarchique & Visualisation en dendrogramme. & Coût quadratique en mémoire : inadapté à 60 000 visiteurs. & Écartée. \\
\hline
\caption{Comparaison des algorithmes de segmentation}
\label{tab:comparaison-seg}
\end{longtable}

\subsection{Configuration}
Chaque visiteur est décrit par neuf variables : nombre de sessions, profondeur des sessions, produits vus, part des recherches, taux d'ajout au panier, taux d'achat, taux d'ajout aux favoris, ancienneté de la dernière visite et nombre de jours actifs. Trois traitements rendent K-Means fiable sur ces données :
\begin{itemize}
  \item \textbf{transformation logarithmique} ($\log(1+x)$) des comptages, très asymétriques ;
  \item \textbf{standardisation}, pour que chaque variable pèse autant ;
  \item \textbf{écrêtage à $\pm 3$ écarts-types} : sans lui, quelques visiteurs extrêmes formaient à eux seuls un groupe d'une seule personne.
\end{itemize}
Le nombre de groupes $K$ est choisi entre 3 et 6 par le meilleur score de silhouette ; $K = 2$ est exclu car deux groupes (en général « dormants » et « tous les autres ») sont trop grossiers pour agir. Chaque groupe reçoit enfin un nom et une recommandation d'action déduits de son centre (par exemple « Abandonnistes de panier : relance panier, code livraison offerte »).

\subsection{Résultats}
Sur les 60 000 visiteurs du jeu Cosmetics, le meilleur score de silhouette est obtenu pour $K = 3$ (0,468, contre 0,385 à 0,418 pour $K$ de 4 à 6). L'indice de Davies-Bouldin vaut 1,26 et la \textbf{stabilité} mesurée par l'ARI moyen entre entraînements vaut 0,951 : les segments ne changent pas d'un entraînement à l'autre. Le tableau~\ref{tab:segments} décrit les trois segments.
\renewcommand{\arraystretch}{1.3}
\begin{longtable}{|p{3.2cm}|p{1.8cm}|p{5.2cm}|p{4.2cm}|}
\hline
\textbf{Segment} & \textbf{Part} & \textbf{Profil moyen} & \textbf{Action suggérée} \\
\hline
\endfirsthead
\hline
\textbf{Segment} & \textbf{Part} & \textbf{Profil moyen} & \textbf{Action suggérée} \\
\hline
\endhead
Visiteurs occasionnels & 76\,\% (45 694) & 1,1 session, 1,4 produit vu, quasiment aucun achat. & Capter l'e-mail (newsletter, première commande). \\
\hline
Acheteurs décidés & 13\,\% (7 615) & 1,4 session, 15 événements par session, 8 produits vus, 1 achat pour 4 ajouts au panier. & Mettre en avant les nouveautés et le réassort. \\
\hline
Explorateurs & 11\,\% (6 691) & 5,6 sessions, 22 produits vus sur 4 jours actifs, achètent moins que les décidés. & Recommandations et contenus inspirants. \\
\hline
\caption{Segments de visiteurs obtenus sur le jeu Cosmetics}
\label{tab:segments}
\end{longtable}
Une silhouette de 0,47 indique des groupes nettement séparés sans l'être parfaitement, ce qui est attendu pour des comportements humains qui forment un continuum.

\section{Prédiction de l'attrition}
\subsection{Définition du problème}
Un client est considéré comme \textbf{perdu} s'il ne passe aucune commande dans les \textbf{90 jours} qui suivent la date d'observation. Le modèle estime cette probabilité pour chaque client ayant déjà acheté, à partir de six variables calculées sur son historique : ancienneté, nombre de commandes, montant total dépensé, panier moyen, fréquence d'achat et nombre de jours depuis la dernière commande. Le jeu H\&M ajoute l'âge du client.

\subsection{Choix du modèle}
Le tableau~\ref{tab:comparaison-churn} compare les modèles candidats. Tous ont été entraînés à T1, puis comparés à T2 selon l'AUC, pour plusieurs réglages (profondeur des arbres, taille minimale des feuilles, taux d'apprentissage).
\renewcommand{\arraystretch}{1.4}
\begin{longtable}{|p{3.1cm}|p{4.3cm}|p{4.3cm}|p{2.5cm}|}
\hline
\textbf{Modèle} & \textbf{Avantages} & \textbf{Contraintes} & \textbf{Décision} \\
\hline
\endfirsthead
\hline
\textbf{Modèle} & \textbf{Avantages} & \textbf{Contraintes} & \textbf{Décision} \\
\hline
\endhead
Régression logistique & Simple, rapide, coefficients lisibles. & Relations linéaires uniquement (après transformation logarithmique). & Référence forte. \\
\hline
Forêt aléatoire & Capture les effets non linéaires et les interactions ; robuste au surapprentissage grâce à la moyenne de nombreux arbres. & Moins lisible qu'une régression ; modèle plus volumineux. & \textbf{Retenue (secteur général).} \\
\hline
Gradient boosting & Souvent le plus précis sur données tabulaires. & Plus sensible aux réglages. & \textbf{Retenu (secteur vêtements).} \\
\hline
Réseau de neurones & Très flexible. & Excessif pour six variables ; peu explicable. & Écarté. \\
\hline
\caption{Comparaison des modèles de prédiction de l'attrition}
\label{tab:comparaison-churn}
\end{longtable}

\subsection{Résultats}
Le tableau~\ref{tab:resultats-churn} présente les résultats sur la période de test T3, jamais utilisée pour choisir le modèle.
\renewcommand{\arraystretch}{1.3}
\begin{longtable}{|p{5.4cm}|p{1.6cm}|p{1.8cm}|p{1.6cm}|p{1.6cm}|}
\hline
\textbf{Modèle} & \textbf{AUC} & \textbf{Précision} & \textbf{Rappel} & \textbf{F1} \\
\hline
\endfirsthead
\hline
\textbf{Modèle} & \textbf{AUC} & \textbf{Précision} & \textbf{Rappel} & \textbf{F1} \\
\hline
\endhead
\multicolumn{5}{|c|}{\textbf{Online Retail II (5 239 clients, taux d'attrition 56,5\,\%)}} \\
\hline
Référence : ancienneté du dernier achat & 0,762 & 0,844 & 0,375 & 0,519 \\
\hline
Régression logistique & 0,791 & 0,737 & 0,804 & 0,769 \\
\hline
\textbf{Forêt aléatoire (profondeur 6), retenue} & \textbf{0,801} & 0,728 & \textbf{0,849} & \textbf{0,784} \\
\hline
\multicolumn{5}{|c|}{\textbf{H\&M (64 871 clients, taux d'attrition 73,1\,\%)}} \\
\hline
Référence : ancienneté du dernier achat & 0,738 & 0,915 & 0,349 & 0,506 \\
\hline
Modèle général appliqué tel quel & 0,787 & 0,831 & 0,887 & 0,858 \\
\hline
Gradient boosting sans l'âge & 0,799 & 0,878 & 0,724 & 0,794 \\
\hline
\textbf{Gradient boosting (taux 0,05), retenu} & \textbf{0,799} & \textbf{0,880} & 0,718 & 0,791 \\
\hline
\caption{Résultats de la prédiction de l'attrition sur la période de test}
\label{tab:resultats-churn}
\end{longtable}

Ces résultats appellent plusieurs remarques :
\begin{itemize}
  \item \textbf{Un gain réel mais mesuré.} Avec une AUC d'environ 0,80, le modèle classe correctement un client parti devant un client resté quatre fois sur cinq. Il dépasse nettement la règle naïve, surtout en rappel : sur Online Retail II, il repère 85\,\% des clients qui partent, contre 38\,\% pour la règle.
  \item \textbf{Des variables explicables.} L'importance par permutation montre que le nombre de jours depuis la dernière commande et la fréquence d'achat dominent sur Online Retail II ; sur H\&M, c'est le nombre de commandes. Le modèle confirme donc l'intuition métier, en la quantifiant.
  \item \textbf{Un modèle par secteur.} Le modèle général fonctionne déjà sur la mode (AUC 0,787), ce qui montre que les comportements d'achat se ressemblent, mais le modèle entraîné sur H\&M fait mieux (0,799) et, surtout, il est plus précis (88\,\% contre 83\,\%) : les alertes envoyées au commerçant sont plus fiables. Sellio choisit donc le modèle selon le secteur de la boutique.
  \item \textbf{L'âge n'apporte rien} (0,7988 sans, 0,7993 avec) : nous n'avons pas besoin de demander cette donnée personnelle aux clients.
\end{itemize}
Le score est enfin traduit en trois niveaux de risque, lisibles par un commerçant : \textbf{élevé} pour les 30\,\% de clients les plus à risque, \textbf{moyen} pour les 30\,\% suivants, \textbf{faible} pour les autres. Les seuils sont calculés pour chaque modèle à partir de la distribution de ses scores.

\subsection{Limites}
Ces modèles ont été entraînés sur des données européennes antérieures à 2021. Le comportement des clients tunisiens peut différer (paiement à la livraison dominant, saisonnalité du Ramadan et des soldes). Le système est conçu pour cela : la recommandation et la segmentation sont ré-entraînées chaque nuit sur les données propres de chaque boutique, et l'attrition pourra l'être dès que les boutiques auront accumulé plusieurs mois d'historique.

%==================================================================
\section{Backlog du Sprint 10}
\Needspace{18\baselineskip}
Le tableau~\ref{tab:backlog-sprint10} présente le \textbf{backlog du Sprint 10}, consacré au suivi comportemental et à la recommandation.

\renewcommand{\arraystretch}{1.05}
\setlength{\tabcolsep}{3pt}
\begin{longtable}{|>{\raggedright\arraybackslash}p{4cm}|>{\raggedright\arraybackslash}p{9cm}|>{\centering\arraybackslash}p{1.2cm}|}
\hline
\textbf{User Story} & \textbf{Tâches} & \textbf{Estimation} \\
\hline
\endfirsthead
\hline
\textbf{User Story} & \textbf{Tâches} & \textbf{Estimation} \\
\hline
\endhead
\hline
\multicolumn{3}{r}{\textit{Suite à la page suivante}} \\
\endfoot
\hline
\endlastfoot

\textbf{En tant que système, je veux collecter les événements de navigation}
& Développer le module \texttt{tracker.js} : identifiants de visiteur et de session, file d'attente et envoi par lots. & 1 \\
\cline{2-3}
& Instrumenter la vitrine (vues, recherches, panier, favoris, commande, achat, clics sur recommandation). & 1 \\
\cline{2-3}
& Développer l'API de réception avec validation des types d'événements. & 1 \\
\cline{2-3}
& Créer la table partitionnée, les index BRIN et partiels, et la gestion automatique des partitions. & 1 \\
\cline{2-3}
& Développer les agrégats quotidiens et la rétention de 13 mois. & 1 \\
\cline{2-3}
& Tester la fonctionnalité. & 1 \\
\hline

\textbf{En tant que visiteur, je souhaite voir des produits similaires et « pour vous »}
& Nettoyer les jeux Online Retail II, Cosmetics et H\&M (scripts et rapports). & 1 \\
\cline{2-3}
& Développer le modèle hybride (signal pondéré, SVD, TF-IDF, mélange adaptatif). & 1 \\
\cline{2-3}
& Développer le réglage automatique de $\lambda$ et $\alpha$ par boutique. & 1 \\
\cline{2-3}
& Évaluer le modèle sur les trois jeux avec découpage temporel et référence. & 1 \\
\cline{2-3}
& Développer \texttt{PythonRunner} et \texttt{MlTrainingService} (entraînement nocturne et à la demande). & 1 \\
\cline{2-3}
& Développer les API « similaires », « pour vous » et « populaires ». & 1 \\
\cline{2-3}
& Développer le composant \texttt{RecommendedProducts} sur la fiche produit et l'accueil. & 1 \\
\hline

\textbf{En tant que visiteur, je souhaite voir les produits souvent achetés ensemble}
& Extraire les règles d'association classées par \textit{lift} et les évaluer. & 1 \\
\cline{2-3}
& Développer l'API « achetés ensemble » et la section du panier. & 1 \\
\hline

\multicolumn{2}{|r|}{\textbf{Total}} & \textbf{15} \\
\hline
\caption{Backlog du Sprint 10}
\label{tab:backlog-sprint10}
\end{longtable}

\section{Diagramme de cas d'utilisation du Sprint 10}
\Needspace{20\baselineskip}
\begin{figure}[H]
    \centering
    \img{0.8}{cu_sprint10.png}
    \caption{Diagramme de cas d'utilisation du Sprint 10}
    \label{fig:cu-sprint10}
\end{figure}

\section{Diagramme de classes du Sprint 10}
\Needspace{20\baselineskip}
La figure~\ref{fig:classes-ml} présente les classes du microservice d'analyse impliquées dans ce sprint (\texttt{TrackingController}, \texttt{TrackingService}, \texttt{EventStore}, \texttt{MlTrainingService}, \texttt{PythonRunner}, \texttt{RecommendationService}).
\begin{figure}[H]
    \centering
    \img{0.95}{classes_analytics.png}
    \caption{Diagramme de classes du microservice d'analyse}
    \label{fig:classes-ml}
\end{figure}

\section{Services web du Sprint 10}
Le tableau~\ref{tab:ws-sprint10} présente les services web du Sprint 10, exposés par \texttt{analytics-service}.
\renewcommand{\arraystretch}{1.15}
\setlength{\tabcolsep}{4pt}
\begin{longtable}{|p{1.5cm}|p{6.8cm}|p{6.2cm}|}
\hline
\textbf{Méthode} & \textbf{URL} & \textbf{Description} \\
\hline
\endfirsthead
\hline
\textbf{Méthode} & \textbf{URL} & \textbf{Description} \\
\hline
\endhead
\multicolumn{3}{|c|}{\textbf{Collecte --- /api/v1/analytics/events}} \\
\hline
POST & /api/v1/analytics/events & Recevoir un lot d'événements (boutique, visiteur, session). \\
\hline
\multicolumn{3}{|c|}{\textbf{Recommandations --- /api/v1/analytics/recommendations}} \\
\hline
GET & /api/v1/analytics/recommendations/similar & Produits similaires à un produit. \\
\hline
GET & /api/v1/analytics/recommendations/for-you & Recommandations personnalisées pour un visiteur. \\
\hline
GET & /api/v1/analytics/recommendations/bought-together & Produits souvent achetés avec ceux du panier. \\
\hline
GET & /api/v1/analytics/recommendations/popular & Produits populaires (repli et démarrage à froid). \\
\hline
\caption{Services web du Sprint 10}
\label{tab:ws-sprint10}
\end{longtable}

\section{Les cas d'utilisation du Sprint 10}
\subsection{Cas d'utilisation « Recevoir des recommandations personnalisées »}
\subsubsection{Description textuelle}
Le tableau~\ref{tab:uc-reco} présente la description textuelle de ce cas.
\renewcommand{\arraystretch}{1.5}
\begin{longtable}{|p{4cm}|p{11cm}|}
\hline
\textbf{Acteurs} & Visiteur ou Client, Moteur d'apprentissage automatique \\
\hline
\textbf{Objectif} & Afficher à chaque visiteur des produits adaptés à ses centres d'intérêt. \\
\hline
\textbf{Pré-condition} & Le modèle de la boutique a été entraîné au moins une fois. \\
\hline
\textbf{Scénario principal} &
1. Le visiteur consulte des produits ; chaque consultation est enregistrée par le module de suivi.\newline
2. À l'affichage de l'accueil ou d'une fiche produit, la vitrine demande les recommandations « pour vous » en transmettant l'identifiant du visiteur.\newline
3. Le système récupère l'historique récent du visiteur, pondéré par type d'événement et par ancienneté.\newline
4. Le système lit les voisins pré-calculés de chaque produit de l'historique, en fait la somme pondérée et la mélange à la popularité avec l'$\alpha$ de la boutique.\newline
5. Le système exclut les produits que le visiteur a déjà mis au panier ou achetés, et renvoie les meilleurs.\newline
6. La vitrine affiche les produits ; un clic est enregistré comme « clic sur recommandation », ce qui permet de mesurer l'efficacité du moteur en production. \\
\hline
\textbf{Scénario alternatif} &
\textbf{Nouveau visiteur sans historique :} à l'étape 3, le système renvoie les produits populaires de la boutique.\newline
\textbf{Modèle pas encore entraîné :} le système se replie sur les produits populaires de la même catégorie, puis sur les nouveautés.\newline
\textbf{Nouveau produit sans interactions :} ses voisins sont calculés d'après sa description (TF-IDF), grâce au mélange adaptatif.\newline
\textbf{Service d'analyse indisponible :} la vitrine masque la section sans bloquer la navigation. \\
\hline
\caption{Description textuelle du cas « Recevoir des recommandations personnalisées »}
\label{tab:uc-reco}
\end{longtable}

\subsubsection{Diagramme de séquence système}
\Needspace{18\baselineskip}
\begin{figure}[H]
    \centering
    \img{0.7}{dss_reco.png}
    \caption{Diagramme de séquence système du cas « Recevoir des recommandations personnalisées »}
    \label{fig:dss-reco}
\end{figure}

\subsubsection{Diagramme de séquence objet}
\Needspace{20\baselineskip}
\begin{figure}[H]
    \centering
    \img{0.95}{dso_reco.png}
    \caption{Diagramme de séquence objet du cas « Recevoir des recommandations personnalisées »}
    \label{fig:dso-reco}
\end{figure}

\section{Réalisation du Sprint 10}
\Needspace{20\baselineskip}
\textbf{Recommandations sur la vitrine :} les figures~\ref{fig:ui-reco-fiche} et~\ref{fig:ui-reco-panier} présentent les sections « Vous aimerez aussi » sur la fiche produit et « Souvent achetés ensemble » dans le panier.
\begin{figure}[H]
    \centering
    \img{0.85}{ui_reco_fiche.png}
    \caption{Recommandations sur la fiche produit}
    \label{fig:ui-reco-fiche}
\end{figure}
\begin{figure}[H]
    \centering
    \img{0.75}{ui_reco_panier.png}
    \caption{Produits souvent achetés ensemble dans le panier}
    \label{fig:ui-reco-panier}
\end{figure}

%==================================================================
\section{Backlog du Sprint 11}
\Needspace{18\baselineskip}
Le tableau~\ref{tab:backlog-sprint11} présente le \textbf{backlog du Sprint 11}.

\renewcommand{\arraystretch}{1.05}
\setlength{\tabcolsep}{3pt}
\begin{longtable}{|>{\raggedright\arraybackslash}p{4cm}|>{\raggedright\arraybackslash}p{9cm}|>{\centering\arraybackslash}p{1.2cm}|}
\hline
\textbf{User Story} & \textbf{Tâches} & \textbf{Estimation} \\
\hline
\endfirsthead
\hline
\textbf{User Story} & \textbf{Tâches} & \textbf{Estimation} \\
\hline
\endhead
\hline
\multicolumn{3}{r}{\textit{Suite à la page suivante}} \\
\endfoot
\hline
\endlastfoot

\textbf{En tant que commerçant, je souhaite voir les segments de mes visiteurs}
& Calculer les neuf variables de comportement par visiteur. & 1 \\
\cline{2-3}
& Développer le K-Means avec transformation, standardisation, écrêtage et choix de $K$ par silhouette. & 1 \\
\cline{2-3}
& Nommer les segments et leur associer une action à partir des centres. & 1 \\
\cline{2-3}
& Évaluer la qualité et la stabilité (silhouette, Davies-Bouldin, ARI) sur le jeu Cosmetics. & 1 \\
\hline

\textbf{En tant que commerçant, je souhaite identifier les clients à risque d'attrition}
& Construire les instantanés T1, T2, T3 et la cible à 90 jours. & 1 \\
\cline{2-3}
& Comparer régression logistique, forêts aléatoires et gradient boosting ; choisir sur T2, évaluer sur T3. & 1 \\
\cline{2-3}
& Entraîner le modèle vêtements sur H\&M et tester l'apport de l'âge. & 1 \\
\cline{2-3}
& Calculer l'importance des variables et les seuils de risque. & 1 \\
\cline{2-3}
& Développer \texttt{predict.py} et \texttt{ChurnPredictionService} (choix du modèle selon le secteur). & 1 \\
\cline{2-3}
& Développer les API de score individuel et par lot. & 1 \\
\hline

\textbf{En tant que commerçant, je souhaite consulter le tableau de bord « Comportement » et ré-entraîner les modèles}
& Développer l'API de synthèse (entonnoir, événements, produits consultés, segments, dernier entraînement). & 1 \\
\cline{2-3}
& Développer la page « Comportement » du back-office. & 1 \\
\cline{2-3}
& Développer le ré-entraînement à la demande et les données de démonstration (création et suppression). & 1 \\
\cline{2-3}
& Restreindre l'accès aux seules données de la boutique du commerçant. & 1 \\
\cline{2-3}
& Rédiger la documentation des modèles (README ML). & 1 \\
\cline{2-3}
& Tester la fonctionnalité. & 1 \\
\hline

\multicolumn{2}{|r|}{\textbf{Total}} & \textbf{16} \\
\hline
\caption{Backlog du Sprint 11}
\label{tab:backlog-sprint11}
\end{longtable}

\section{Diagramme de cas d'utilisation du Sprint 11}
\Needspace{20\baselineskip}
\begin{figure}[H]
    \centering
    \img{0.8}{cu_sprint11.png}
    \caption{Diagramme de cas d'utilisation du Sprint 11}
    \label{fig:cu-sprint11}
\end{figure}

\section{Services web du Sprint 11}
Le tableau~\ref{tab:ws-sprint11} présente les services web du Sprint 11.
\renewcommand{\arraystretch}{1.15}
\setlength{\tabcolsep}{4pt}
\begin{longtable}{|p{1.5cm}|p{6.8cm}|p{6.2cm}|}
\hline
\textbf{Méthode} & \textbf{URL} & \textbf{Description} \\
\hline
\endfirsthead
\hline
\textbf{Méthode} & \textbf{URL} & \textbf{Description} \\
\hline
\endhead
\multicolumn{3}{|c|}{\textbf{Analyse comportementale --- /api/v1/analytics/behavior}} \\
\hline
GET & /api/v1/analytics/behavior/overview & Synthèse : entonnoir, événements, segments, dernier entraînement. \\
\hline
POST & /api/v1/analytics/behavior/train & Ré-entraîner les modèles de la boutique. \\
\hline
POST & /api/v1/analytics/behavior/demo-data & Générer des données de démonstration. \\
\hline
DELETE & /api/v1/analytics/behavior/demo-data & Supprimer les données de démonstration. \\
\hline
\multicolumn{3}{|c|}{\textbf{Attrition --- /api/v1/analytics/churn}} \\
\hline
GET & /api/v1/analytics/churn/\{userId\} & Score et niveau de risque d'un client. \\
\hline
GET & /api/v1/analytics/churn/\{userId\}/features & Variables utilisées pour le score d'un client. \\
\hline
POST & /api/v1/analytics/churn/batch-ids & Scores d'une liste de clients. \\
\hline
GET & /api/v1/analytics/churn/batch & Scores de tous les clients de la boutique. \\
\hline
\caption{Services web du Sprint 11}
\label{tab:ws-sprint11}
\end{longtable}

\section{Les cas d'utilisation du Sprint 11}
\subsection{Cas d'utilisation « Consulter l'analyse comportementale et les risques d'attrition »}
\subsubsection{Description textuelle}
Le tableau~\ref{tab:uc-comportement} présente la description textuelle de ce cas.
\renewcommand{\arraystretch}{1.5}
\begin{longtable}{|p{4cm}|p{11cm}|}
\hline
\textbf{Acteurs} & Commerçant, Moteur d'apprentissage automatique \\
\hline
\textbf{Objectif} & Comprendre le comportement des visiteurs et agir sur les clients à risque. \\
\hline
\textbf{Pré-condition} & Le commerçant est connecté ; sa boutique est active. \\
\hline
\textbf{Scénario principal} &
1. Le commerçant ouvre la page « Comportement ».\newline
2. Le système vérifie que la boutique demandée est bien la sienne.\newline
3. Le système renvoie l'entonnoir (visites, produits vus, ajouts au panier, commandes), les produits les plus consultés, la répartition des segments avec leur action suggérée et la date du dernier entraînement.\newline
4. Le commerçant ouvre la liste des clients ; le système calcule les scores d'attrition par lot avec le modèle du secteur de la boutique et affiche pour chacun un risque élevé, moyen ou faible.\newline
5. Le commerçant ouvre la fiche d'un client à risque élevé, consulte les variables qui expliquent son score et le contacte (e-mail, code promo ciblé).\newline
6. S'il le souhaite, il relance l'entraînement ; le système ré-entraîne recommandation et segmentation et affiche les nouvelles métriques. \\
\hline
\textbf{Scénario alternatif} &
\textbf{Accès à une autre boutique :} à l'étape 2, le système répond 403.\newline
\textbf{Boutique sans trafic :} pour une démonstration, le commerçant peut générer environ 60 jours de visites fictives (identifiants préfixés « demo- »), puis les supprimer en un clic. \\
\hline
\caption{Description textuelle du cas « Consulter l'analyse comportementale et les risques d'attrition »}
\label{tab:uc-comportement}
\end{longtable}

\subsubsection{Diagramme de séquence système}
\Needspace{18\baselineskip}
\begin{figure}[H]
    \centering
    \img{0.7}{dss_comportement.png}
    \caption{Diagramme de séquence système du cas « Consulter l'analyse comportementale »}
    \label{fig:dss-comportement}
\end{figure}

\subsubsection{Diagramme de séquence objet}
\Needspace{20\baselineskip}
\begin{figure}[H]
    \centering
    \img{0.95}{dso_comportement.png}
    \caption{Diagramme de séquence objet du cas « Consulter l'analyse comportementale »}
    \label{fig:dso-comportement}
\end{figure}

\section{Réalisation du Sprint 11}
\Needspace{20\baselineskip}
\textbf{Tableau de bord « Comportement » :} la figure~\ref{fig:ui-comportement} présente la page d'analyse comportementale : entonnoir de conversion, segments de visiteurs et actions suggérées, produits les plus consultés et état des modèles.
\begin{figure}[H]
    \centering
    \img{0.9}{ui_comportement.png}
    \caption{Tableau de bord « Comportement »}
    \label{fig:ui-comportement}
\end{figure}

\Needspace{20\baselineskip}
\textbf{Clients à risque :} la figure~\ref{fig:ui-churn} présente le niveau de risque d'attrition affiché dans la liste et la fiche des clients.
\begin{figure}[H]
    \centering
    \img{0.85}{ui_churn.png}
    \caption{Niveau de risque d'attrition des clients}
    \label{fig:ui-churn}
\end{figure}

\section*{Conclusion}
Cette dernière release a doté Sellio d'une véritable intelligence comportementale : collecte d'événements pensée pour le volume, recommandations hybrides, segmentation et prédiction de l'attrition. Chaque modèle a été choisi par comparaison, évalué sur des données réelles avec un protocole temporel, et comparé à une référence simple. Les résultats sont volontairement présentés tels quels : ils montrent un gain réel mais mesuré, cohérent avec la difficulté des tâches.


\thispagestyle{plain}
\chapter*{
  \vspace{-1.7cm}
  \centering
  \rule{\textwidth}{0.5pt} \\[0.4cm]
  Conclusion et perspectives
}
\addcontentsline{toc}{chapter}{\numberline{}Conclusion et perspectives}
Au terme de ce projet de fin d'études, nous avons conçu et développé \textbf{Sellio}, une plateforme transactionnelle intelligente qui permet à des commerçants de cosmétiques et de vêtements de créer, personnaliser et exploiter leur propre boutique en ligne. Partie d'une boutique unique, NaturEssence, la solution est devenue une plateforme multi-boutiques fondée sur une architecture microservices : cinq services Spring Boot derrière une passerelle, deux applications React et une base PostgreSQL. \newline

La plateforme couvre l'ensemble du cycle d'une boutique : intégration et vérification des commerçants (KYC par carte bancaire, pièce d'identité et selfie), catalogue aux attributs normalisés, trois gabarits de vitrine personnalisables sans code, tunnel de commande avec paiement Stripe, gestion des commandes, des retours, de la TVA et de la livraison, promotions, e-mailing et programme de fidélité.\newline

Elle se distingue par deux volets intelligents. Le premier est une \textbf{cabine d'essayage virtuel en temps réel}, fondée sur le modèle génératif \textit{Lucy VTON} de Decart et intégrée de façon sécurisée par des jetons éphémères. Le second est un ensemble de \textbf{modèles d'apprentissage automatique} alimentés par le suivi comportemental des visiteurs : un moteur de recommandation hybride, une segmentation des visiteurs par K-Means et une prédiction de l'attrition adaptée au secteur de chaque boutique. Ces modèles ont été entraînés et évalués sur trois jeux de données publics réels, nettoyés par des scripts reproductibles, avec un découpage temporel et une comparaison systématique à une référence simple. Les gains obtenus sont réels et mesurés : +10 à +16\,\% de taux de réussite pour la recommandation, une AUC d'environ 0,80 pour l'attrition, des segments stables et interprétables.\newline

Sur le plan professionnel, ce projet nous a permis de consolider nos compétences en architecture distribuée, en sécurité des applications web, en intégration de services tiers (paiement, IA générative, e-mail), en ingénierie des données et en apprentissage automatique, tout en appliquant une démarche Agile fondée sur Scrum. Il nous a surtout appris qu'un modèle ne vaut que par la rigueur de son évaluation : un score présenté sans référence ni protocole ne dit rien.\newline

\textbf{Perspectives.} Plusieurs évolutions peuvent être envisagées. Sur le plan des paiements, l'intégration d'un \textbf{prestataire de paiement local} prenant en charge le dinar tunisien, la confirmation des commandes par \textit{webhook} et la revérification serveur des prix et du stock sont les prochaines étapes indispensables avant une mise en production. La vérification des commerçants pourrait être automatisée par une \textbf{reconnaissance faciale} comparant le selfie à la photo de la pièce d'identité, l'administrateur n'intervenant plus que sur les cas douteux.

Sur le plan de l'architecture, le passage à une \textbf{base de données par microservice}, une communication asynchrone par messages (par exemple pour les événements de commande), la conteneurisation et le \textbf{déploiement dans le Cloud} amélioreraient l'isolation, la disponibilité et la capacité de montée en charge. Sur le plan de l'intelligence, le modèle d'attrition pourrait être \textbf{ré-entraîné sur les données tunisiennes} dès que les boutiques auront accumulé plusieurs mois d'historique, et les recommandations évaluées en production par des \textbf{tests A/B} mesurant leur effet réel sur le chiffre d'affaires. Enfin, une \textbf{application mobile} pour les clients et pour les commerçants élargirait l'usage de la plateforme.

Ainsi, Sellio pourrait évoluer d'une plateforme de création de boutiques vers une véritable solution de commerce intelligent, locale, sécurisée et pilotée par les données de chaque commerçant.


\bibliographystyle{unsrt}
\bibliography{References}

\IfFileExists{lastpage.pdf}
  {\includepdf[pages=-,pagecommand={\thispagestyle{empty}}]{lastpage.pdf}}
  {% Quatrième de couverture provisoire : remplacée automatiquement par
% lastpage.pdf si ce fichier (modèle officiel ESPRIT) est ajouté au dossier.
\clearpage
\thispagestyle{empty}
\noindent\textbf{\large Résumé}\\[0.3cm]
Ce rapport présente le travail réalisé au sein d'Antigone dans le cadre de notre projet de fin d'études. Il porte sur la conception et le développement de Sellio, une plateforme transactionnelle intelligente de type SaaS qui permet à des commerçants de cosmétiques et de vêtements de créer et d'exploiter leur boutique en ligne. La plateforme repose sur une architecture microservices (Spring Boot, Spring Cloud Gateway, PostgreSQL) et des interfaces React. Elle intègre la vérification d'identité des commerçants, trois gabarits de vitrine personnalisables, le paiement par Stripe, des outils marketing et de fidélisation, un essayage virtuel en temps réel fondé sur l'API Decart, ainsi qu'un suivi comportemental exploité par apprentissage automatique : recommandation hybride, segmentation des visiteurs et prédiction de l'attrition, évaluées sur des jeux de données publics réels. Le projet a été conduit selon la méthode Scrum.

\vspace{0.3cm}
\noindent\textbf{Mots clés :} e-commerce, SaaS, microservices, Spring Boot, React, Stripe, essayage virtuel, apprentissage automatique, système de recommandation, attrition, Scrum.

\vspace{1.2cm}
\noindent\textbf{\large Abstract}\\[0.3cm]
This report presents the work carried out at Antigone as our end-of-studies project. It covers the design and development of Sellio, an intelligent transactional SaaS platform that lets cosmetics and clothing merchants create and run their own online store. The platform is built on a microservices architecture (Spring Boot, Spring Cloud Gateway, PostgreSQL) with React front ends. It includes merchant identity verification, three customisable storefront templates, Stripe payments, marketing and loyalty tools, real-time virtual try-on based on the Decart API, and behavioural tracking used by machine-learning models: hybrid recommendation, visitor segmentation and churn prediction, evaluated on real public datasets. The project followed the Scrum framework.

\vspace{0.3cm}
\noindent\textbf{Keywords:} e-commerce, SaaS, microservices, Spring Boot, React, Stripe, virtual try-on, machine learning, recommender system, churn, Scrum.
}

\end{document}
